\documentclass[master=ewtk,masteroption=,openany]{kulemt}
\setup{
  title={Topology Optimization of Flow Configurations Operating in the Turbulent Regime},
  author={Tiebert Lefebure},
  promotor={Prof. Dr. Ir. Niels Horsten\and Prof. Dr. Ir. Maarten Blommaert},
  assessor={Prof. Dr. Ir. Elke Deckers\and Ir. Olivier Renders},
  assistant={Ir. Esteban Foronda Obando}}
% Remove the "%" on the next line for generating the cover page
%\setup{coverpageonly}
% Remove the "%" before the next "\setup" to generate only the first pages
% (e.g., if you are a Word user).
%\setup{frontpagesonly}

% If you want to include other LaTeX packages, do it here. 

% Finally the hyperref package is used for pdf files.
% This can be commented out for printed versions.
\usepackage[pdfusetitle,colorlinks,plainpages=false]{hyperref}
\usepackage[table]{xcolor} % 'table'optie is handig voor cell backgrounds
\usepackage{xltabular}    % combinatie van longtable + tabularx
\usepackage{enumitem}
\usepackage{stmaryrd}
\usepackage[section]{placeins}
\usepackage{xltabular,enumitem} 
\usepackage{pdflscape}
\usepackage{subcaption}
\usepackage{float}
\usepackage{amsmath}
\usepackage{listings}
\usepackage{xcolor}
\usepackage{graphicx}
\graphicspath{%
	{Figures/Figures-Ch1/}%
	{Figures/Figures-Ch2/}%
	{Figures/Figures-Ch3/}%
	{Figures/Figures-Ch4/}%
	{Figures/Figures-Ch5/}%
	{Figures/Figures-Ch6/}%
	{Figures/Figures-Ch7/}%
	{Figures/Figures-AppA/}%
	{Figures/Figures-AppB/}%
	{Figures/Figures-AppC/}%
}
\usepackage[utf8]{inputenc}
\DeclareUnicodeCharacter{00A0}{ }
\DeclareUnicodeCharacter{200B}{ }
\DeclareUnicodeCharacter{202F}{ }
\definecolor{codebg}{HTML}{F6F8FB}
\definecolor{codeframe}{HTML}{9FB3C8}
\definecolor{codekeyword}{HTML}{005F9E}
\definecolor{codecomment}{HTML}{3B7A45}
\definecolor{codestring}{HTML}{A33D2D}
\definecolor{codeemph}{HTML}{6D4AA2}
\definecolor{codenumber}{HTML}{687385}
\lstset{
    language=Python,
    alsoletter={_},
    basicstyle=\ttfamily\footnotesize,
    keywordstyle=\color{codekeyword}\bfseries,
    commentstyle=\color{codecomment}\itshape,
    stringstyle=\color{codestring},
    emphstyle=\color{codeemph}\bfseries,
    emph={
        active_mask,ActiveDV,rho_raw,rho_filtered,rho_effective,
        alpha,nu_tilde,nu_t_frozen,wall_distance,dJ_drho_raw,
        SA_NU_TILDE_PENALTY_ALPHA,SA_NU_TILDE_PENALTY_N,
        SA_WALL_PENALTY_ALPHA,SA_WALL_PENALTY_N
    },
    showstringspaces=false,
    columns=fullflexible,
    keepspaces=true,
    breaklines=true,
    backgroundcolor=\color{codebg},
    frame=single,
    rulecolor=\color{codeframe},
    framerule=0.6pt,
    framesep=5pt,
    xleftmargin=0.5em,
    xrightmargin=0.5em,
    numbers=left,
    numberstyle=\scriptsize\color{codenumber},
    numbersep=7pt,
    captionpos=b
}
%%%%%%%
% The lipsum package is used to generate random text.
% You never need this in a real master's thesis text!
\IfFileExists{lipsum.sty}%
 {\usepackage{lipsum}\SetLipsumDefault{11-13}}%
 {\newcommand{\lipsum}[1][11-13]{\par And some text: lipsum ##1.\par}}
%%%%%%%
\usepackage{longtable}
\setlength{\headheight}{28pt}
%%%%%%%
%\includeonly{chapter-n}
\begin{document}

\begin{preface}

This thesis marks the culmination of my Master of Science in Mechanical Engineering at KU Leuven. The journey of exploring computational fluid dynamics and mathematical optimization has been both intensely challenging and deeply rewarding. Bridging the gap between theoretical mathematics and practical engineering applications required countless hours of coding, debugging, and analyzing flow fields, but seeing the optimizer successfully carve out organic, aerodynamic fluid channels made the entire process worthwhile. 

I would like to express my sincere gratitude to my supervisors, Prof. Dr. Ir. Niels Horsten and Prof. Dr. Ir. Maarten Blommaert, for providing me with the opportunity to work on this fascinating topic. Their academic guidance, critical insights, and continuous support were instrumental in shaping the direction and rigor of this research. A very special thanks goes to my assistant-supervisor, Ir. Esteban Foronda Obando, for his day-to-day patience, technical expertise, and willingness to help me navigate the complex numerical hurdles within FEniCS.

Furthermore, I would like to thank the assessors, Prof. Dr. Ir. Elke Deckers and Ir. Olivier Renders, for their time and valuable feedback in evaluating this work. 

Finally, I want to thank my family and friends for their unwavering encouragement, patience, and perspective throughout the course of my studies. 

\end{preface}

\tableofcontents*

\begin{abstract}

As compact electronic, automotive, and energy systems become more power-dense, thermal-management hardware increasingly depends on low-loss internal-flow paths. This thesis addresses the hydraulic part of that problem: turbulent flow-only topology optimization for bends, U-bends, and diffusers before conjugate heat transfer is added. A density-based FEniCS framework is developed using Brinkman penalization, wall-resolved Spalart-Allmaras (SA) turbulence modeling, a relaxed Eikonal wall-distance calculation, and a frozen-turbulence adjoint in which the eddy-viscosity and wall-distance fields are updated in the forward solve but fixed during sensitivity assembly.

The work first defines the modeling limits of SA for curved and separated internal flows, then verifies the custom FEniCS-SA solver against Ansys Fluent for a periodic turbulent channel and a two-dimensional U-bend. These cases show qualitative and profile-level agreement sufficient for use as the forward solver in the optimization framework, while retaining the known limitations of a two-dimensional, one-equation eddy-viscosity model.

The framework is then evaluated on turbulent pipe-bend, U-bend, and diffuser benchmarks. For the completed pipe-bend verification, frozen sensitivities agree with central finite differences to a maximum coordinate discrepancy of $5.92\times10^{-5}$, and the directional Taylor check shows approximately second-order remainder convergence. Ansys Fluent validation of the extracted pipe-bend design shows moderate closure sensitivity: replacing Ansys-SA by SST $k-\omega$ increases pressure drop by about $4.5\%$ and viscous dissipation by about $8.5\%$. Under identical SST $k-\omega$ operating conditions, the turbulent design reduces pressure drop by approximately $72\%$ and viscous dissipation by approximately $47\%$ relative to the laminar-optimized reference. The U-bend is retained as the companion $180^\circ$ return-flow benchmark and is interpreted using the same validation standard. The pressure-outlet diffuser gives a more cautionary result: the turbulent SA design improves on the laminar reference under Ansys-SA, but the ranking reverses under SST $k-\omega$, showing that separated diffuser performance is strongly closure-sensitive.

The main conclusion is that frozen-SA topology optimization can produce credible low-loss turbulent internal-flow designs when its assumptions are explicit and final geometries are independently validated. The main remaining limitations are the frozen-adjoint approximation, two-dimensional suppression of secondary-flow physics, and sensitivity of final performance to geometry extraction.

\end{abstract}

% A list of figures and tables is optional
%\listoffigures
%\listoftables
% If there are only a few figures and tables, it is better to use the following:
\listoffiguresandtables


% The list of symbols is also optional.
% This list must be created manually, e.g. as follows:
\raggedbottom
\chapter{List of Abbreviations and Symbols}

\section*{Abbreviations}
\begin{flushleft}
	\small
	\renewcommand{\arraystretch}{1.0}
	\begin{tabularx}{\textwidth}{@{}p{35mm}X@{}}
		AMG & Algebraic Multigrid \\
		BiCGSTAB & Bi-Conjugate Gradient Stabilized method \\
		CAD & Computer-Aided Design \\
		CEV & Constant Eddy Viscosity \\
		CFD & Computational Fluid Dynamics \\
		CG1 & Continuous Galerkin (Linear Basis Functions) \\
		CG2 & Continuous Galerkin (Quadratic Basis Functions) \\
		CHT & Conjugate Heat Transfer \\
		DG0 & Discontinuous Galerkin (Piecewise Constant Elements) \\
		DNS & Direct Numerical Simulation \\
		DOF & Degree of Freedom \\
		FEM & Finite Element Method \\
		FD & Finite Difference \\
		FVM & Finite Volume Method \\
		GMRES & Generalized Minimal Residual method \\
		IPCS & Incremental Pressure Correction Scheme \\
		ILU & Incomplete LU factorization \\
		MMA & Method of Moving Asymptotes \\
		MUMPS & Multifrontal Massively Parallel Sparse direct Solver \\
		NS-SA & Navier-Stokes-Spalart-Allmaras coupling \\
		PDE & Partial Differential Equation \\
		RAMP & Rational Approximation of Material Properties \\
		RANS & Reynolds-Averaged Navier-Stokes \\
			SA  & Spalart-Allmaras \\
			SIMP & Solid Isotropic Material with Penalization \\
			SNES & Scalable Nonlinear Equations Solver \\
			SUPG & Streamline-Upwind/Petrov--Galerkin stabilization \\
			SST $k-\omega$ & Shear Stress Transport $k-\omega$ model \\
			TMR & NASA Turbulence Modeling Resource \\
			TO  & Topology Optimization \\
			UFL & Unified Form Language \\
			VMFL & Ansys Fluent verification benchmark case \\
		XFEM & Extended Finite Element Method \\
	\end{tabularx}
\end{flushleft}

\section*{Symbols}

\subsection*{Fluid Dynamics \& Turbulence}
\begin{flushleft}
	\renewcommand{\arraystretch}{1.1}
	\begin{tabularx}{\textwidth}{@{}p{35mm}X@{}}
		$D, D_h$ & Pipe diameter and hydraulic diameter [m] \\
		$De$ & Dean number [-] \\
		$H, L$ & Channel height and characteristic length [m] \\
		$k$ & Turbulent kinetic energy [m$^2$/s$^2$] \\
		$p, \bar{p}, p'$ & Instantaneous, mean, and fluctuating static pressure [Pa] \\
		$\Delta P$ & Macroscopic total pressure drop [Pa] \\
		$R_c$ & Radius of curvature [m] \\
		$Re$ & Bulk Reynolds number [-] \\
			$S, \tilde{S}$ & Scalar vorticity magnitude and modified SA vorticity magnitude [1/s] \\
			$t$ & Time / Pseudo-time [s] \\
			$\tau_w$ & Wall shear stress [Pa] \\
		$\mathbf{u}, \bar{u}_i$ & Velocity vector / Mean velocity component [m/s] \\
		$u'_i$ & Fluctuating velocity component [m/s] \\
		$U_{\mathrm{ref}}, U_{\mathrm{bulk}}, U_{\max}$ & Reference, bulk, and maximum velocity [m/s] \\
		$u^+$ & Non-dimensional mean velocity [-] \\
		$u_\tau$ & Friction velocity [m/s] \\
		$y, d_w$ & Distance to the nearest wall [m] \\
		$y^+$ & Non-dimensional wall distance [-] \\
		$C_f$ & Skin-friction coefficient [-] \\
		$c_{b1}, c_{b2}, c_{v1}$ & Spalart-Allmaras production and viscosity-ratio constants [-] \\
		$c_{w1}, c_{w2}, c_{w3}$ & Spalart-Allmaras wall-destruction constants [-] \\
		$f_{v1}, f_{v2}, f_w$ & Spalart-Allmaras damping and wall-destruction functions [-] \\
		$r_{\mathrm{SA}}, g_{\mathrm{SA}}$ & Spalart-Allmaras wall-destruction auxiliary variables [-] \\
		$\delta_{ij}$ & Kronecker delta [-] \\
		$\mu, \mu_t$ & Dynamic and turbulent eddy viscosity [Pa$\cdot$s] \\
		$\nu, \nu_t$ & Kinematic and turbulent eddy kinematic viscosity [m$^2$/s] \\
		$\tilde{\nu}$ & Spalart-Allmaras working variable [m$^2$/s] \\
		$\rho$ & Fluid density [kg/m$^3$] \\
		$\omega$ & Scalar vorticity in two-dimensional flow [1/s] \\
			$\Omega_{ij}$ & Vorticity tensor [1/s] \\
			$\chi$ & Spalart-Allmaras viscosity ratio, $\tilde{\nu}/\nu$ [-] \\
			$\kappa$ & von Kármán constant [-] \\
			$\sigma$ & Spalart-Allmaras turbulent Prandtl number [-] \\
		\end{tabularx}
	\end{flushleft}

\subsection*{Numerical Implementation of the Spalart-Allmaras Model}
\begin{flushleft}
	\renewcommand{\arraystretch}{1.1}
	\begin{tabularx}{\textwidth}{@{}p{35mm}X@{}}
		$G, G_0$ & Relaxed Eikonal field variable and boundary value [1/m] \\
			$h, h_{avg}$ & Finite element cell diameter and average cell size [m] \\
			$p_k$ & Kinematic pressure [m$^2$/s$^2$] \\
			$s$ & Arclength coordinate along a profile line [m] \\
			$e_i$ & Pointwise profile error or design-space basis vector, depending on context [-] \\
			$E_\infty, E_{L_2}$ & Maximum and relative $L_2$ profile discrepancies [-] \\
			$C_{\mathrm{SUPG}}$ & User-selected SUPG stabilization factor [-] \\
			$\Delta h$ & Finite-difference design perturbation [-] \\
			$\Delta t_{\tilde{\nu}}$ & Pseudo-time step used for optional SA transport damping [s] \\
			$S_i^{\mathrm{FD}}, S_i^{\mathrm{fr}}$ & Finite-difference and frozen-adjoint sensitivity of sampled design variable $i$ \\
			$\varepsilon_i^{\mathrm{fr}}$ & Relative frozen-adjoint sensitivity discrepancy [-] \\
			$||\cdot||_{L_2}$ & $L_2$ (Euclidean) norm [-] \\
			$\llbracket \cdot \rrbracket$ & Jump operator across an internal element facet [-] \\
			$\Gamma_{\mathrm{in}}, \Gamma_{\mathrm{out}}, \Gamma_D$ & Inlet, outlet, and profile-extraction boundaries \\
			$\mathbf{n}$ & Outward-pointing normal vector [-] \\
			$\mathcal{R}_{\tilde{\nu}}$ & Strong residual of the Picard-linearized SA transport equation [-] \\
			$T_h, V_h, Q_h$ & Discrete turbulence, velocity, and pressure finite-element spaces \\
			$\tau_{\tilde{\nu}}$ & SUPG stabilization parameter for SA working-variable transport [s] \\
			$v, \xi$ & Test functions for pressure and turbulence equations [-] \\
			$\alpha_u, \alpha_{\tilde{\nu}}$ & Relaxation factors for velocity and turbulence [-] \\
			$\Delta t$ & Pseudo-time step [s] \\
		$\epsilon_{tol}$ & Convergence tolerance [-] \\
		$\phi$ & General scalar field variable [-] \\
		$\sigma_w$ & Eikonal relaxation constant [-] \\
	\end{tabularx}
\end{flushleft}

\subsection*{Topology Optimization}
\begin{flushleft}
	\renewcommand{\arraystretch}{1.1}
	\begin{tabularx}{\textwidth}{@{}p{25mm}X@{}}
		$\alpha(\gamma)$ & Continuous Brinkman penalization term / inverse permeability \\
		$\alpha_{solid}, \alpha_{fluid}$ & Maximum solid and minimum fluid Brinkman penalties \\
		$\alpha_{dg}$ & Artificial diffusion scaling factor for DG0 filter [-] \\
		$\beta$ & Heaviside projection steepness parameter [-] \\
		$\gamma, \gamma_{raw}$ & Continuous and raw design variable (1 for fluid, 0 for solid) [-] \\
		$\tilde{\gamma}, \bar{\gamma}$ & Filtered and Heaviside-projected design variables [-] \\
		$\eta$ & Heaviside projection threshold [-] \\
		$\Pi_\beta$ & Heaviside projection operator [-] \\
		$\mathcal{H}$ & Helmholtz filter operator [-] \\
		$M$ & Active design mask [-] \\
		$J$ & Generic objective function used in the optimization problem [problem dependent] \\
		$J_D$ & Dissipation objective, evaluated from viscous dissipation [W, or W/m for two-dimensional unit-depth cases] \\
		$J_p$ & Average inlet-pressure objective [Pa] \\
		$\Phi_D$ & Sharp-geometry viscous dissipation metric [W, or W/m for two-dimensional unit-depth cases] \\
		$q$ & RAMP penalization tuning parameter [-] \\
		$r$ & Continuous Helmholtz filter radius [m] \\
		$r_V^{(k)}$ & Volume residual at optimization iteration $k$ [-] \\
		$v_f^{(k)}$ & Realized fluid volume fraction at optimization iteration $k$ [-] \\
		$\mathbf{s}, \mathbf{s}^*$ & Monolithic forward state and frozen-flow adjoint state vectors \\
		$V_f, V_{domain}$ & Physical fluid volume and total design domain volume [m$^3$] \\
		$V_{\mathrm{mask}}$ & Volume-mask measure [m$^3$] \\
		$V_{max}$ & Maximum allowable fluid volume fraction constraint [-] \\
		$\Omega, \partial\Omega, \Gamma_{int}$ & Computational domain, domain boundary, and internal element facets \\
		$\Omega_D, \Omega_V$ & Objective-integration and volume-integration regions \\
	\end{tabularx}
\end{flushleft}


% Now comes the main text
\flushbottom
\mainmatter

\chapter{Introduction}
\label{cha:intro}

\section{Background and Motivation}
Energy systems increasingly rely on compact flow devices that must move heat, mass, or working fluids through limited space with minimal loss. Examples include liquid-cooled electronics, battery and power-electronics thermal management, fuel-cell and electrolyzer flow fields, compact heat exchangers, and coolant distribution hardware in automotive and aerospace applications. In these systems, the thermal design is often constrained before any heat-transfer equation is considered: the fluid must first be distributed, collected, redirected, or expanded without creating an excessive pressure drop. A cooling plate or heat exchanger with excellent local heat-transfer potential can still be impractical if the hydraulic routing requires too much pumping power or if the flow distribution is poor.

This thesis is positioned within that broader energy and thermal-management context, but its actual framework is deliberately narrower. It does not solve the energy equation, does not perform conjugate heat-transfer optimization, and does not optimize temperature directly. Instead, it studies representative flow-only configurations: bends, U-bends, and diffusers whose performance is governed by pressure loss, viscous dissipation, separation, and turbulence-model behavior. These flow configurations are a necessary stepping stone toward turbulent heat-transfer optimization because reliable thermal design requires a credible and stable prediction of the hydraulic losses first.

Mathematical optimization offers a way to design such internal-flow layouts without prescribing the final channel topology in advance. Among these strategies, topology optimization has emerged as a powerful computational design tool for fluid and thermal-management problems \cite{Fawaz2022}. In a fully coupled thermal problem, the eventual objective may involve heat-transfer rate, maximum temperature, or thermal resistance. In the flow-only problem treated here, the relevant objectives are instead pressure drop and dissipated mechanical energy. This restriction is not merely a simplification for convenience: turbulent density-based topology optimization remains challenging even before the energy equation is added, because the flow solver, turbulence model, wall treatment, geometry representation, and adjoint sensitivity calculation all interact with the evolving topology.

\begin{figure}[!ht]
	\centering
	\includegraphics[width=\textwidth]{Heat-Sink-TO}
	\caption{Topology-optimized heat sink from Barroca et al.~\cite{Barroca2023}.}
	\label{fig:barroca_heat_sink}
\end{figure}

The heat-sink example in Figure~\ref{fig:barroca_heat_sink} illustrates the type of hydraulic layout that motivates the present work: the optimized fin structure guides the velocity field while the thermal field responds to that flow distribution. The central gap addressed in this thesis is therefore that density-based fluid topology optimization is far more mature for laminar flows than for turbulent internal flows. Laminar assumptions are attractive because they simplify both the governing equations and the associated adjoint sensitivities, but they are often not representative of compact engineering flow paths operating at industrial Reynolds numbers. In turbulent internal flows, the chosen Reynolds-Averaged Navier-Stokes (RANS) closure, near-wall treatment, and response to curvature or separation can influence the optimized topology itself. Before turbulent conjugate heat-transfer optimization can be treated as a reliable design tool, the underlying turbulent flow-only optimization problem must first be stabilized, verified, validated, and interpreted critically.

\section{Conceptual Overview of Density-Based Topology Optimization}
To understand the challenges of applying optimization to turbulent fluids, it is first necessary to distinguish topology optimization from traditional design methodologies. Historically, computational design relied on size optimization, which changes the thickness of predefined components, or shape optimization, which moves the boundary nodes of a predefined geometry. While effective, these methods are inherently limited by their starting topology. A shape optimization algorithm cannot spontaneously create a new hole or split a single fluid channel into two.

Density-based topology optimization, pioneered by Bendsoe and Kikuchi in 1988 for solid mechanics \cite{Bendsoe1988}, approaches the problem differently. Instead of moving boundaries, the design space is discretized into a fixed grid of finite elements. The algorithm then determines the optimal distribution of material across this grid. In fluid problems, this is achieved by assigning a continuous design variable, $\gamma \in [0, 1]$, to every element in the domain. In the convention adopted for this thesis, $\gamma = 1$ represents a pure fluid channel, while $\gamma = 0$ represents a solid, impermeable structure.

Because gradient-based optimizers require continuous and differentiable functions, this binary choice is mathematically relaxed. The design variable is allowed to take intermediate ``gray'' values between 0 and 1. To give physical meaning to these intermediate states, the fluid equations are modified through porous-medium Brinkman penalization, where intermediate densities correspond to varying levels of flow permeability. The optimization algorithm iteratively updates the design field to minimize a specific objective function, such as pressure loss or dissipated power, while satisfying constraints such as a maximum allowable fluid volume fraction. Figure~\ref{fig:to_iterative_process} summarizes this evolution: the optimizer starts from a uniform density field, passes through a smooth intermediate distribution, and converges toward a discrete fluid topology. In this introductory figure, the displayed density should be interpreted as the physical projected density used by the flow equations. Chapter~\ref{ch:topology_optimization_methodology} later distinguishes this field from the raw optimizer variable and the filtered density.

\begin{figure}[htpb]
	\centering
	\begin{subfigure}[t]{0.28\textwidth}
		\centering
		\vspace{0pt}
		
\includegraphics[height=0.13\textheight,keepaspectratio]{Initial-Density-Field}
		\caption{Initial uniform density}
		\label{fig:to_initial}
	\end{subfigure}
	\hfill
	\begin{subfigure}[t]{0.28\textwidth}
		\centering
		\vspace{0pt}
		\includegraphics[height=0.13\textheight,keepaspectratio]{Intermediate-Density-Field}
		\caption{Continuous distribution}
		\label{fig:to_intermediate}
	\end{subfigure}
	\hfill
	\begin{subfigure}[t]{0.38\textwidth}
		\centering
		\vspace{0pt}
		\includegraphics[height=0.13\textheight,keepaspectratio]{Final-Density-Field}
		\caption{Discrete topology}
		\label{fig:to_final}
	\end{subfigure}
	\caption{Density-field evolution during topology optimization.}
	\label{fig:to_iterative_process}
\end{figure}

\section{From Laminar to Turbulent Fluid Topology Optimization}
While density-based topology optimization became a standard tool in structural engineering in the late 1990s, its application to fluid dynamics is relatively recent. The optimization of internal-flow layouts remains an active and highly challenging area of research \cite{Fawaz2022}. The genesis of fluid topology optimization is widely attributed to Borrvall and Petersson (2003), who successfully applied the technique to minimize power dissipation in creeping Stokes flow \cite{Borrvall2003}. This was quickly followed by Gersborg-Hansen et al. (2005), who extended the framework to steady, laminar Navier-Stokes flows \cite{Gersborg2005}. Over the past decade, the literature surrounding laminar flow optimization expanded extensively, adapting the methodology to standard incompressible Navier-Stokes equations \cite{Deng2013, Alexandersen2020, Zhou2017}.

However, the assumption of laminar flow significantly simplifies both the governing equations and the derivation of adjoint sensitivities required for gradient-based optimization. Unfortunately, this assumption prevents the application of these methodologies to many industrial and high-performance applications, where flow velocities are high and turbulence dominates the flow physics. Only recently has the computational mechanics community begun to tackle the severe non-linearities and numerical instabilities associated with optimizing geometries directly in the turbulent regime \cite{Yoon2016, Dilgen2018, Zhao2021}.

Extending topology optimization to the turbulent regime is highly non-trivial. It requires coupling optimization algorithms with Reynolds-Averaged Navier-Stokes (RANS) turbulence models. These models introduce steep velocity gradients, highly non-linear source terms, near-wall resolution requirements, and severe numerical instabilities during the iterative optimization process. To mitigate this, gradient-based turbulent optimizations frequently rely on the ``frozen turbulence'' assumption, where the turbulent viscosity field is updated in the forward solve but treated as fixed during sensitivity assembly. While this assumption provides necessary solver stability and has driven success in industrial adjoint shape optimization for automotive aerodynamics \cite{Othmer2014}, advancing toward fully coupled turbulent sensitivities in topology optimization remains an ongoing mathematical challenge.

\begin{figure}[htpb]
	\centering
	\begin{minipage}[c]{0.37\textwidth}
		\includegraphics[width=\textwidth]{Manifold-TO-domain}
	\end{minipage}
	\hfill
	\begin{minipage}[c]{0.60\textwidth}
		\includegraphics[width=\textwidth]{Manifold-TO-solution}
	\end{minipage}
		\caption{Turbulent-flow topology optimization example. Adapted from Dilgen et al.~\cite{Dilgen2018}.}
	\label{fig:dilgen_composite}
\end{figure}

Pioneering works by Yoon (2016) \cite{Yoon2016} and Dilgen et al. (2018) \cite{Dilgen2018} have pushed the boundaries of turbulent frameworks. As seen in Figure~\ref{fig:dilgen_composite}, integrating turbulence models allows the optimizer to generate complex internal-flow layouts tailored to guide turbulent flow distributions efficiently.

\section{Objectives and Research Questions}
The main objective of this Master's thesis is to extend density-based topology optimization to internal flow configurations operating in the turbulent regime. The phrase ``flow configurations'' in the title is intentional: the work targets the hydraulic part of energy and thermal-management devices, not the complete heat-transfer device. To answer the research questions below, the thesis couples incompressible flow, RANS turbulence modeling, Brinkman topology penalization, and gradient-based optimization within a finite-element framework. By isolating the fluid mechanics from heat-transfer calculations, the thesis bounds its scope to establish a stable and verifiable baseline for future turbulent thermal design.

Rather than relying solely on black-box commercial CFD software, the core framework is implemented in \textbf{FEniCS}, an open-source computing platform for solving partial differential equations \cite{Alnaes2015}. FEniCS and UFL provide the mathematical flexibility required to write finite-element weak forms directly, verify the one-equation Spalart-Allmaras (SA) turbulence model, and assemble the residual derivatives used by the frozen-turbulence adjoint.

At this stage, it is important to distinguish two different comparisons that appear later in the thesis. A comparison between laminar-optimized and turbulent-optimized layouts is useful as supporting design context, because it shows how inertia and turbulence alter the preferred topology. However, this comparison is not treated as the principal CFD validation step. The decisive assessment instead consists of validating the final SA-optimized geometries in Ansys Fluent, first with the same SA model and then with the SST $k-\omega$ closure.

The study is organized around the following research questions:
\begin{enumerate}
	\item \textbf{How accurately does the Spalart-Allmaras turbulence model predict curved and separated internal turbulent flows, and where do its physical limitations appear?}

	Chapter~\ref{ch:turbulence_modeling} answers this question by combining turbulence-model theory, curved-pipe and separated-flow literature evidence, and Ansys Fluent comparisons on curved internal-flow benchmarks. The Dean number is used only as a rough physical scale for the bend cases, not as a universal validity threshold for SA.

	\item \textbf{Can a custom FEniCS implementation of the Spalart-Allmaras turbulence model achieve quantitative agreement with Ansys Fluent for turbulent internal-flow benchmarks?}

	Chapters~\ref{ch:fenics_implementation} and~\ref{ch:sa_verification} answer this question by presenting the weak forms, implementing the standalone SA solver, and comparing it against Ansys Fluent on a turbulent channel and a U-bend benchmark.

	\item \textbf{To what extent do SA-driven topology optimizations produce final designs whose predicted performance remains consistent under independent Ansys Fluent validation with SA and SST $k-\omega$, and how do they compare with laminar-optimized designs under identical turbulent operating conditions?}

	Chapter~\ref{ch:benchmark_results} answers this question using the pipe-bend, U-bend, and diffuser benchmarks. It reports the FEniCS optimizations, sharp-geometry Ansys-SA validation, SST $k-\omega$ validation, and laminar-versus-turbulent design comparisons. Appendix~\ref{app:C_topology_optimization_benchmarks} contains the supporting frozen-sensitivity verification.
\end{enumerate}

\section{Outline of the Thesis}
The remainder of this thesis is structured into the following six chapters:
\begin{itemize}
	\item \textbf{Chapter~\ref{ch:turbulence_modeling} (Turbulence Modeling):} Provides the theoretical background on fluid dynamics and turbulence. It discusses the Navier-Stokes equations, RANS modeling, near-wall treatment, the physical characteristics of the Spalart-Allmaras model, and the curvature and separation limits relevant to this thesis.
	\item \textbf{Chapter~\ref{ch:fenics_implementation} (Implementation of the Spalart-Allmaras Model in FEniCS):} Details the numerical implementation of the standalone SA solver in FEniCS, including weak forms, IPCS flow solves, Picard coupling, and wall-distance computation.
	\item \textbf{Chapter~\ref{ch:sa_verification} (Verification of the Spalart-Allmaras Solver for Turbulent Internal Flows):} Verifies the standalone FEniCS-SA forward solver against Ansys Fluent on turbulent channel-flow and U-bend benchmarks.
	\item \textbf{Chapter~\ref{ch:topology_optimization_methodology} (Turbulent Topology Optimization Methodology):} Introduces the density-based optimization framework. It details Brinkman momentum penalization, topology penalties, reciprocal wall-distance penalization, filtering, projection, MMA, and the frozen-turbulence adjoint.
	\item \textbf{Chapter~\ref{ch:benchmark_results} (Turbulent Topology Optimization Results):} Presents the turbulent topology optimization results for the pipe-bend and U-bend benchmarks from Bayat et al.~\cite{Bayat2026} and the diffuser benchmark from Yoon~\cite{Yoon2016}, and evaluates extracted designs with both SA and SST $k-\omega$ turbulence models. Appendix~\ref{app:C_topology_optimization_benchmarks} documents the pipe-bend sensitivity verification and numerical support material.
	\item \textbf{Chapter~\ref{ch:conclusion} (Conclusion and Future Work):} Summarizes the main findings, answers the research questions conservatively, and outlines remaining limitations and future extensions toward three-dimensional and conjugate heat-transfer applications.
\end{itemize}

\chapter{Turbulence Modeling}
\label{ch:turbulence_modeling}

\section{Introduction to Fluid Dynamics and Turbulence}
Fluid dynamics is governed by the Navier-Stokes equations, which describe the conservation of mass and momentum for a viscous continuum fluid \cite{White2006}. In this introductory text, these equations are presented using index notation (Einstein summation convention) to clearly illustrate the individual tensor components. In this convention, the indices $i$ and $j$ represent the three spatial dimensions (1, 2, and 3, corresponding to the $x$, $y$, and $z$ axes). Consequently, $x_i$ denotes the Cartesian coordinates, and $u_i$ represents the corresponding components of the velocity vector. A repeated index within a single term implies a mathematical summation over all three spatial directions. (Note: In subsequent chapters dealing with finite element implementation, this notation will shift to standard vector calculus to align with FEniCS unified form language syntax).

The fundamental principle of mass conservation is expressed by the continuity equation. For a general, compressible fluid, this is defined as:
\begin{equation}
	\frac{\partial \rho}{\partial t} + \frac{\partial (\rho u_i)}{\partial x_i} = 0
\end{equation}
where $\rho$ is the fluid density. However, for most liquid cooling applications and low-speed fluid manifolds, the fluid can be safely assumed to be incompressible. Under the incompressible flow assumption, the density $\rho$ is constant in both space and time. Therefore, the time derivative of density becomes zero, and the constant density can be factored out of the spatial derivative. Dividing the equation by $\rho$ yields the simplified, incompressible continuity equation:
\begin{equation}
	\frac{\partial u_i}{\partial x_i} = 0
	\label{eq:continuity}
\end{equation}

Similarly, for an incompressible, Newtonian fluid with a constant dynamic viscosity $\mu$, the conservation of momentum is expressed as:
\begin{equation}
	\rho \left( \frac{\partial u_i}{\partial t} + u_j \frac{\partial u_i}{\partial x_j} \right) = -\frac{\partial p}{\partial x_i} + \frac{\partial}{\partial x_j} \left( \mu \frac{\partial u_i}{\partial x_j} \right)
	\label{eq:navier_stokes}
\end{equation}
where $p$ represents the static pressure of the fluid. In principle, solving Equation \ref{eq:navier_stokes} directly yields the exact flow field. This approach is known as Direct Numerical Simulation (DNS). However, as the Reynolds number ($Re$) increases, the flow transitions into a chaotic turbulent state characterized by a broad spectrum of spatial and temporal scales \cite{Pope2000}. Resolving the smallest turbulent eddies requires extremely fine spatial grids. For industrial applications, including the design of electronics cooling manifolds, DNS is computationally prohibitive. Therefore, engineers rely on turbulence modeling to predict the averaged effects of these fluctuations without resolving the complete turbulent spectrum \cite{Versteeg2007, Altair}.

\section{Reynolds-Averaged Navier-Stokes (RANS) Equations}
The most common approach to turbulence modeling in engineering is the Reynolds-Averaged Navier-Stokes (RANS) framework, based on the theories originally introduced by Osborne Reynolds \cite{Pope2000}. The core principle of RANS is the Reynolds decomposition. This mathematical operation separates the instantaneous flow variables into a time-averaged mean component ($\bar{u}_i$, $\bar{p}$) and a fluctuating component ($u'_i$, $p'$):

\begin{equation}
	u_i = \bar{u}_i + u'_i \quad \text{and} \quad p = \bar{p} + p'
\end{equation}

Substituting these decomposed variables into the standard Navier-Stokes equations and taking the time average yields the RANS momentum equations:

\begin{equation}
	\rho \left( \frac{\partial \bar{u}_i}{\partial t} + \bar{u}_j \frac{\partial \bar{u}_i}{\partial x_j} \right) = -\frac{\partial \bar{p}}{\partial x_i} + \frac{\partial}{\partial x_j} \left( \mu \frac{\partial \bar{u}_i}{\partial x_j} - \rho \overline{u'_i u'_j} \right)
	\label{eq:rans_momentum}
\end{equation}

The averaging process introduces a new term, $-\rho \overline{u'_i u'_j}$, known as the Reynolds stress tensor. This tensor represents the momentum transfer caused by turbulent fluctuations. The introduction of this term creates a fundamental mathematical closure problem because the RANS equations now contain more unknowns than available physical equations \cite{Versteeg2007}.

\subsection{The Boussinesq Hypothesis and the Isotropic Assumption}
To close the system and solve for the mean flow, most RANS models rely on the Boussinesq eddy viscosity hypothesis \cite{Pope2000}. This hypothesis assumes that the turbulent momentum transfer is analogous to molecular momentum transfer. It relates the Reynolds stresses linearly to the mean velocity gradients using a proportionality constant called the turbulent eddy viscosity ($\mu_t$):

\begin{equation}
	-\rho \overline{u'_i u'_j} = \mu_t \left( \frac{\partial \bar{u}_i}{\partial x_j} + \frac{\partial \bar{u}_j}{\partial x_i} \right) - \frac{2}{3} \rho k \delta_{ij}
	\label{eq:boussinesq}
\end{equation}
where $k$ is the turbulent kinetic energy and $\delta_{ij}$ is the Kronecker delta. The primary task of any Boussinesq-based turbulence model is to calculate this local eddy viscosity. It is crucial to highlight a fundamental physical limitation of Equation \ref{eq:boussinesq}: the Boussinesq hypothesis inherently assumes that $\mu_t$ is an \textit{isotropic} scalar quantity. It mathematically enforces that turbulence behaves identically in all spatial directions and aligns perfectly with the mean strain rate. While this assumption holds for simple shear flows, it systematically fails in predicting anisotropic phenomena, such as the strong secondary flows induced by streamline curvature.

\section{Boundary Layer Theory and Near-Wall Treatment}
Before selecting a specific turbulence equation to calculate $\mu_t$, the physical structure of a turbulent boundary layer must be considered, as the model must accurately resolve these near-wall regions. The boundary layer is traditionally divided into three distinct regions based on the non-dimensional wall distance ($y^+$) and non-dimensional velocity ($u^+$):
\begin{equation}
	y^+ = \frac{y u_\tau}{\nu}, \quad u^+ = \frac{\bar{u}}{u_\tau}
\end{equation}
where $y$ is the absolute distance from the wall, $\nu$ is the kinematic viscosity ($\mu/\rho$), and $u_\tau = \sqrt{\tau_w/\rho}$ is the friction velocity derived from the wall shear stress ($\tau_w$) \cite{Pope2000}.
\begin{enumerate}
	\item \textbf{The Viscous Sublayer ($y^+ < 5$):} Immediately adjacent to the wall, turbulent fluctuations are damped out by kinematic constraints. Viscous shear dominates, and the velocity profile is strictly linear ($u^+ = y^+$).
	\item \textbf{The Buffer Layer ($5 < y^+ < 30$):} A transition region where both viscous and turbulent shear stresses are of equal magnitude.
	\item \textbf{The Log-Law Region ($y^+ > 30$):} Farther from the wall, turbulence dominates, and the velocity follows a logarithmic distribution predicted by the von Kármán constant \cite{Versteeg2007}.
\end{enumerate}

High-Reynolds models (like the standard $k-\epsilon$ model) use empirical "wall functions" to bridge the gap between the wall and the log-law region, requiring the first computational mesh element to be placed at $y^+ > 30$. However, for topology optimization, where solid-fluid boundaries continuously evolve during the optimization loop, empirical wall functions introduce severe mathematical discontinuities and numerical instabilities. Therefore, a "wall-resolved" model that integrates the governing equations all the way through the viscous sublayer ($y^+ \approx 1$) is strictly required for robust gradient-based optimization.

\section{The Spalart-Allmaras One-Equation Model}
While two-equation models (such as $k-\epsilon$ or $k-\omega$) solve separate transport equations for the turbulent kinetic energy and a turbulent length scale, they introduce highly non-linear source terms. These terms frequently cause numerical instabilities during the adjoint sensitivity analysis required for topology optimization. To achieve a trade-off between physical accuracy and computational efficiency, this thesis utilizes the Spalart-Allmaras (SA) model, originally developed in 1992 for aerodynamic flows with adverse pressure gradients \cite{Spalart1992}. 

The SA model is a one-equation model that solves a single transport equation directly for a working variable, $\tilde{\nu}$, which is subsequently linked to the true kinematic eddy viscosity ($\nu_t = \mu_t / \rho$) \cite{Spalart1994}. The transport equation for the standard SA model is given by:

\begin{equation}
	\frac{\partial \tilde{\nu}}{\partial t} + \bar{u}_j \frac{\partial \tilde{\nu}}{\partial x_j} = c_{b1} \tilde{S} \tilde{\nu} - c_{w1} f_w \left( \frac{\tilde{\nu}}{d} \right)^2 + \frac{1}{\sigma} \left[ \frac{\partial}{\partial x_j} \left( (\nu + \tilde{\nu}) \frac{\partial \tilde{\nu}}{\partial x_j} \right) + c_{b2} \left( \frac{\partial \tilde{\nu}}{\partial x_j} \right)^2 \right]
	\label{eq:sa_transport}
\end{equation}

The terms on the right-hand side represent the production, destruction, and diffusion of the turbulent working variable, respectively. Crucially, the production term relies on a modified vorticity scalar, $\tilde{S}$. This indicates that turbulence generation in the SA model is fundamentally driven by fluid shear, a physical characteristic that introduces unique boundary-condition sensitivities during advanced adjoint calculations (discussed further in Chapter~\ref{ch:topology_optimization_methodology}). 

Furthermore, the destruction term is explicitly a function of $d$, the absolute distance to the nearest wall. This makes the SA model a highly effective "wall-resolved" low-Reynolds model. It is inherently sensitive to near-wall damping effects and integrates smoothly through the viscous sublayer without requiring complex empirical wall functions or a second length-scale equation \cite{NasaTurb}. Detailed derivations of the model constants and wall damping functions are deferred to Chapter~\ref{ch:fenics_implementation}, where the numerical implementation in FEniCS is discussed.

\section{Evaluation of Curvature Effects and the Dean Number}
\label{sec:curvature_effects}
Despite its computational advantages for topology optimization, the standard SA model relies on a single length scale tied to the nearest wall distance, making it mathematically blind to the anisotropic physics of streamline curvature. When fluid passes through a bend, centrifugal forces induce secondary flows perpendicular to the main flow direction. To quantify the strength of these secondary flows, the dimensionless Dean number ($De$) is utilized:
\begin{equation}
	De = Re_d \sqrt{\frac{d}{2R_c}}
\end{equation}
where $d$ is the characteristic length of the cross-section (e.g., the pipe diameter $D$ or hydraulic diameter $D_h$), $R_c$ is the radius of curvature, and $Re_d$ is the bulk Reynolds number based strictly on that exact same characteristic length. A higher Dean number indicates a stronger cross-sectional influence of centrifugal forces. 

A review of the literature demonstrates that the predictive capability of the SA model is heavily dependent on this Dean number. For mildly curved flows, the SA model performs relatively well. Apalowo \cite{Apalowo2022} evaluated a $90^\circ$ circular pipe bend with a Dean number of $De \approx 9,000$ (operating at $Re_D \approx 34,000$, based on the pipe diameter). Under these moderate conditions, where the bend radius was seven times larger than the pipe diameter ($R_c = 7D$), the SA model provided a highly acceptable approximation of the streamwise velocity profiles compared to higher-order models, successfully capturing the primary flow physics.

\begin{figure}[htpb]
	\centering
	\includegraphics[width=0.8\textwidth]{example-image}
	\caption{Streamwise velocity profiles in a $90^\circ$ pipe bend ($R_c = 7D$, $Re_D \approx 34,000$, $De \approx 9,000$). The SA model performs relatively well, closely matching the experimental data and two-equation models for mild curvature. Reprinted from Apalowo \cite{Apalowo2022}.}
	\label{fig:apalowo_profiles}
\end{figure}

However, as the curvature tightens and the Dean number increases, the limitations of the isotropic Boussinesq assumption become severe. Rumsey et al. \cite{Rumsey2000} evaluated the internal flow through a strongly curved $180^\circ$ turnaround duct (the Monson U-Duct). While this benchmark operated at a higher bulk Reynolds number ($Re_H = 100,000$, based on the duct height $H$), the extremely tight geometry ($R_c \approx 0.83H$) generated massive centrifugal forces and an extreme equivalent Dean number. 

It is critical to note that the Spalart-Allmaras model was fundamentally designed for high-Reynolds aerospace applications and remains mathematically stable well beyond $Re = 100,000$ in straight or mildly curved flows. Therefore, the subsequent failure of the model in this benchmark cannot be attributed to inertial scaling, but is strictly driven by the extreme geometric curvature. Under these severe centrifugal conditions, the uncorrected SA model drastically overpredicted the production of eddy viscosity on the inner radius of the bend. This artificial excess of viscosity suppressed the adverse pressure gradients, causing the standard SA formulation to completely fail at predicting the massive flow separation bubble that physically occurs on the inner convex wall (Figure~\ref{fig:rumsey_uduct}).

\begin{figure}[htpb]
	\centering
	\includegraphics[width=0.7\textwidth]{example-image} % Replace with an image from Rumsey 2000
	\caption{Turbulence model predictions in a strongly curved internal U-duct ($R_c \approx 0.83H$, $Re_H = 100,000$, $De \approx 109,500$). Due to the extreme curvature, the standard Spalart-Allmaras model overpredicts eddy viscosity and fails to capture the inner-wall separation. Adapted from Rumsey et al. \cite{Rumsey2000}.}
	\label{fig:rumsey_uduct}
\end{figure}

While the literature clearly establishes that the SA model is viable at $De \approx 9,000$ and fails catastrophically in extreme bends, the specific operational threshold between these extremes must be defined to justify its use in topology optimization. Therefore, to explicitly isolate the threshold where the model breaks down, rigorous benchmark studies were subsequently conducted within Ansys Fluent.

\section{Benchmark Studies in Ansys Fluent}
\label{sec:ansys_benchmarks}
To investigate the exact physical breaking point of the SA model, benchmark simulations were performed using \textbf{Ansys Fluent}. The study focused on two distinct, fully three-dimensional (3D) internal flow geometries: a $180^\circ$ U-bend and a tighter $90^\circ$ pipe bend. To ensure a standardized and rigorously tested numerical baseline, the computational meshes, fluid properties, and boundary conditions for these simulations were adopted directly from the official Ansys Fluid Dynamics Verification Manual \cite{AnsysVerification2024}. Specifically, the $90^\circ$ bend corresponds to verification case VMFL030, while the 3D U-bend corresponds to case VMFL048. 

These cases evaluate industrial flow regimes ($Re_D > 40,000$) to ensure the turbulence model is rigorously verified against dominant inertial physics before being implemented into the custom FEniCS solver. To evaluate the accuracy of the one-equation model, the results were compared directly against the $k-\omega$ SST model, an industry standard for accurately predicting adverse pressure gradients and flow separation.
\subsection{Mild Curvature: The U-Bend Case}
The first benchmark evaluated a 3D U-bend geometry characterized by a moderate curvature ratio of $R_c \approx 4.46D$. Operating at a bulk Reynolds number of $Re_D \approx 44,700$, this geometry resulted in a Dean number of $De \approx 15,000$. Consistent with the findings of Apalowo for moderate bends, the SA model performed well under these conditions. As shown in Figure~\ref{fig:ansys_ubend_comparison}(b), the velocity contours generated by the SA model closely matched those of the high-fidelity $k-\omega$ SST model (Figure~\ref{fig:ansys_ubend_comparison}(a)), accurately capturing the bulk flow distribution without severe artificial diffusion.

\begin{figure}[htpb]
	\centering
	\begin{subfigure}[b]{0.48\textwidth}
		\centering
		\includegraphics[width=\textwidth]{example-image}
		\caption{$k-\omega$ SST Model}
		\label{fig:ubend_komega}
	\end{subfigure}
	\hfill
	\begin{subfigure}[b]{0.48\textwidth}
		\centering
		\includegraphics[width=\textwidth]{example-image}
		\caption{Spalart-Allmaras Model}
		\label{fig:ubend_sa}
	\end{subfigure}
	\caption{Comparison of velocity contours in a $180^\circ$ U-bend ($R_c \approx 4.46D$, $Re_D \approx 44,700$, $De \approx 15,000$). The SA model performs well, closely matching the high-fidelity prediction.}
	\label{fig:ansys_ubend_comparison}
\end{figure}

\subsection{Strong Curvature: The $90^\circ$ Bend Case}
The second benchmark evaluated a tighter $90^\circ$ pipe bend characterized by a severe curvature ratio of $R_c = 2.8D$. Despite operating at a very similar bulk Reynolds number of $Re_D \approx 43,000$ compared to the U-bend, this significantly tighter bend radius pushed the Dean number up to $De \approx 18,000$. Because the bulk flow speed and Reynolds number remain effectively constant between these two benchmark cases, this comparison successfully isolates the geometric effect of streamline curvature. It clearly demonstrates that the subsequent failure of the SA model is driven purely by the tighter bend radius, which elevates the Dean number and amplifies anisotropic centrifugal effects.

\begin{figure}[htpb]
	\centering
	\begin{subfigure}[b]{0.48\textwidth}
		\centering
		\includegraphics[width=\textwidth]{example-image}
		\caption{$k-\omega$ SST Model}
		\label{fig:90bend_komega}
	\end{subfigure}
	\hfill
	\begin{subfigure}[b]{0.48\textwidth}
		\centering
		\includegraphics[width=\textwidth]{example-image}
		\caption{Spalart-Allmaras Model}
		\label{fig:90bend_sa}
	\end{subfigure}
	\caption{Comparison of velocity contours in a $90^\circ$ pipe bend ($R_c = 2.8D$, $Re_D \approx 43,000$, $De \approx 18,000$). The SA model over-smooths the velocity gradients compared to the $k-\omega$ SST model.}
	\label{fig:ansys_90bend_comparison}
\end{figure}

\section{Conclusion}
This chapter directly addresses the first research question regarding the predictive capabilities and physical limitations of the Spalart-Allmaras model. The theoretical breakdown and Ansys benchmark studies confirm that for internal geometries with moderate curvature (up to $De \approx 15,000$), the SA model performs accurately and successfully captures the primary flow physics. Furthermore, because it is a wall-resolved model, it avoids the numerical pitfalls of empirical wall functions, making it exceptionally well-suited for boundary-evolving topology optimization. 

However, its physical limitation lies in the isotropic Boussinesq assumption. When streamline curvature induces strong anisotropic secondary flows (heuristically observed to cause model breakdown beyond a threshold of $De \approx 18,000$), the uncompensated SA model systematically overpredicts eddy viscosity and fails to capture flow separation. Consequently, while the SA model provides the necessary numerical stability to drive the optimization algorithms in FEniCS, the final topological designs must be critically verified. Any complex manifold generated by the SA optimizer must be subsequently evaluated in commercial software using a higher-fidelity two-equation model to confirm its true aerodynamic performance.
\chapter{Implementation and Verification of the Spalart-Allmaras Model in FEniCS}
\label{ch:fenics_implementation}

\section{Introduction to FEniCS and the Finite Element Method}
With the theoretical foundation of the Spalart-Allmaras (SA) model established in Chapter~\ref{ch:turbulence_modeling}, this chapter details its numerical translation into a working computational solver. The numerical simulations and eventual topology optimization in this thesis are built upon FEniCS, an open-source computing platform for solving partial differential equations (PDEs). Unlike commercial computational fluid dynamics (CFD) software such as Ansys, which typically relies on the Finite Volume Method (FVM) and provides a "black box" graphical interface, FEniCS uses the Finite Element Method (FEM) and requires the user to explicitly define the mathematical formulation of the problem \cite{Logg2012}. 

To solve a PDE in FEniCS, the governing equations (the "strong form") must be mathematically transformed into a "weak formulation" (or variational form). FEniCS utilizes the Unified Form Language (UFL) to express these weak forms in Python code that closely mirrors the mathematical notation \cite{Langtangen2016}. This high-level mathematical scripting is particularly advantageous for topology optimization, as it allows for the automated derivation of the adjoint equations required for gradient-based optimization algorithms. 

To ensure complete transparency and reproducibility, the entire custom codebase developed for this thesis is available in the public GitHub repository \texttt{TurbulentTO}. To maintain modularity, the repository is divided into two primary directories. The \texttt{TurbulenceModels} directory contains the standalone Spalart-Allmaras fluid dynamics engine and the verification benchmarks discussed throughout this chapter. Conversely, the \texttt{FluidTO} directory houses the optimization algorithms, penalty formulations, and objective functions that will be detailed in Chapters~\ref{ch:topology_optimization_methodology} and \ref{ch:benchmark_results}.

\subsection{Derivation of a Weak Formulation}
Before diving into the complex turbulence equations, it is instructive to demonstrate how a continuous PDE is translated into a discrete FEniCS formulation. To illustrate this methodology, consider the Poisson equation for pressure, which is a fundamental component of incompressible fluid solvers. The strong form of the Poisson equation is given by:
\begin{equation}
	-\nabla^2 p = f \quad \text{in} \, \Omega
\end{equation}
where $p$ is the pressure field, $f$ is a source term (derived from the velocity divergence), and $\Omega$ is the computational domain. To obtain the weak form, the equation is multiplied by a sufficiently smooth test function $v$ and integrated over the entire domain $\Omega$, following standard Galerkin finite element procedures \cite{Logg2012, Langtangen2016}:
\begin{equation}
	-\int_{\Omega} (\nabla^2 p) v \, dx = \int_{\Omega} f v \, dx
\end{equation}
Applying integration by parts (Green's first identity) to the left-hand side yields:
\begin{equation}
	\int_{\Omega} \nabla p \cdot \nabla v \, dx - \int_{\partial \Omega} (\nabla p \cdot \mathbf{n}) v \, ds = \int_{\Omega} f v \, dx
\end{equation}
where $\partial \Omega$ is the boundary of the domain and $\mathbf{n}$ is the outward-pointing normal vector. This integral formulation fundamentally lowers the continuity requirements of the solution, requiring only first-order derivatives to be square-integrable. Furthermore, Neumann boundary conditions (e.g., a prescribed pressure gradient at an outlet) can be naturally substituted directly into the boundary integral $\int_{\partial \Omega} (\nabla p \cdot \mathbf{n}) v \, ds$. This variational approach forms the foundation of all fluid and turbulence equations implemented in this thesis.

\section{Numerical Solution Strategy for Fluid Dynamics}
Solving the full Reynolds-Averaged Navier-Stokes (RANS) equations is numerically challenging due to the non-linear convective terms, the saddle-point nature of the coupled velocity-pressure system, and the intense cross-coupling with the turbulence model. To achieve a stable and efficient steady-state solution in FEniCS that seamlessly translates into the topology optimization framework, the governing equations are rigorously decoupled.

\subsection{Decoupling the Physics: The Picard Iteration Scheme}
Because the momentum equation relies on the turbulent eddy viscosity ($\nu_t$), and the turbulence transport equation relies on the convective velocity ($\mathbf{u}$), solving the entire fluid-turbulence system simultaneously in a massive monolithic matrix frequently leads to severe ill-conditioning and solver divergence. 

To resolve this, the forward physics are decoupled using a \textbf{Picard iteration scheme} (successive substitution), a highly robust approach for segregated RANS solvers \cite{Ferziger2012, Versteeg2007}. During each Picard step, the flow solver advances using the eddy viscosity from the previous iteration. The newly updated velocity field is then passed to the Spalart-Allmaras solver to calculate a new eddy viscosity field. This alternating sequence forms the outer loop of the solver, iteratively updating the decoupled equations until the fully coupled physical system reaches a steady-state equilibrium.

\subsection{Incremental Pressure Correction Scheme (IPCS)}
Within the Picard loop, the Navier-Stokes equations themselves must be solved. Rather than solving for velocity and pressure simultaneously, this implementation utilizes an Incremental Pressure Correction Scheme (IPCS). This algorithmic splitting is fundamentally based on the fractional-step projection methods originally pioneered by Chorin \cite{Chorin1968} and Temam \cite{Temam1969}. To nurse the non-linear convective momentum terms to a stable steady state, a pseudo-time stepping approach ($\Delta t$) is employed strictly for the flow variables. The IPCS algorithm splits the fluid equations into three sequential steps per pseudo-time step:
\begin{enumerate}
	\item \textbf{Tentative Velocity Step:} The momentum equation is solved using the pressure field from the previous iteration ($p^n$) to compute an intermediate, non-divergence-free velocity field ($\mathbf{u}^*$). Because this is a RANS framework, the diffusion term utilizes the effective kinematic viscosity, $(\nu + \nu_t)$. Treating the convective term explicitly and the diffusive term implicitly, the semi-discrete equation is:
	\begin{equation}
		\frac{\mathbf{u}^* - \mathbf{u}^n}{\Delta t} + (\mathbf{u}^n \cdot \nabla) \mathbf{u}^* - \nabla \cdot \left( (\nu + \nu_t^n) \nabla \mathbf{u}^* \right) = -\nabla p^n
	\end{equation}
	
	\item \textbf{Pressure Correction Step:} A Poisson equation is solved to find the pressure variation that enforces the continuity constraint ($\nabla \cdot \mathbf{u}^{n+1} = 0$). The equation for the new pressure field ($p^{n+1}$) is derived as:
	\begin{equation}
		\nabla^2 p^{n+1} = \nabla^2 p^n + \frac{1}{\Delta t} \nabla \cdot \mathbf{u}^*
	\end{equation}
	
	\item \textbf{Velocity Update Step:} The tentative velocity is strictly corrected using the gradient of the pressure difference to yield a mass-conserving velocity field ($\mathbf{u}^{n+1}$) for the current iteration:
	\begin{equation}
		\mathbf{u}^{n+1} = \mathbf{u}^* - \Delta t \nabla (p^{n+1} - p^n)
	\end{equation}
\end{enumerate}
The simulation is deemed fully converged to a steady state when the $L_2$ norm of the residual difference between successive Picard iterations falls below a strictly defined tolerance limit ($\epsilon_{tol}$) for both the fluid and turbulence fields.

\section{Detailed Formulation of the Spalart-Allmaras Model}
As concluded in Chapter~\ref{ch:turbulence_modeling}, the Spalart-Allmaras (SA) one-equation model was selected for its balance of physical accuracy and numerical stability. Unlike the transient flow equations, the SA transport equation is solved in its purely steady-state form within the Picard loop. The steady transport equation for the working variable $\tilde{\nu}$ is given by:
\begin{equation}
	\mathbf{u} \cdot \nabla \tilde{\nu} = c_{b1} \tilde{S} \tilde{\nu} - c_{w1} f_w \left( \frac{\tilde{\nu}}{y} \right)^2 + \frac{1}{\sigma} \left[ \nabla \cdot \left( (\nu + \tilde{\nu}) \nabla \tilde{\nu} \right) + c_{b2} (\nabla \tilde{\nu})^2 \right]
	\label{eq:sa_steady_transport}
\end{equation}
where $\nu$ is the molecular kinematic viscosity, $\sigma = 2/3$ is the turbulent Prandtl number, $c_{b1}$ and $c_{w1}$ are calibrated closure constants, $\tilde{S}$ is a modified vorticity magnitude, $f_w$ is the wall-damping function, and $y$ is the physical distance to the nearest solid boundary \cite{Spalart1994}.

\subsection{Closure Terms and Calibration Constants}
The Spalart-Allmaras model relies on several interconnected closure functions to accurately capture boundary layer physics. The production term relies on a modified vorticity magnitude, $\tilde{S}$, which prevents the destruction of turbulence in regions where the vorticity goes to zero, while maintaining correct scaling near the wall \cite{Spalart1994}. It is defined as:
\begin{equation}
	\tilde{S} = S + \frac{\tilde{\nu}}{\kappa^2 y^2} f_{v2}
\end{equation}
where $S = \sqrt{2 \Omega_{ij}\Omega_{ij}}$ is the scalar magnitude of the vorticity tensor, $\kappa$ is the von Kármán constant, and $y$ is the distance to the nearest wall. The function $f_{v2}$ is a near-wall modifier:
\begin{equation}
	f_{v2} = 1 - \frac{\chi}{1 + \chi f_{v1}} \quad \text{where} \quad \chi = \frac{\tilde{\nu}}{\nu} \quad \text{and} \quad f_{v1} = \frac{\chi^3}{\chi^3 + c_{v1}^3}
\end{equation}

Similarly, the wall destruction term is modulated by the non-dimensional function $f_w$, which accelerates the decay of turbulent viscosity inside the viscous sublayer:
\begin{equation}
	f_w = g \left( \frac{1 + c_{w3}^6}{g^6 + c_{w3}^6} \right)^{1/6}
\end{equation}
where the intermediate variables $g$ and $r$ are evaluated as:
\begin{equation}
	g = r + c_{w2}(r^6 - r) \quad \text{and} \quad r = \min \left( \frac{\tilde{\nu}}{\tilde{S} \kappa^2 y^2}, 10 \right)
\end{equation}
The upper bound of $10$ on $r$ ensures numerical stability during solver initialization when $\tilde{S}$ may momentarily approach zero. 

It should be explicitly noted that the original formulation by Spalart and Allmaras \cite{Spalart1994} includes additional geometric trip functions ($f_{t1}$ and $f_{t2}$) designed to manually trigger the transition from laminar to turbulent flow at specific user-defined coordinates. However, following the topology optimization methodology established by Yoon \cite{Yoon2016}, these transition functions are deliberately omitted from this thesis. Because the topological boundaries are unknown *a priori* and evolve during optimization, tracking discrete trip points is mathematically unfeasible; therefore, the model assumes a fully turbulent flow regime.

To close the system, the standard empirical calibration constants formulated by Spalart and Allmaras are utilized throughout the FEniCS implementation \cite{NasaTurb}:
\begin{align*}
	c_{b1} &= 0.1355 & \sigma &= 2/3 & c_{b2} &= 0.622 \\
	\kappa &= 0.4187 & c_{w2} &= 0.3 & c_{w3} &= 2.0 \\
	c_{v1} &= 7.1 & c_{w1} &= \frac{c_{b1}}{\kappa^2} + \frac{1 + c_{b2}}{\sigma}
\end{align*}

\subsection{Wall-Distance Calculation via the Relaxed Eikonal Equation}
A major computational challenge in the SA model is the calculation of the wall distance $y$. In a standard static CFD simulation, $y$ is purely geometric and can be pre-computed. However, to ensure this verification solver translates seamlessly into the topology optimization code (where solid boundaries evolve continuously), the implementation abandons discrete geometric searches in favor of a PDE-based approach. 

The true geometric wall distance satisfies the standard Eikonal equation, $|\nabla y| = 1$. However, this formulation contains singularities at the medial axes of the domain. To guarantee continuous differentiability—a strict requirement for later adjoint derivations—a relaxed Eikonal equation proposed by Yoon \cite{Yoon2016} is solved:
\begin{equation}
	\nabla G \cdot \nabla G + \sigma_w G (\nabla \cdot \nabla G) = (1 + 2\sigma_w) G^4
\end{equation}
Here, $G$ is an intermediate, non-dimensional field variable. The parameter $\sigma_w$ is a relaxation constant (typically set to $0.1$) that controls the smoothness of the field. The PDE is subjected to a Dirichlet boundary condition at the physical walls, $G = G_0$, where $G_0$ is an arbitrarily large constant (e.g., $G_0 = 20.0$). Once the field $G$ is solved over the domain, the physical wall distance $y$ is recovered algebraically at every node using the inverse relationship:
\begin{equation}
	y = \frac{1}{G} - \frac{1}{G_0}
\end{equation}

This system is implemented in the core solver suite within the \texttt{Utilities.py} script. As shown in Listing~\ref{lst:eikonal_fenics}, once the non-linear variational problem for $G$ is solved, the physical distance $y$ is recovered and strictly bounded from below to prevent numerical singularities.

\begin{lstlisting}[language=Python, caption={Implementation of Yoon's relaxed Eikonal equation for wall-distance calculation, extracted from \texttt{Utilities.py}.}, label={lst:eikonal_fenics}]
	# Weak form corresponding to Yoon 2016 Eq. (19):
	F = (
	(Constant(1.0) - sigma_w_const) * inner(grad(G), grad(G)) * z * dx
	- sigma_w_const * G * inner(grad(G), grad(z)) * dx
	- (Constant(1.0) + Constant(2.0) * sigma_w_const) * G**4 * z * dx
	)
	problem = NonlinearVariationalProblem(F, G, bcs=wall_bcs, J=derivative(F, G))
	solver = NonlinearVariationalSolver(problem)
	solver.solve()
	
	# Bound G to prevent division by zero
	G.vector()[:] = np.maximum(G.vector()[:], float(g_floor))
	
	# Recover physical wall distance (y)
	y = project(Constant(1.0) / G - Constant(1.0) / g0_const, Space)
	return bound_from_bellow(y, 0.0)
\end{lstlisting}

\section{Variational Formulation and FEniCS Implementation}
The standard Galerkin weak formulation is constructed by multiplying the steady transport Equation~\ref{eq:sa_steady_transport} by a test function $\xi$ and applying integration by parts to the diffusive terms. Because the equation is highly non-linear, robust convergence is heavily dependent on the quality of the mesh and the stability of the external velocity field provided by the Picard outer loop.

The core SA physics, boundary conditions, continuous Eikonal formulations, and the overarching Picard iteration scheme were derived independently for this thesis. However, the underlying FEniCS solver architecture was structurally inspired by an open-source $k-\epsilon$ repository developed by Marcibál \cite{Marcibal2024, joove_keps}, specifically regarding the procedural execution of the variational forms and the implementation of the Incremental Pressure Correction Scheme (IPCS). The explicit translation of the SA mathematical formulation is handled by the \texttt{SpalartAllmarasSteady} class within the \texttt{TurbulenceModel\_SpalartAllmaras.py} script. As demonstrated in Listing~\ref{lst:sa_fenics}, the Unified Form Language (UFL) mirrors the analytical weak form precisely, utilizing the most recent convective velocity field (\texttt{external\_velocity}) passed from the IPCS flow solver.

\begin{lstlisting}[language=Python, caption={FEniCS implementation of the steady-state SA standard Galerkin weak formulation. Python dictionary lookups have been simplified for readability.}, label={lst:sa_fenics}]
	def construct_forms(self, external_velocity):
	self._construct_turbulent_quantities(external_velocity)
	sigma = Constant(2.0 / 3.0)
	
	# Steady-state SA transport equation weak form
	FNT = (
	dot(external_velocity, nabla_grad(self._nu_tilde)) * self._xi * self._dx
	+ inner((self._nu_laminar + self._nu_tilde0) / sigma * grad(self._nu_tilde), grad(self._xi)) * self._dx
	+ self._react_nt * self._nu_tilde * self._xi * self._dx
	- self._source_nt * self._xi * self._dx
	)
	
	self._a_nt = lhs(FNT)
	self._l_nt = rhs(FNT)
\end{lstlisting}

\section{General Solver Configuration and Numerical Parameters}
Before detailing the specific verification benchmarks, it is necessary to establish the general numerical configuration of the FEniCS solver. Both the channel flow and the U-bend verifications were simulated using the segregated algorithmic structure detailed earlier: an outer Picard fixed-point loop couples the fluid momentum equations to the Spalart-Allmaras transport equation, while an inner IPCS resolves the velocity-pressure subsystem via pseudo-time stepping. 

Because the U-bend domain introduces severe geometric curvature, strong convective transport, and adverse pressure gradients, it represents a significantly more demanding numerical challenge than a fully developed, unidirectional channel flow. Consequently, the U-bend simulations necessitated stricter numerical controls to maintain stability, including a heavily refined IPCS pseudo-time step ($\Delta t = 1.0 \times 10^{-4}$ s) and much tighter under-relaxation of the turbulence field updates ($\alpha_{\tilde{\nu}} = 0.01$). To ensure strict scientific reproducibility without disrupting the narrative flow of this verification, the comprehensive matrices detailing all physical boundary conditions, geometric dimensions, and exact numerical solver parameters for both benchmarks have been compiled and are documented extensively in Sections~\ref{sec:app_channel} and \ref{sec:app_ubend} of Appendix~\ref{app:baseline_configurations}.

\section{Verification Benchmark 1: Channel Flow}
To verify the fundamental correctness of the custom fluid and turbulence solvers, the code was first tested on a standard 2D turbulent channel flow. To cleanly isolate and verify turbulence production, wall damping, and Picard coupling without the artifacts of a developing boundary layer, this benchmark was strictly formulated as a pressure-driven periodic channel rather than a conventional inlet-outlet duct. 

The velocity field ($\mathbf{u}$) and Spalart-Allmaras working variable ($\tilde{\nu}$) were constrained periodically between the streamwise boundaries, while the top and bottom boundaries were treated as standard no-slip walls ($\mathbf{u} = 0, \tilde{\nu} = 0$). Flow was driven entirely by a fixed pressure drop from the upstream plane ($p=2$ Pa) to the downstream plane ($p=0$ Pa). Based on a reference bulk velocity of $U_{\mathrm{ref}}=18.5$ m/s and utilizing the full channel height ($H=2.0$ m) as the characteristic length scale, the solver was targeted to resolve a fully turbulent flow regime at a channel Reynolds number of $Re_H \approx 20,300$. By explicitly defining the Reynolds number based on the channel height rather than the mathematical hydraulic diameter, this formulation directly aligns with standard academic conventions for fundamental 2D channel flow verification. To rigorously verify the mathematical accuracy of the formulation, the FEniCS results were compared directly against the commercial Ansys Fluent solver operating under identical periodic boundary conditions.

\subsection{Mesh Resolution Requirements}
Because the SA formulation is a "low-Reynolds" model that does not utilize empirical wall functions, it requires mathematical integration directly through the viscous sublayer. Consequently, standard uniform computational grids are physically insufficient. To satisfy this strict boundary layer requirement, a custom structured inflation mesh was generated for the channel flow. Because the domain is a simple periodic rectangle, this highly anisotropic grid was constructed programmatically utilizing NumPy arrays and converted into a FEniCS-compatible format via the open-source \texttt{meshio} library. Based on the target bulk Reynolds number, this approach explicitly enforced a physical first-cell height of approximately $1.7 \times 10^{-3}$ meters, guaranteeing a non-dimensional wall distance of $y^+ \approx 1$ across the entire domain. 

Furthermore, to guarantee that the numerical results were not influenced by spatial discretization, a formal grid independence study was conducted. The full results of this spatial verification, proving convergence, are documented in Appendix~\ref{subsec:app_channel_mesh}.

\subsection{FEniCS vs. Ansys Comparison}
The channel flow was simulated in FEniCS using the decoupled Picard iteration method and compared against the Ansys Fluent baseline. To validate the solver accuracy both qualitatively and quantitatively, velocity contours were compared, and flow data was extracted along a wall-normal line probe spanning the width of the fully developed channel.

\begin{figure}[htpb]
	\centering
	\begin{subfigure}[b]{0.48\textwidth}
		\centering
		\includegraphics[width=\textwidth]{example-image} % Replace with Ansys channel contour
		\caption{Ansys Fluent SA Contour}
		\label{fig:channel_ansys}
	\end{subfigure}
	\hfill
	\begin{subfigure}[b]{0.48\textwidth}
		\centering
		\includegraphics[width=\textwidth]{example-image} % Replace with FEniCS channel contour
		\caption{FEniCS SA Contour}
		\label{fig:channel_fenics}
	\end{subfigure}
	
	\vspace{0.5cm}
	
	\begin{subfigure}[b]{0.8\textwidth}
		\centering
		\includegraphics[width=\textwidth]{example-image} % Replace with channel velocity profile plot
		\caption{Velocity Profiles}
		\label{fig:channel_profiles}
	\end{subfigure}
	\caption{Verification of the 2D Channel Flow ($Re_H \approx 20,300$). Comparison of the qualitative velocity contours between the Ansys Fluent baseline (a) and the custom FEniCS implementation (b), alongside the quantitative velocity profiles extracted across the channel (c).}
	\label{fig:channel_benchmark}
\end{figure}

As shown in Figure~\ref{fig:channel_benchmark}, the qualitative velocity contours of the custom FEniCS implementation visually replicate the Ansys Fluent baseline. More importantly, the quantitative line plots confirm that the velocity profiles generated by the custom FEniCS model match the Ansys reference data impeccably. The custom solver successfully resolves the logarithmic region of the turbulent boundary layer and accurately predicts the turbulent diffusion toward the channel center. This perfect alignment provides strict mathematical proof that the core Navier-Stokes solver, the relaxed Eikonal field, and the highly non-linear SA production/destruction terms are functioning correctly before introducing the geometric complexities of the U-bend.

\section{Verification Benchmark 2: U-Bend Geometry}
Following the successful channel flow verification, the solver was tested on the highly convective, separating secondary flows of the U-bend geometry evaluated in Chapter~\ref{ch:turbulence_modeling}. It is important to distinguish the physical scope of this specific verification test from the standard Ansys Verification Manual (VMFL048) case. 

While the original VMFL048 case is a full 3D simulation, the verification performed here utilized a strictly modified computational domain to ensure efficiency. First, the domain was reduced to a 2D geometry corresponding directly to the symmetry mid-plane of the 3D case. Second, to reduce unnecessary computational overhead during eventual topology optimization, the highly extended straight inlet and outlet sections ($1.555$ m) of the standard benchmark were explicitly truncated to strictly $30D$ ($0.840$ m).

To drive the flow, a uniform prescribed velocity ($\mathbf{u}_{in} = 1.42$ m/s) was imposed at the inlet alongside a fully turbulent boundary value ($\tilde{\nu}_{in} = 9.34 \times 10^{-5}$ m$^2$/s), while a zero relative static pressure ($p=0$ Pa) was assigned to the outlet. While the FEniCS domain is mathematically a 2D planar approximation of the symmetry mid-plane, the 3D physical pipe diameter ($D = 0.028$ m) and geometric radius of curvature ($R_c = 0.125$ m) were strictly retained as the characteristic length scales. Conserving these parameters ensures that the computed Reynolds number ($Re_D \approx 44,700$) and Dean number ($De \approx 15,000$) identically match the 3D Ansys reference case, maintaining non-dimensional consistency and correctly governing the internal secondary flows.

\subsection{Mesh Resolution Requirements}
Because this formulation utilizes a standard Galerkin weak form without spatial upwind stabilization, numerical stability in highly convective flows is strictly dependent on adequate grid resolution. As established in the channel flow benchmark, the SA model requires integration entirely through the viscous sublayer. 

Unlike the simple rectangular channel where a programmatic NumPy script was sufficient, the U-bend geometry features curved inner and outer arcs, complex corner transitions, and anisotropic boundary-layer inflation around non-orthogonal walls. To manage this severe geometric complexity while satisfying the $y^+ \approx 1$ resolution requirement, a custom mesh was generated using the advanced open-source mesh generator \textbf{Gmsh} \cite{Geuzaine2009}. The first element height adjacent to the wall was strictly controlled to $1.20 \times 10^{-5}$ meters across the entire domain. As documented comprehensively in Appendices~\ref{subsec:app_ubend_truncation} and \ref{subsec:app_ubend_mesh}, strict independence studies were conducted to explicitly prove that neither the $30D$ domain truncation nor the spatial grid discretization artificially influenced the secondary flow physics.

\subsection{FEniCS vs. Ansys Comparison}
Having mathematically validated the spatial resolution and boundary placement, an independent Ansys Fluent simulation was executed on the exact shortened 2D domain. This served as a strict 1:1 baseline to verify the accuracy of the custom FEniCS implementation. 

\begin{figure}[htpb]
	\centering
	\begin{subfigure}[b]{0.48\textwidth}
		\centering
		\includegraphics[width=\textwidth]{example-image} % Replace with Ansys U-bend contour
		\caption{Ansys Fluent SA Contour}
		\label{fig:ubend_ansys}
	\end{subfigure}
	\hfill
	\begin{subfigure}[b]{0.48\textwidth}
		\centering
		\includegraphics[width=\textwidth]{example-image} % Replace with FEniCS U-bend contour
		\caption{FEniCS SA Contour}
		\label{fig:ubend_fenics}
	\end{subfigure}
	
	\vspace{0.5cm}
	
	\begin{subfigure}[b]{0.8\textwidth}
		\centering
		\includegraphics[width=\textwidth]{example-image} % Replace with U-bend velocity profile plot
		\caption{Velocity Profiles}
		\label{fig:ubend_profiles}
	\end{subfigure}
	\caption{Verification of the 2D U-Bend Flow ($Re_D \approx 44,700$, $De \approx 15,000$). Comparison of the qualitative velocity contours between the Ansys Fluent baseline (a) and the custom FEniCS implementation (b), showing successful capture of the separation bubble, alongside the quantitative velocity profiles extracted at the bend exit (c).}
	\label{fig:ubend_verification}
\end{figure}

As shown in Figure~\ref{fig:ubend_verification}, the global flow fields and secondary flow separation bubbles observed in the qualitative velocity contours demonstrate excellent visual agreement. However, Research Question 2 strictly demands \textit{quantitative} validation. To satisfy this requirement, velocity data was extracted along a cross-sectional line probe positioned directly at the bend exit, marking the exact transition from the end of the curvature to the start of the straight outlet section. 

The quantitative velocity profile computed by the custom FEniCS implementation traces the commercial Ansys results closely. The FEniCS implementation accurately replicates the boundary layer growth, the momentum shift toward the outer wall, and the turbulent diffusion predicted by the commercial solver, thereby confirming the mathematical integrity of the code.

\section{Conclusion}
This chapter has detailed the mathematical formulation, stabilization strategy, and rigorous verification of the steady-state Spalart-Allmaras model within the FEniCS framework. By successfully replicating the benchmark results of Ansys Fluent, this implementation directly answers the second research question of this thesis: achieving numerical stability and quantitative agreement in a custom finite element solver, laying a fully consistent foundation for the topology optimization routines developed in the subsequent chapters. 

First, numerical stability was successfully achieved without the need for complex mathematical stabilization artifacts. By decoupling the Navier-Stokes and Spalart-Allmaras equations via a rigorous Picard iteration scheme, and utilizing an Incremental Pressure Correction Scheme (IPCS) to manage the convective fluid terms, the highly non-linear system was robustly driven to equilibrium. Supported by high-density mesh resolution and Yoon's relaxed Eikonal equation, this standard Galerkin formulation completely avoided the numerical oscillations that typically plague turbulent finite-element solvers. 

Second, quantitative agreement was thoroughly verified. The solver produced velocity and eddy viscosity profiles that perfectly matched the commercial Ansys Fluent baseline in both simple periodic channel flows ($Re_H \approx 20,300$) and the highly convective, geometrically complex U-bend ($Re_D \approx 44,700$, $De \approx 15,000$). Consequently, this decoupled Picard FEniCS solver is proven to be a mathematically sound and reliable fluid dynamics engine, specifically structured to support the continuous adjoint topology optimization campaigns presented in Chapters~\ref{ch:topology_optimization_methodology} and \ref{ch:benchmark_results}.
\chapter{Topology Optimization for Turbulent Flows: Methodology}
\label{ch:topology_optimization_methodology}

\section{Introduction}
Building upon the verified Spalart-Allmaras (SA) fluid dynamics solver \cite{Spalart1992, Spalart1994} developed in Chapter~\ref{ch:fenics_implementation}, this chapter introduces the mathematical framework required to perform Topology Optimization (TO) in the turbulent regime. The core objective of this methodology is to algorithmically determine the optimal distribution of fluid channels within a given design domain to reduce a prescribed hydraulic loss metric, such as power dissipation or pressure loss \cite{Alexandersen2020}. 

This chapter details the density-based fictitious-domain formulation used to represent solid and fluid regions on a fixed computational mesh \cite{Borrvall2003, Gersborg2005}. It then explains the specific turbulent-flow ingredients retained for the remainder of the thesis: a Brinkman-type momentum resistance, a power-law topology penalty for the SA and wall-distance equations, and the reciprocal penalized wall-distance formulation used by the Alexandersen pipe-bend computations.

To solve this optimization problem in a computationally tractable manner, continuous adjoint sensitivity analysis is required \cite{Papoutsis2016, Othmer2014}. The optimization results reported in this thesis use the frozen-turbulence, or Constant Eddy Viscosity (CEV), assumption \cite{Wu2023, Schramm2018}. This assumption is stated explicitly because it is both the practical reason the turbulent optimization loop remains stable and the main limitation of the sensitivity model.

Following the derivation of these sensitivities, the continuous optimization problem is formally defined, establishing the admissible objective functions and the strict fluid volume constraints over explicitly defined design regions. Finally, the mathematical techniques for active-domain density filtering \cite{Lazarov2011} and Heaviside projection \cite{Wang2011} are detailed alongside their FEniCS unified form language (UFL) implementation \cite{Alnaes2015, Logg2012} to ensure the generation of manufacturable designs. The algorithms, penalty formulations, and execution scripts are housed within the \texttt{FluidTO} directory of the public repository, with the frozen-turbulence workflow implemented in \texttt{TurbulentTO\_Frozen.py}.

\section{The Fictitious Domain Method and Brinkman Penalization}
\label{sec:fictitious_domain}
In standard shape optimization, the solid-fluid boundary is explicitly defined by the mesh, and the nodes are physically moved to improve the design. In density-based topology optimization, however, the computational domain ($\Omega$) remains strictly fixed. To create solid walls and fluid channels within this fixed grid, the Fictitious Domain Method is employed. 

Originally applied to Stokes flow by Borrvall and Petersson \cite{Borrvall2003} and extended to the Navier-Stokes equations by Gersborg-Hansen et al. \cite{Gersborg2005}, this method assigns a continuous design variable to the spatial coordinates in the domain. In this specific implementation convention, a physical topology value of \textbf{$\bar{\gamma} = 1$} represents a pure \textbf{fluid} channel, while \textbf{$\bar{\gamma} = 0$} represents a \textbf{solid}, impermeable structure. To force the fluid velocity to zero inside the solid regions without mathematically removing the finite elements, the fluid is treated as a fictitious porous medium. 

This flow resistance is quantified by an inverse permeability term, or Brinkman penalty function, denoted as $\alpha(\bar{\gamma})$. To ensure the optimization algorithm (which relies on continuous gradients) can transition smoothly between fluid and solid states, $\alpha$ must be a continuous interpolation function. This solver utilizes a convex Rational Approximation of Material Properties (RAMP) interpolation \cite{Stolpe2001} scaled between a maximum solid penalty ($\alpha_{solid}$) and a minimum fluid penalty ($\alpha_{fluid}$):
\begin{equation}
	\alpha(\bar{\gamma}) = \alpha_{solid} + (\alpha_{fluid} - \alpha_{solid}) \frac{\bar{\gamma} (1 + q)}{\bar{\gamma} + q}
	\label{eq:brinkman_interpolation}
\end{equation}
where $\alpha_{solid}$ is a sufficiently large penalty factor that completely stagnates the flow when $\bar{\gamma} = 0$, and $q$ is a tuning parameter that controls the convexity of the penalization curve. Unlike traditional structural optimization which relies on the SIMP power-law \cite{Bendsoe1988}, fluid optimization generally avoids SIMP for momentum because its derivative vanishes at $\bar{\gamma} = 0$. By utilizing the RAMP formulation, a strictly non-zero gradient is guaranteed across the entire design domain.

\subsection{The Density Chain and Domain Bounding}
While conceptually represented by the physical topology field in the governing equations, the computational implementation dictates that the layout is managed through a distinct chain of density representations to manage design and non-design domains. To explicitly avoid mathematical confusion with the fluid mass density ($\rho$) utilized in the Reynolds-Averaged Navier-Stokes equations, this thesis exclusively utilizes the $\gamma$ notation for all topological variables. However, these variables correspond directly to the exact implementation code variables as follows:
\begin{itemize}
	\item \textbf{Raw design density} ($\gamma \equiv \rho_{\mathrm{raw}}$): The baseline parameter manipulated by the optimizer.
	\item \textbf{Filtered density} ($\tilde{\gamma} \equiv \rho_f$): The spatially smoothed design field.
	\item \textbf{Projected effective density} ($\bar{\gamma} \equiv \rho_{\mathrm{eff}}$): The bounded, physical topology field evaluated by the governing equations.
\end{itemize}

The raw design density, $\gamma$, is defined as a Discontinuous Galerkin (DG0) \cite{Hesthaven2007} function over the entire mesh. The optimization algorithm strictly manipulates the degrees of freedom corresponding to the actively optimized design space (\texttt{ActiveDV}), while passive (non-design) cells are subsequently restored to their prescribed bounding values. 

This raw variable is subjected to a spatial PDE filter \cite{Lazarov2011} to yield a smoothed field $\tilde{\gamma}$. To recover sharp boundaries and strictly preserve required non-design regions, the filtered density is then projected and mapped into the bounded effective density ($\bar{\gamma}$) \cite{Guest2004}:
\begin{equation}
	\bar{\gamma} = \gamma_{\mathrm{lower}} + (\gamma_{\mathrm{upper}} - \gamma_{\mathrm{lower}}) \Pi_\beta(\tilde{\gamma})
\end{equation}
where $\Pi_\beta$ is the projection operator, and $\gamma_{\mathrm{lower}}$ and $\gamma_{\mathrm{upper}}$ define the absolute permissible bounds for any given cell. This rigorous bounding guarantees that predefined features, such as non-design fluid inlet extensions ($\gamma_{\mathrm{lower}} = \gamma_{\mathrm{upper}} = 1$) or explicitly constrained structural walls ($\gamma_{\mathrm{lower}} = \gamma_{\mathrm{upper}} = 0$), remain rigidly fixed without interfering with the mathematical behavior of the active optimization space.

The computational mesh is therefore partitioned into an active design domain and passive non-design domains. The active design domain contains the cells whose raw densities are updated by MMA. Fluid non-design domains are fixed as permanent fluid, typically to preserve inlet and outlet extensions, buffer channels, or measurement regions. Solid non-design domains are fixed as permanent solid, typically to impose external walls, blocked regions, or geometric features that must not be modified. These passive cells still participate in the forward flow, turbulence, and wall-distance solves, but they are excluded from the optimization masks: the volume constraint is evaluated only over the configured volume mask $\Omega_V$, and domain-integral objectives are evaluated only over the configured objective mask $\Omega_D$. Consequently, prescribed inlet/outlet buffers and mandatory solid walls do not count toward the optimized fluid fraction and do not directly bias the domain-integrated objective.

\section{Penalization of the Governing Equations}
\label{sec:penalization_governing}
To effectively optimize turbulent manifolds, the respective penalties must be rigorously integrated into the turbulence model and the wall-distance calculations to ensure the physics remain consistent as the solid boundaries evolve. All physical penalties strictly utilize the projected effective density, $\bar{\gamma}$.

\subsection{Modified Reynolds-Averaged Navier-Stokes Equations}
The primary application of the Brinkman penalty is the modification of the steady-state, incompressible RANS momentum equation. The RAMP-interpolated inverse permeability $\alpha(\bar{\gamma})$ is added as a linear sink term (Darcy friction) to the momentum conservation equation:
\begin{equation}
	\rho (\mathbf{u} \cdot \nabla) \mathbf{u} = -\nabla p + \nabla \cdot \left[ (\mu + \mu_t) \left( \nabla \mathbf{u} + (\nabla \mathbf{u})^T \right) \right] - \alpha(\bar{\gamma}) \mathbf{u}
	\label{eq:penalized_momentum}
\end{equation}
It is critical to note the retention of the full symmetric rate-of-strain tensor. Here, $\rho$ strictly denotes the fluid mass density. Because the turbulent eddy viscosity ($\mu_t$) varies drastically across the spatial domain \cite{Pope2000}, the effective fluid viscosity is no longer constant, and the transposed term exerts a non-zero physical influence on the momentum balance. In regions where the optimizer places solid material ($\bar{\gamma} \to 0$), the enormous magnitude of $\alpha_{solid} \mathbf{u}$ dominates all inertial and viscous forces, forcing the local velocity mathematically to zero. 

\subsection{Modified Spalart-Allmaras Transport Equation}
As the Fictitious Domain Method forces velocity to zero in the solid regions, the turbulence generation inside those regions should theoretically cease. However, due to numerical diffusion, turbulent eddy viscosity ($\tilde{\nu}$) can inadvertently bleed into the solid regions. This unphysical behavior destabilizes the solver and heavily distorts the adjoint sensitivity analysis.

To strictly enforce a laminar, zero-turbulence state inside the solid structures, a corresponding penalty sink term is added to the Spalart-Allmaras transport equation \cite{Yoon2016, Dilgen2018}:
\begin{equation}
	\mathbf{u} \cdot \nabla \tilde{\nu} = c_{b1} \tilde{S} \tilde{\nu} - c_{w1} f_w \left( \frac{\tilde{\nu}}{y} \right)^2 + \frac{1}{\sigma} \left[ \nabla \cdot \left( (\nu + \tilde{\nu}) \nabla \tilde{\nu} \right) + c_{b2} (\nabla \tilde{\nu})^2 \right] - R_{\tilde{\nu}}(\bar{\gamma}) \tilde{\nu}
	\label{eq:penalized_sa}
\end{equation}
Here $R_{\tilde{\nu}}(\bar{\gamma})$ is an artificial topology-dependent reaction coefficient. It is not part of the original SA model; it is introduced only to suppress spurious values of $\tilde{\nu}$ in the fictitious solid and in intermediate porous material.

\subsection{Power-Law Topology Penalty}
The SA and wall-distance penalties require an interpolation from the physical density field $\bar{\gamma}$ to a solid indicator. For the computations emphasized in this thesis, this interpolation is the positive-part power-law indicator
\begin{equation}
	S_{\mathrm{pow}}(\bar{\gamma};N) =
	[\max(1-\bar{\gamma},0)]^N ,
	\label{eq:power_penalty_indicator}
\end{equation}
so that a generic reaction penalty is written as $R_{\psi}(\bar{\gamma})=\alpha_{\psi}S_{\mathrm{pow}}(\bar{\gamma};N_{\psi})$, with $\psi \in \{\tilde{\nu},G\}$ depending on the equation being penalized. The indicator is zero in pure fluid and increases smoothly toward one as the density approaches solid. For thresholded wall-distance updates, the same construction is applied to a shifted solid indicator,
\begin{equation}
	S_{\mathrm{pow},\theta}(\bar{\gamma};N,\theta_s) =
	\left[\max\left(\frac{\theta_s-\bar{\gamma}}{\theta_s},0\right)\right]^N .
	\label{eq:threshold_power_penalty_indicator}
\end{equation}

\subsection{Reciprocal Penalized Wall-Distance Equation}
The Spalart-Allmaras model is highly sensitive to the wall distance, $y$. While differential equations for wall distance computation are well-established \cite{Tucker2003}, topology optimization introduces a unique challenge. The location of the wall is not a static mesh boundary, but rather the dynamically evolving $\bar{\gamma} = 0$ contour. Therefore, the relaxed Eikonal equation must be penalized to dynamically recognize the fictitious solid as a physical boundary \cite{Yoon2016}:
\begin{equation}
	\nabla G \cdot \nabla G + \sigma_w G (\nabla \cdot \nabla G) = (1 + 2\sigma_w) G^4 + R_G(\bar{\gamma}) (G - G_0) 
	\label{eq:penalized_eikonal}
\end{equation}
Consistent with the SA transport modifications, the reciprocal wall-distance equation uses a staged multiplier paired with the power-law solid indicator. In the Alexandersen pipe-bend configuration, $R_G(\bar{\gamma})=\alpha_G S_{\mathrm{pow},\theta}(\bar{\gamma};N_G,\theta_s)$ with a design-driven wall-density source. This is the implementation mode denoted \texttt{reciprocal\_penalized}. The algebraic recovery function
\begin{equation}
	y = \frac{1}{G} - \frac{1}{G_0}
	\label{eq:reciprocal_wall_distance_recovery}
\end{equation}
then collapses the SA wall distance toward zero inside topology-created solid regions while retaining the physical wall distance in fluid regions.

\section{Boundary Conditions and Turbulence Initialization}
\label{sec:boundary_conditions}
To stabilize the topology optimization and ensure physical accuracy, boundary constraints are strictly managed across both fluid and turbulence regimes. Non-design fluid buffer regions are strictly mandated at inlets and outlets to isolate the active optimization domain from artificial numerical shear. The framework robustly handles boundary interactions by incorporating pressure outlets coupled with optional velocity-component constraints, as well as pressure pin controls to fully constrain the discrete PDE problem.

At the domain inlets, establishing accurate turbulent boundary conditions natively bridges physical intent with numerical stability. The prescribed turbulence intensity and the associated turbulence length scale are systematically converted into an intermediate eddy-viscosity ratio. This ratio directly calculates the requisite boundary value for the Spalart-Allmaras working variable $\tilde{\nu}$. Throughout the computation, to protect the solver from chaotic geometric fluctuations during intermediate design cycles, rigorous numerical floors, ceilings, and internal clipping functions are imposed to bind the working variable safely within physical limits.

\section{Sensitivity Analysis and the Frozen Turbulence Assumption}
\label{sec:frozen_turbulence}
With the forward state equations fully defined and penalized, the topology optimization algorithm requires the gradient (sensitivity) of the objective function with respect to the design variable. Because evaluating finite differences for every mesh element is computationally impossible, the continuous adjoint method is employed \cite{Othmer2014}. 

However, calculating the fully coupled, exact adjoint for turbulent flows is mathematically formidable. As demonstrated by Zymaris et al. \cite{Zymaris2009}, formulating the exact continuous adjoint for the Spalart-Allmaras model requires differentiating the highly non-linear production and destruction terms, accounting for the variations in the distance function, and managing the intense cross-coupling with the momentum equations. The complexity of these exact adjoints frequently introduces severe numerical stiffness and divergence.

To make the gradient calculation both numerically robust and computationally feasible, the primary optimization workflows in this implementation rely on the \textbf{Frozen Turbulence Assumption}, frequently referred to in the literature as the Constant Eddy Viscosity (CEV) assumption \cite{Wu2023, Schramm2018}.

\subsection{Mechanism of Frozen Turbulence}
The frozen turbulence assumption simplifies the sensitivity analysis by decoupling the turbulence transport from the design update. During the adjoint formulation, the turbulent eddy viscosity field ($\nu_t$) generated by the forward SA solver is assumed to be locally independent of infinitesimally small changes in the physical topology ($\bar{\gamma}$). Mathematically, this implies that during the calculation of the adjoint sensitivities:
\begin{equation}
	\frac{\partial \nu_t}{\partial \bar{\gamma}} \approx 0
\end{equation}
By treating the eddy viscosity field as frozen (a spatially varying but constant parameter), the complex adjoints of the Spalart-Allmaras and Eikonal equations do not need to be solved. The optimizer only solves the standard fluid adjoint equations (utilizing the effective viscosity $\mu + \mu_t$) to determine how a change in topology affects the viscous dissipation. 

\subsection{Justification and Limitations}
The primary advantage of the frozen turbulence approach is not that it is universally better, but that it makes the optimization problem solvable within a reasonable numerical and implementation budget. Relative to a fully coupled turbulent adjoint, it offers greater numerical stability, a smaller implementation burden, and a materially lower computational cost per design iteration. This is why frozen turbulence remains common in industrial adjoint workflows and why it is used as the baseline in this thesis \cite{Othmer2014, Schramm2018}.

The price of this robustness is that the gradients are approximate reduced derivatives. Because the algorithm ignores how an infinitesimal topology change would alter future turbulence production, the sensitivity can be inaccurate when the objective is strongly controlled by separated shear layers, stall-like behavior, or other regimes where the eddy viscosity changes sharply with geometry. Wu et al.~\cite{Wu2023} show that the CEV assumption can still achieve useful optimization goals in benign flows, but its sensitivity accuracy and adjoint convergence can deteriorate near stall and in unsteady settings. Zymaris et al.~\cite{Zymaris2009} further demonstrate, for SA-based continuous adjoints, that including the turbulence-equation adjoint improves derivative accuracy for loss functionals. The present thesis therefore defends the frozen assumption as a practical and verifiable approximation, not as an exact turbulence sensitivity model. Attempts to include the SA turbulence equation in the adjoint were explored during the implementation phase, but the resulting solves were not sufficiently stable for the thesis optimization campaign.

\section{Optimization Problem Formulation}
\label{sec:optimization_problem}
Using the frozen-turbulence sensitivities described above, the topology optimization task is formulated as a mathematical programming problem, where the optimizer must find a spatial distribution that minimizes a specific performance metric subject to physical and geometric constraints.

\subsection{Objective Functions}
For fluid manifolds, the primary engineering goal is to guide the fluid from the inlets to the outlets with the minimum possible hydraulic loss. In incompressible flow, this loss can be expressed either through a domain-integrated power dissipation objective or through a pressure-based objective, depending on the benchmark configuration and the quantity reported in the reference study.

For dissipation-driven cases, it is crucial to explicitly distinguish the algorithmic optimization objective ($J_D$) from the true physical validation metrics analyzed later in Chapter~\ref{ch:benchmark_results}. The dissipation objective includes artificial Brinkman drag and is integrated exclusively over the configured objective mask ($\Omega_D$):
\begin{equation}
	J_D(\mathbf{u}, \bar{\gamma}) = \int_{\Omega_D} \left[ \frac{1}{2} (\mu + \mu_t) \left( \nabla \mathbf{u} + (\nabla \mathbf{u})^T \right) : \left( \nabla \mathbf{u} + (\nabla \mathbf{u})^T \right) + \alpha(\bar{\gamma}) \mathbf{u} \cdot \mathbf{u} \right] d\Omega
	\label{eq:objective_function}
\end{equation}
The first term in the integral represents the physical dissipation caused by the effective fluid viscosity. The second term represents the mathematical energy lost when the fluid attempts to penetrate the fictitious solid regions. 

For pressure-driven benchmarks with prescribed inflow and fixed outlet pressure, the optimization objective can instead be defined as the average inlet pressure,
\begin{equation}
	J_p(p) = \frac{1}{|\Gamma_{\mathrm{in}}|} \int_{\Gamma_{\mathrm{in}}} p \, d\Gamma .
	\label{eq:inlet_pressure_objective}
\end{equation}
When the outlet pressure is used as the gauge reference, minimizing $J_p$ is equivalent to minimizing the pressure loss required to drive the imposed flow rate through the topology. Unlike Eq.~\ref{eq:objective_function}, this is a boundary integral rather than a volume integral over $\Omega_D$. Its sensitivity is nevertheless mapped back only to the active design degrees of freedom through the same active-mask filter and transpose-filter machinery, so the prescribed non-design inlet region is used only to measure the pressure and is not itself optimized.

\subsection{Volume Constraint}
Because the goal is to design a targeted manifold or pipe system, a strict volume constraint must be imposed. The volume fraction constraint is enforced via a discrete volume region mask ($\Omega_V$), strictly dictating the proportion of material evaluated by the optimizer:
\begin{equation}
	g(\bar{\gamma}) = \frac{\int_{\Omega_V} \bar{\gamma} \, d\Omega}{V_{\mathrm{mask}}} - V_{\max} \leq 0
	\label{eq:volume_constraint}
\end{equation}
where $V_{\mathrm{mask}}$ is the total volume of the designated integration space, and $V_{\max}$ is the maximum allowable fluid volume fraction.

\subsection{The Complete Continuous Optimization Problem}
Combining the selected objective function, the governing equations, and the constraints, the continuous topology optimization problem for turbulent flow is formally stated as:
\begin{equation}
	\begin{aligned}
		& \underset{\gamma(\mathbf{x}) \in [0, 1]}{\text{minimize}}
		& & J(\mathbf{u}, p, \bar{\gamma}) \\
		& \text{subject to}
		& & \text{Modified RANS Equations (Eq.~\ref{eq:penalized_momentum})} \\
		& & & \text{Modified Spalart-Allmaras Equation (Eq.~\ref{eq:penalized_sa})} \\
		& & & \text{Configured Wall-Distance Equation} \\
		& & & g(\bar{\gamma}) \leq 0
	\end{aligned}
\end{equation}
where $J$ denotes the selected benchmark objective, either the dissipation objective $J_D$ in Eq.~\ref{eq:objective_function} or the inlet-pressure objective $J_p$ in Eq.~\ref{eq:inlet_pressure_objective}. The wall-distance constraint used for the Alexandersen-focused computations is the reciprocal penalized relaxed Eikonal equation in Eq.~\ref{eq:penalized_eikonal}, with the wall distance recovered by Eq.~\ref{eq:reciprocal_wall_distance_recovery}.

To solve this constrained minimization problem, the FEniCS implementation utilizes the Method of Moving Asymptotes (MMA) \cite{Svanberg1987}, an optimization algorithm specifically designed for handling large-scale, gradient-based topology optimization problems. MMA approximates the highly non-linear design space using convex asymptotes, relying on strict move limits which cap the maximum allowable change in the active raw design variable ($\gamma$) during any single optimization step.

\section{Density Filtering and Projection}
\label{sec:filtering}
A well-known mathematical complication in density-based topology optimization is the lack of existence of a classical solution. If the design space is left unconstrained, the optimizer will tend to create infinitely thin micro-structures (checkerboarding). To guarantee a mesh-independent and manufacturable design, a minimum length scale must be enforced.

\subsection{The Discontinuous Galerkin Helmholtz Filter}
To prevent checkerboarding, a density filter is applied to the raw design variable generated by the optimizer. Because a continuous nodal definition causes the Brinkman penalty to interpolate smoothly across the interior of boundary elements, creating a semi-permeable gray gradient that allows mass leakage, this implementation discretizes the density field using piecewise constant elements (Discontinuous Galerkin, DG0) \cite{Arnold2002}. By assigning a single, uniform density value to the entire element, the structural boundaries remain perfectly sharp at the element edges.

Consequently, the standard continuous Helmholtz filter cannot be applied. The filter is mathematically adapted to account for the discontinuous boundaries utilizing an interior penalty method across the element facets ($\Gamma_{int}$). Standard volume diffusion is mathematically replaced by interior facet jump penalties:
\begin{equation}
	\int_{\Gamma_{int}} r^2 \frac{\alpha_{dg}}{h_{avg}} \llbracket \tilde{\gamma} \rrbracket \cdot \llbracket v \rrbracket \, dS + \int_{\Omega} \tilde{\gamma} v \, dx = \int_{\Omega} \gamma v \, dx
	\label{eq:dg_helmholtz_filter}
\end{equation}
where $v$ is the test function, $\llbracket \cdot \rrbracket$ denotes the jump operator across a facet, $h_{avg}$ is the average cell size, and $r = \frac{R}{2\sqrt{3}}$ is the scaled filter radius.

\subsubsection{Active-Mask Normalization}
To appropriately handle the severe boundary interactions near designated inlet and outlet buffers, this filtering is strictly active-mask-normalized. Rather than applying the filter blindly across the entire domain, the filtered output evaluates as:
\begin{equation}
	\tilde{\gamma} = \frac{H(M \gamma)}{H(M)}
\end{equation}
where $H$ represents the Helmholtz PDE operator and $M$ represents the isolated active design indicator mask. Crucially, during sensitivity analysis, the mathematical transpose of this normalized filter must be explicitly derived and applied to map the calculated objective gradients back into the active DG0 design space. The explicit programmatic integration of this interior penalty filter, active-mask normalization, and the corresponding adjoint transpose logic is demonstrated in Listing~\ref{lst:pde_filter}.

\begin{lstlisting}[language=Python, caption={Schematic pseudocode of the DG0 Helmholtz filter with active-mask normalization and interior facet jump penalties. Extracted from \texttt{TurbulentTO\_Frozen.py}.}, label={lst:pde_filter}]
	# r is the internal PDE coefficient scaled by 1/(2*sqrt(3))
	r = r_filter / (2.0 * 3.0**0.5)  
	
	u_filter = TrialFunction(DensitySpace)
	v_filter = TestFunction(DensitySpace)
	filter_in = Function(DensitySpace)
	
	h_avg = (CellDiameter(mesh)("+") + CellDiameter(mesh)("-")) / 2.0
	
	def pde_filter_raw(input_field, output_field):
	alpha_dg = 4.0
	helmholtz = (
	r**2 * (alpha_dg / h_avg * dot(jump(v_filter, n), jump(u_filter, n))) * dS
	+ u_filter * v_filter * dx
	- filter_in * v_filter * dx
	)
	assign(filter_in, input_field)
	solve(lhs(helmholtz) == rhs(helmholtz), output_field)
	return output_field
	
	def pde_filter_design_density(input_field, output_field):
	"""Filter active design variables and normalize by the filtered mask."""
	# Active indicator M is filtered to create the denominator H(M)
	filter_work.vector()[ActiveDV] = input_field.vector()[ActiveDV]
	pde_filter_raw(filter_work, output_field)
	
	# rho_f = H(M rho) / H(M); passive cells restored after filtering
	output_field.vector()[:] /= filter_denominator_values
	output_field.vector()[PassiveDV] = input_field.vector()[PassiveDV]
	return output_field
	
	def pde_filter_design_gradient(input_field, output_field):
	"""Apply the transpose of the active mask-normalized density filter."""
	# Adjoint of rho -> H(M rho) / H(M), restricted to active DOFs
	filter_work.vector()[ActiveDV] = input_field.vector()[ActiveDV] / filter_denominator_values[ActiveDV]
	pde_filter_raw(filter_work, output_field)
	
	output_field.vector()[PassiveDV] = 0.0
	return output_field
\end{lstlisting}

\subsection{Heaviside Projection}
While the Helmholtz filter successfully imposes a minimum length scale, it inherently smooths the design field, resulting in a large gray transition zone between the solid and fluid regions. To recover a crisp, discrete topology, a continuous Heaviside projection is applied to the filtered variable \cite{Wang2011, Guest2004}:
\begin{equation}
	\Pi_\beta(\tilde{\gamma}) = \frac{\tanh(\beta \eta) + \tanh(\beta (\tilde{\gamma} - \eta))}{\tanh(\beta \eta) + \tanh(\beta (1 - \eta))}
	\label{eq:tanh_projection}
\end{equation}
where $\eta \in [0, 1]$ is the projection threshold (typically 0.5), and $\beta$ is a steepness parameter. As $\beta \to \infty$, the function approaches a strict Heaviside step function.

\section{Algorithmic Workflow and Solver Implementation}
\label{sec:algorithmic_workflow}
Executing the optimization requires a highly robust discrete, iterative algorithm. To support the severe non-linearities of the turbulent regime, the solver architecture natively manages the segregated Spalart-Allmaras turbulence model, explicit physics decoupling via Picard iterations, and deeply hybridized forward flow routines.

\subsection{Boundary Conditions and Turbulence Initialization}
At the domain inlets, establishing accurate turbulent boundary conditions natively bridges physical intent with numerical stability. The prescribed turbulence intensity and the associated turbulence length scale are systematically converted into an intermediate eddy-viscosity ratio. This ratio directly calculates the requisite boundary value for the Spalart-Allmaras working variable $\tilde{\nu}$. Throughout the computation, to protect the solver from chaotic geometric fluctuations during intermediate design cycles, rigorous numerical floors, ceilings, and internal clipping functions are imposed to bind the working variable safely within physical limits.

Furthermore, the solver dynamically handles pressure outlets coupled with optional velocity-component constraints, allowing the flow field to organically adapt to the evolving structural viscous dissipation without artificial boundary blockages.

\subsection{SA Numerical Solution Controls}
Although Streamline Upwind Petrov-Galerkin stabilization is a common treatment for advection-dominated Galerkin formulations \cite{Brooks1982}, the computations reported in this thesis do not rely on SUPG terms or pseudo-time substepping in the SA equation. Instead, numerical control is provided through Picard-level coupling, carefully bounded turbulence variables, continuation of the topology penalties, and the hybrid forward-flow recovery strategy described below.

\subsection{Hybrid Forward Flow Recovery and Picard Decoupling}
Rather than relying solely on simple fixed-point decoupled sweeps, the forward fluid solver utilizes an advanced hybrid architecture integrating both Incremental Pressure Correction Scheme (IPCS) methods and monolithic SNES workflows. 

During the optimization loop, a Picard iteration scheme strictly decouples the fluid and turbulence physics. The flow solver advances the momentum field using the eddy viscosity from the previous iteration, and this updated velocity is then passed back to the SA solver. During volatile intermediate design steps, the fluid framework initiates convection continuation strategies and leverages Stokes-Brinkman warm starts to slowly reintroduce inertial physics. If convergence stalls, the SNES solver autonomously steps through multiple recovery schedules (including dynamic tolerance reductions and altered line-search methods). 

The optimization architecture supports non-converged acceptance, allowing the algorithm to passively absorb intermediate errors and maintain topological momentum. However, prior to the continuous adjoint assembly, optional configured flow solves are executed, becoming strictly bound where enabled. When utilized, these strict solves confirm that the sensitivities, and specifically the frozen turbulent fields, are mathematically derived from a fully equilibrated state. Otherwise, relying on the accepted residual criteria serves as a necessary pragmatic compromise to maintain solver momentum in highly separated regimes.

\subsection{Python Implementation and Continuation Strategy}
To ensure convergence to a physically meaningful binary design, the algorithm iterates over a strict continuation schedule. As the optimization progresses, the Brinkman penalty parameter ($q$) and the projection steepness ($\beta$) are iteratively increased, while the MMA move limits are proportionally tightened.

The global optimization routine, reflecting the hybrid solver calls, explicit gradient filtering, frozen adjoint solve, and penalty schedules, is outlined in Listing~\ref{lst:to_loop} as a schematic pseudocode. Note that while the repository source code uses the variable name \texttt{rho} for the design density, it is denoted here as \texttt{rho\_raw} to explicitly avoid confusion with the fluid mass density $\rho$.

\begin{lstlisting}[language=Python, caption={Schematic pseudocode representing the algorithmic architecture of the iterative topology optimization loop using the frozen turbulence assumption. Code structure has been simplified to highlight the core physics, hybrid solvers, and adjoint pipelines.}, label={lst:to_loop}]
	for stage_idx, q_val in enumerate(Q_PENAL_SCHEDULE):
	# Update continuation parameters for the current stage
	BETA_PROJ.assign(BETA_PROJ_SCHEDULE[stage_idx])
	q_penal.assign(q_val)
	
	while inner_count < max_iters_now and not objective_converged:
	# 1. Spatial Filtering & Penalized Wall-Distance Update
	rho_f = pde_filter_design_density(rho_raw, rho_f)
	update_wall_distance_field()
	
	if iter_count == 0:
	initialize_forward_guess_with_stokes()
	
	# 2. Forward Physics (Outer Picard Loop for Decoupling)
	for picard_idx in range(PICARD_STEPS):
	solve_forward_picard(...) # Robust intermediate flow solve
	
	sa_model.construct_forms(w_fwd.sub(0))
	sa_model.solve_turbulence_model()
	sa_model.update_variables(relaxation=TURBULENCE_RELAXATION)
	
	# Lock in the updated eddy viscosity for the frozen assumption
	nu_tilde_frozen.assign(sa_model.nu_tilde0)
	
	# Final configured flow solve, strict where enabled
	solve_forward_snes(...) 
	
	# 3. Sensitivity Analysis (Continuous Frozen Adjoint)
	solve_adjoint(objective_adjoint_form)
	
	# 4. Gradient Assembly & Transpose Filtering
	unfiltered_gradient.vector()[:] = assemble(objective_ddx)[:]
	filtered_gradient = pde_filter_design_gradient(unfiltered_gradient, filtered_gradient)
	
	# 5. Method of Moving Asymptotes (MMA) Design Update
	(xmma, ...) = mmasub(mmma, num_mma, iter_count, xval, xmin, xmax, ...)
	
	# Update raw density array with strict bounding
	rho_raw.vector()[ActiveDV] = xmma[:, 0]
	
	inner_count += 1
\end{lstlisting}

\section{Conclusion}
This chapter has established the mathematical and algorithmic framework used for the turbulent topology optimization results in this thesis. By employing the Fictitious Domain Method with Brinkman momentum resistance \cite{Borrvall2003, Gersborg2005}, the evolving structural topology is integrated into the fixed computational grid. The SA working-variable and reciprocal wall-distance equations are coupled to the topology through the power-law penalty, while the DG0 Helmholtz interior penalty filter \cite{Lazarov2011} and continuous Heaviside projection \cite{Wang2011} enforce a minimum length scale and reduce mesh-dependent numerical artifacts.

The optimization is driven by the frozen-turbulence continuous adjoint \cite{Wu2023, Schramm2018}. This assumption is defensible as a stable engineering approximation, but the thesis explicitly treats it as an approximation whose limitations must be checked through finite-difference sensitivity verification and independent CFD validation. Embedded within a Picard-decoupled Python architecture utilizing hybrid IPCS/SNES recovery schemes, this methodology transforms the verified fluid dynamics engine from Chapter~\ref{ch:fenics_implementation} into the automated design tool evaluated in Chapter~\ref{ch:benchmark_results}.

\chapter{Topology Optimization for Turbulent Flows: Results and Cross-Code CFD Verification}
\label{ch:benchmark_results}

\section{Introduction}
With the mathematical framework for turbulent topology optimization established in Chapter~\ref{ch:topology_optimization_methodology}, this chapter presents the resulting optimized topologies and rigorously evaluates their fluid dynamic performance. The objective is not merely to demonstrate that turbulent optimization yields different geometries than laminar assumptions, but to verify the physical credibility of the Spalart-Allmaras (SA) driven designs through independent computational fluid dynamics (CFD) analysis. 

To achieve this, a three-stage evaluation pipeline is utilized. First, the turbulent topology optimization is executed within the custom FEniCS framework, analyzing the objective function convergence, fluid routing, and eddy viscosity distributions. Second, the optimized continuous design fields are extracted as sharp, discrete geometries and exported to Ansys Fluent. These geometries are re-analyzed using the identical SA model to quantify cross-code consistency and assess the impact of geometry extraction. Finally, the designs are evaluated in Ansys using the higher-fidelity SST $k-\omega$ (and standard $k-\epsilon$) models to determine the sensitivity of the performance metrics to the underlying turbulence closure, particularly in regimes characterized by high Dean numbers and flow separation.

\section{Boundary Conditions and Evaluation Metrics}
\label{sec:boundary_conditions_metrics}
To accurately capture the physics of internal turbulent manifolds, the boundary conditions for all benchmark cases strictly utilize uniform inlet velocity profiles combined with a prescribed turbulence intensity (TI), diverging from the fully developed parabolic profiles typical of laminar topology optimization. Furthermore, to prevent the optimizer from exploiting artificial shear layers generated directly at the boundary constraints, non-design fluid regions ($\gamma = 1$) are enforced at all inlets and outlets. This allows the turbulent flow to naturally develop before interacting with the actively optimized design domain.

The primary performance metric is the total power dissipation (or equivalently, the total pressure drop for incompressible, fixed-flow-rate systems). Direct comparison of the raw objective function $J$ between FEniCS and Ansys requires careful interpretation. The FEniCS objective integrates both physical viscous dissipation and the artificial Brinkman penalty across a diffuse field. Conversely, Ansys evaluates purely physical dissipation on a body-fitted mesh bounded by sharp walls. Therefore, while the FEniCS objective history demonstrates the internal convergence of the optimizer, the ultimate validation metric relies on the post-processed hydraulic power loss calculated on the extracted Ansys geometries.

Geometry extraction inherently introduces uncertainties. The continuous design variable $\bar{\gamma}$ is thresholded at a specific iso-value to define the solid-fluid interface. Any discrepancies in dissipation between FEniCS and Ansys stem from a combination of discretization differences, solver architecture, and this thresholding process itself.

\section{Benchmark 1: The Diffuser}
\label{sec:benchmark_diffuser}
The diffuser serves as the foundational benchmark to evaluate separation control in a purely expanding flow, isolating shear layer dynamics from curvature-induced secondary flows. 

\subsection{FEniCS-SA Optimization Results}
The optimization relies on the frozen turbulence adjoint to suppress recirculation zones and minimize the pressure drop across the expansion. [Insert text analyzing the objective function plot: note how rapidly the objective drops initially before asymptoting]. 

Figure [X] illustrates the final projected topology, alongside the velocity magnitude and eddy viscosity fields. [Insert critical analysis of the SA field: Does the optimizer successfully maintain attached flow? Does the eddy viscosity field indicate an over-prediction of diffusion that artificially stabilizes the flow?]

\subsection{Cross-Code Verification and Model Sensitivity}
The thresholded geometry was imported into Ansys Fluent. Table [X] details the power dissipation evaluated across the different solvers and turbulence models. 

Comparing FEniCS-SA to Ansys-SA reveals a relative difference of [X]\%. This discrepancy highlights the impact of moving from a Brinkman-penalized diffuse interface to a wall-resolved, body-fitted boundary. When analyzed with the SST $k-\omega$ model, the total dissipation shifts by [X]\%. [Insert critical analysis: Does the SST model predict flow separation that the SA model missed during optimization? If the SA design fails to perform well under SST conditions, acknowledge this openly as a limitation of the one-equation closure in expanding adverse pressure gradients].

\section{Benchmark 2: The Double Pipe}
\label{sec:benchmark_double_pipe}
The double pipe configuration introduces complex merging and splitting dynamics, assessing the optimizer's ability to mitigate internal junction losses and mixing dissipation without the complications of global curvature.

\subsection{FEniCS-SA Optimization Results}
[Insert analysis of FEniCS results. Focus on the central junction. How does the optimizer route the flow to avoid stagnation zones? Reference the $\nu_t$ field to show where turbulence production is localized.]

\subsection{Cross-Code Verification and Model Sensitivity}
Table [X] presents the comparative CFD metrics for the double pipe geometry. [Insert analysis of the Ansys results. Evaluate if the sharp geometry extraction preserved the smooth junction routing. Compare SA vs SST $k-\omega$: Are the mixing losses at the intersection heavily dependent on the chosen turbulence model?]

\section{Benchmark 3: The Pipe Bend}
\label{sec:benchmark_pipe_bend}
To strictly evaluate the framework's performance in curvature-dominated regimes, the pipe bend benchmark proposed by Alexandersen et al. (2026) is utilized. This case incorporates explicit non-design regions and forces the fluid through a severe geometric deflection, triggering strong streamline curvature.

\subsection{Dean Number and Theoretical Limits}
As established in Chapter~\ref{ch:turbulence_modeling}, the standard Boussinesq approximation utilized by the SA model struggles to accurately capture the anisotropic stress tensors generated by strong secondary flows. The operational Dean number ($De$) for this benchmark is calculated as [Insert Value], based on the characteristic inlet diameter and the mean radius of curvature generated by the optimizer.

\subsection{Optimization and Re-analysis}
[Insert analysis of the FEniCS TO results]. 

Re-analysis in Ansys provides critical insight into the physical limits of the SA-optimized design. As shown in Table [X], the performance gap between Ansys-SA and Ansys SST $k-\omega$ is [Insert Result]. [If the gap is large: This discrepancy confirms the theoretical limitations outlined in Chapter~\ref{ch:turbulence_modeling}. The SA-driven optimizer, blind to curvature-induced anisotropy, generates a topology that underperforms when subjected to higher-fidelity turbulence models. This is not a failure of the numerical implementation, but a demonstrated physical limit of utilizing a one-equation linear eddy viscosity model for highly curved manifolds.]

\section{Benchmark 4: The U-Bend (Advanced Separation)}
\label{sec:benchmark_u_bend}
% Note: Include this section only if you successfully complete the Alexandersen U-bend case.
The U-bend benchmark forces a 180-degree flow reversal, combining intense curvature with inevitable flow separation at the inner radius. [Present the FEniCS results and Ansys comparisons. Focus on whether the optimizer managed to delay separation, and whether SST $k-\omega$ validates or rejects the SA-generated inner wall contour].

\section{Evaluation of the Coupled ``Semi-Frozen'' Adjoint}
\label{sec:semifrozen_evaluation}
While the frozen adjoint serves as the primary optimization engine due to its numerical tractability, an advanced ``semi-frozen'' adjoint was formulated to explicitly capture turbulence production sensitivities. 

However, evaluating this coupled formulation reveals significant practical limitations. As observed in Figure [X], while the coupled adjoint mathematically detects shear-driven turbulence, exact continuous adjoints for turbulence models are notoriously susceptible to numerical noise at domain boundaries. Even with non-design regions implemented, the coupled sensitivities frequently become dominated by artificial shear concentrations at fluid-solid transitions. Consequently, the optimizer tends to prioritize mitigating these localized mathematical artifacts rather than optimizing the global fluid routing. Therefore, despite its theoretical superiority, the frozen adjoint remains the more robust and reliable approach for the manifolds studied in this thesis.
% ... and so on until
\chapter{Turbulent Topology Optimization Results}
\label{ch:benchmark_results}

\section{Introduction}
\label{sec:ch6_introduction}
Building upon the FEniCS-SA solver verified in Chapter~\ref{ch:sa_verification} and the topology optimization methodology developed in Chapter~\ref{ch:topology_optimization_methodology}, this chapter evaluates the turbulent topology optimization solver on the pipe-bend and U-bend internal-flow benchmarks from Bayat et al.~\cite{Bayat2026} and the diffuser benchmark from Yoon~\cite{Yoon2016}. Table~\ref{tab:ch6_benchmark_cases} gives the operating points and validation analyses. Each benchmark section then introduces only the geometry, objective, and interpretation needed for that case.

The optimized layouts are interpreted through projected density fields, velocity fields, and independent sharp-geometry CFD validation in Ansys Fluent. Appendix~\ref{app:C_topology_optimization_benchmarks} retains the supporting frozen-sensitivity verification, FEniCS convergence histories, numerical settings, and Ansys grid-independence studies, grouped by benchmark case.

\section{Benchmark Cases and Evaluation Metrics}
\label{sec:ch6_strategy_metrics}

The three benchmarks cover curved internal flows and a diffuser with a prescribed outlet contraction. The Bayat cases use reciprocal wall-distance penalization, topology-dependent SA and wall-distance reaction terms, an average inlet-pressure objective, and a frozen-turbulence adjoint. The diffuser uses the same frozen-SA optimization machinery, but it minimizes the dissipation objective $J_D$.

\begin{table}[!ht]
	\centering
	\footnotesize
	\caption{Benchmark cases evaluated in Chapter~\ref{ch:benchmark_results}. Detailed physical and numerical parameters are listed in Appendix~\ref{app:C_topology_optimization_benchmarks}.}
	\label{tab:ch6_benchmark_cases}
	\renewcommand{\arraystretch}{1.2}
		\begin{tabularx}{\textwidth}{l c c >{\raggedright\arraybackslash}X}
		\hline
		\textbf{Case} & \textbf{$V_f$} & \textbf{$Re$} & \textbf{Analyses reported} \\
		\hline
			Pipe bend~\cite{Bayat2026} & $0.25$ & $5{,}000$ & FEniCS optimization; Ansys-SA and Ansys SST $k-\omega$ CFD validation; laminar-design comparison \\
			U-bend~\cite{Bayat2026} & $0.27$ & $5{,}000$ & FEniCS optimization; Ansys-SA and Ansys SST $k-\omega$ CFD validation; laminar-design comparison \\
			Diffuser~\cite{Yoon2016} & $0.30$ & $3{,}000$ & Velocity-outlet baseline; pressure-outlet FEniCS optimization with $J_D$; Ansys-SA and Ansys SST $k-\omega$ CFD validation; laminar-design comparison for the pressure-outlet case \\
		\hline
	\end{tabularx}
\end{table}

The pipe-bend and U-bend cases minimize the average inlet-pressure objective $J_p$ defined in Eq.~\ref{eq:inlet_pressure_objective}, which keeps their setup comparable with the Bayat benchmark objective. The diffuser case minimizes the dissipation objective $J_D$ defined in Eq.~\ref{eq:dissipation_objective}, matching the formulation used by Yoon~\cite{Yoon2016}. All extracted geometries are then evaluated as sharp-wall Ansys Fluent models, where the Brinkman penalty is removed and the validation metrics are recomputed on the exported flow domain.

The primary validation metric is the physical viscous dissipation integrated over the full exported CFD domain, including the non-design inlet, outlet, and fixed regions that are retained in the final geometry,
\begin{equation}
	\Phi_D = \int_{\Omega}
	\frac{1}{2} \mu_{\mathrm{eff}}
	\left(\nabla\mathbf{u}+\nabla\mathbf{u}^{T}\right):
	\left(\nabla\mathbf{u}+\nabla\mathbf{u}^{T}\right)
	\,d\Omega ,
	\label{eq:ch6_validation_dissipation}
\end{equation}
where $\Omega$ denotes the complete sharp-geometry flow domain exported to Ansys and $\mathbf{S}=\frac{1}{2}\left(\nabla\mathbf{u}+\nabla\mathbf{u}^{T}\right)$ is the symmetric strain-rate tensor. The integrand in Eq.~\ref{eq:ch6_validation_dissipation} can therefore also be written as $2\mu_{\mathrm{eff}}\mathbf{S}:\mathbf{S}$. In the Ansys Fluent post-processing, the scalar strain-rate magnitude is $\dot{\gamma}=\sqrt{2\mathbf{S}:\mathbf{S}}$, so the custom field function $\mu_{\mathrm{eff}}\dot{\gamma}^{2}$ is equivalent to the integrand in Eq.~\ref{eq:ch6_validation_dissipation}. Because the reported Ansys validations all re-solve the extracted geometries as turbulent flows, the effective viscosity used in these post-processed metrics is $\mu_{\mathrm{eff}}=\mu+\mu_t$.

The second global validation metric is the static pressure drop between the physical inlet and outlet sections,
\begin{equation}
	\Delta p =
	\bar{p}_{\mathrm{in}}-\bar{p}_{\mathrm{out}},
	\qquad
	\bar{p}_{\Gamma} =
	\frac{1}{|\Gamma|}
	\int_{\Gamma} p_s\,d\Gamma ,
	\label{eq:ch6_static_pressure_drop}
\end{equation}
where $p_s$ is the static pressure and $|\Gamma|=\int_{\Gamma}1\,d\Gamma$ is the boundary length in the two-dimensional model. In Ansys Fluent, $p_s$ is taken directly from the area-weighted average static pressure on the inlet and outlet surfaces. For FEniCS results that use kinematic pressure, the corresponding static pressure is obtained as $p_s=\rho p_k$ before evaluating Eq.~\ref{eq:ch6_static_pressure_drop}. Pressure drop is retained throughout Chapter~\ref{ch:benchmark_results} because it is the most direct engineering interpretation of the same mechanical losses measured by $\Phi_D$. All two-dimensional dissipation quantities are evaluated per unit depth.

\section{Pipe Bend Benchmark}
\label{sec:alexandersen_pipebend}
The pipe-bend benchmark introduces a direct comparison with the turbulent topology optimization results of Bayat et al.~\cite{Bayat2026}. While the reference study employs a standard $k$-$\varepsilon$ RANS model with wall functions, the present framework tackles the same problem using a wall-resolved SA model under the frozen-turbulence assumption. This comparison tests whether the fundamental topological solution, a streamlined and separation-avoiding bend, remains consistent across different turbulence closures.

\subsection{Geometry and Boundary Conditions}
\label{subsec:alexandersen_pipebend_setup}
The boundary conditions and geometric constraints replicate the setup shown in Fig.~10 of Bayat et al.~\cite{Bayat2026}, specifically targeting the $V_f=0.25$ volume-fraction case. The inlet is placed on the left side of the non-design extension, the outlet is placed along the lower boundary, and the outlet pressure is fixed to zero. The optimization minimizes the average pressure on this prescribed inlet boundary required to drive the flow through the evolving bend.
The local benchmark geometry used for the present implementation is shown in Figure~\ref{fig:alexandersen_pipebend_setup}.

\begin{figure}[!ht]
	\centering
	\includegraphics[width=0.85\textwidth]{PipeBend-Geometry-BCs}
	\caption{Pipe-bend benchmark geometry and boundary conditions, adapted from Fig.~10 of Bayat et al.~\cite{Bayat2026} \protect\ChatGPTFigureRef.}
	\label{fig:alexandersen_pipebend_setup}
\end{figure}

\subsection{Optimization Results}
\label{subsec:alexandersen_pipebend_fenics}
Figure~\ref{fig:alexandersen_pipebend_fenics_results} illustrates the projected density field $\bar{\gamma}$ and the corresponding velocity magnitude. The supporting coordinate finite-difference and Taylor checks for the frozen adjoint are reported in Appendix~\ref{app:C_topology_optimization_benchmarks}. The FEniCS-SA result shows the same broad routing strategy as the $k$-$\varepsilon$ reference from Bayat et al.~\cite{Bayat2026}: the main high-speed path bypasses the sharpest corners to reduce momentum loss. A direct overlay is provided in Figure~\ref{fig:app_C_alexandersen_pipebend_paper_overlay}. The exact gray-scale density distribution near the walls still differs because the turbulence closure and near-wall treatment are not the same.

\begin{figure}[!ht]
	\centering
	\begin{subfigure}[b]{0.47\textwidth}
		\centering
		\includegraphics[width=\textwidth]{PipeBend-Turbulent-Topology-Jp}
		\caption{Projected density field $\bar{\gamma}$.}
		\label{fig:alexandersen_pipebend_topology}
	\end{subfigure}
	\hfill
	\begin{subfigure}[b]{0.49\textwidth}
		\centering
		\includegraphics[width=\textwidth]{PipeBend-Turbulent-Velocity-FEniCS-Jp}
		\caption{Velocity magnitude and vector field.}
		\label{fig:alexandersen_pipebend_velocity}
	\end{subfigure}
	\caption{FEniCS-SA optimization result for the pipe-bend benchmark at $V_f=0.25$. Final objective value $J_p= 0.2130$ Pa.}
	\label{fig:alexandersen_pipebend_fenics_results}
\end{figure}

For the same $V_f=0.25$, $Re=5{,}000$ operating point, Bayat et al.~\cite{Bayat2026} report $J_p=0.4023$ Pa for their conventional $k$-$\varepsilon$ topology optimization. This value is used only as a reference scale, not as a direct validation target, because the turbulence closure and near-wall treatment differ from the present wall-resolved SA formulation.

The reported FEniCS value is the endpoint of a short $70$-iteration run without continuation. It is used here as a provisional design state for comparing topology and validation metrics, not as a converged optimum. The convergence history in Figure~\ref{fig:alexandersen_pipebend_histories} shows that the normalized pressure objective is still monotonically decreasing, so the run is promising but should be extended with more MMA iterations and continuation.


\subsection{CFD Validation}
\label{subsec:alexandersen_pipebend_cfd}
The extracted $70$-iteration pipe-bend topology is evaluated in Table~\ref{tab:alexandersen_pipebend_validation_metrics}. Because the SST $k-\omega$ model is notably more sensitive to adverse pressure gradients than SA, evaluating the extracted geometry under both closures serves as a stress test for the provisional shape.

\begin{table}[!ht]
	\centering
	\footnotesize
	\caption{CFD validation metrics for the extracted pipe-bend topology.}
	\label{tab:alexandersen_pipebend_validation_metrics}
	\renewcommand{\arraystretch}{1.2}
	\begin{tabularx}{\textwidth}{l c c c X}
		\hline
		\textbf{Solver} & \textbf{Model} & \textbf{$\Phi_D$ [W/m]} & \textbf{$\Delta p$ [Pa]} & \textbf{Comment} \\
		\hline
		FEniCS diffuse design & SA & 0.0138 & 0.2130 & 70-iteration non-continuation state \\
		Ansys sharp geometry & SA & 0.0153 & 0.1457 & Cross-code comparison \\
		Ansys sharp geometry & SST $k-\omega$ & 0.0166 & 0.1522 & Turbulence-model sensitivity \\
		\hline
	\end{tabularx}
\end{table}

Using the Ansys-SA result as the reference, the SST $k-\omega$ evaluation increases the pressure drop by approximately $4.5\%$ and the viscous dissipation by approximately $8.5\%$. The SST $k-\omega$ model therefore predicts slightly higher losses, consistent with a closure that is more sensitive to adverse pressure gradients, but the changes remain modest. This suggests that the extracted pipe-bend geometry is not narrowly tuned to the SA model. It does not, however, establish that the underlying FEniCS optimization has converged.

\subsection{Comparison with Laminar Design}
\label{subsec:alexandersen_pipebend_laminar}
The laminar reference uses the same pipe-bend mesh, design-domain restrictions, inlet and outlet treatment, objective type, and volume fraction, but replaces the SA turbulent state equations by the laminar flow formulation and is optimized at $Re=1$. This laminar optimization favors a sharper, shorter turn that is appropriate for the creeping-flow reference problem but less suitable once the extracted geometry is operated at the target turbulent Reynolds number. The comparison is therefore interpreted under identical turbulent operating conditions, rather than by comparing the laminar and turbulent optimization objectives directly.

Figure~\ref{fig:alexandersen_pipebend_laminar_topologies} compares the laminar and turbulent optimized topologies before sharp-geometry extraction. The laminar-design mesh-independence study supporting the SST $k-\omega$ result is reported in Table~\ref{tab:app_C_alexandersen_pipebend_lam_mesh} of Appendix~\ref{app:C_topology_optimization_benchmarks}.

\begin{figure}[!ht]
	\centering
	\begin{subfigure}[b]{0.47\textwidth}
		\centering
		\includegraphics[width=\textwidth]{PipeBend-Laminar-Topology-Jp}
		\caption{Projected density field $\bar{\gamma}$: laminar-optimized design ($Re=1$).}
		\label{fig:alexandersen_pipebend_laminar_topo}
	\end{subfigure}
	\hfill
	\begin{subfigure}[b]{0.47\textwidth}
		\centering
		\includegraphics[width=\textwidth]{PipeBend-Turbulent-Topology-Jp}
		\caption{Projected density field $\bar{\gamma}$: turbulent-optimized design ($Re=5{,}000$).}
		\label{fig:alexandersen_pipebend_turbulent_topo}
	\end{subfigure}
	\caption{Laminar and turbulent topology comparison for the pipe-bend benchmark.}
	\label{fig:alexandersen_pipebend_laminar_topologies}
\end{figure}


Figure~\ref{fig:alexandersen_pipebend_laminar_velocity_comparison} shows the corresponding SST $k-\omega$ velocity contours through the extracted laminar and turbulent pipe-bend designs. These contours provide the flow-field evidence behind the scalar metrics: sharper turns and corners in the laminar design promote stronger local separation and momentum loss when the geometry is operated at the turbulent Reynolds number.

\begin{figure}[!ht]
	\centering
	\begin{subfigure}[b]{0.47\textwidth}
		\centering
		\includegraphics[width=\textwidth]{PipeBend-Laminar-VelocityMagnitude-SSTkomega}
		\caption{Velocity magnitude field $u$: laminar-optimized design.}
		\label{fig:alexandersen_pipebend_laminar_velocity_sst}
	\end{subfigure}
	\hfill
	\begin{subfigure}[b]{0.47\textwidth}
		\centering
		\includegraphics[width=\textwidth]{PipeBend-Turbulent-VelocityMagnitude-SSTkomega}
		\caption{Velocity magnitude field $u$: turbulent-optimized design.}
		\label{fig:alexandersen_pipebend_turbulent_velocity_sst}
	\end{subfigure}
	\caption{Velocity-contour comparison for the extracted pipe-bend designs under identical turbulent operating conditions (SST $k-\omega$ model).}
	\label{fig:alexandersen_pipebend_laminar_velocity_comparison}
\end{figure}

The final SST $k-\omega$ comparison in Table~\ref{tab:alexandersen_pipebend_laminar_perf} evaluates the extracted laminar and turbulent pipe-bend designs under identical turbulent operating conditions. The turbulent design reduces the viscous dissipation from $0.0313$ to $0.0166~\mathrm{W/m}$ and the pressure drop from $0.5444$ to $0.1522~\mathrm{Pa}$. These reductions, $47.0\%$ and $72.0\%$ respectively, show that the turbulent topology optimization framework outperforms the laminar surrogate for the target turbulent operating point.

\begin{table}[!ht]
	\centering
	\footnotesize
	\caption{Performance comparison of the extracted pipe-bend topologies under identical turbulent operating conditions.}
	\label{tab:alexandersen_pipebend_laminar_perf}
	\renewcommand{\arraystretch}{1.2}
	\begin{tabularx}{\textwidth}{X c c c X}
		\hline
		\textbf{Design origin} & \textbf{Model} & \textbf{$\Phi_D$ [W/m]} & \textbf{$\Delta p$ [Pa]} & \textbf{Relative performance} \\
		\hline
		Laminar TO, $Re=1$ & SST $k-\omega$ & $0.0313$ & $0.5444$ & Baseline \\
		Turbulent TO, SA & SST $k-\omega$ & $0.0166$ & $0.1522$ & $47.0\%$ lower $\Phi_D$; $72.0\%$ lower $\Delta p$ \\
		\hline
	\end{tabularx}
\end{table}

\section{U-Bend Benchmark}
\label{sec:alexandersen_ubend}
The U-bend is the companion benchmark to the pipe bend. It uses the same frozen-SA topology optimization framework, but the geometry forces a $180^\circ$ return path between two ports on the same side of the domain. This makes the case useful for checking whether the optimizer can form a smooth turbulent return duct without relying on the shorter $90^\circ$ routing available in the pipe-bend case.

\subsection{Geometry and Boundary Conditions}
\label{subsec:alexandersen_ubend_setup}
The geometry and boundary conditions follow Fig.~11 and Table~4 of Bayat et al.~\cite{Bayat2026}. The square design region has side length $L$, with a non-design lead duct of length $0.2L$ on the left. The inlet and outlet ports both have height $0.2L$; the top port acts as the inlet and the lower port acts as the pressure outlet. A fixed internal separator bar of thickness $0.10L$ extends from the lead duct into the design domain, forcing the flow to turn around its rounded tip. The target volume fraction is $V_f=0.27$, and the Reynolds number is $Re=5{,}000$ based on the port height and maximum inlet velocity.
Figure~\ref{fig:alexandersen_ubend_setup} shows the corresponding local geometry and boundary-condition layout used for this benchmark.

\begin{figure}[!ht]
	\centering
	\includegraphics[width=0.85\textwidth]{UBend-TO-Geometry-BCs}
	\caption{Geometry and boundary conditions for the U-bend benchmark, adapted from Fig.~11 and Table~4 of Bayat et al.~\cite{Bayat2026} \protect\ChatGPTFigureRef.}
	\label{fig:alexandersen_ubend_setup}
\end{figure}

\subsection{Optimization Results}
\label{subsec:alexandersen_ubend_fenics}
Figure~\ref{fig:alexandersen_ubend_fenics_results} reports the final projected density and velocity-magnitude field for the U-bend. The main feature to assess is whether the design forms a rounded return passage around the separator tip, rather than a short path with tight local curvature. Because the objective is average inlet pressure, a successful layout reduces the pressure required to maintain the imposed inlet flow while respecting the prescribed volume fraction.

\begin{figure}[!ht]
	\centering
	\begin{subfigure}[b]{0.45\textwidth}
		\centering
		\includegraphics[width=\textwidth]{UBend-Turbulent-Topology-Jp}
		\caption{Projected density field $\bar{\gamma}$.}
		\label{fig:alexandersen_ubend_topology}
	\end{subfigure}
	\hfill
	\begin{subfigure}[b]{0.53\textwidth}
		\centering
		\includegraphics[width=\textwidth]{UBend-Turbulent-Velocity-FEniCS-Jp}
		\caption{Velocity magnitude and vector field.}
		\label{fig:alexandersen_ubend_velocity}
	\end{subfigure}
	\caption{FEniCS-SA optimization result for the U-bend benchmark at $V_f=0.27$. Final objective value $J_p=0.8250~\mathrm{Pa}$.}
	\label{fig:alexandersen_ubend_fenics_results}
\end{figure}

For the matching $V_f=0.27$, $Re=5{,}000$ U-bend case, the conventional $k$-$\varepsilon$ result reported by Bayat et al.~\cite{Bayat2026} is $J_p=0.8327$. As for the pipe bend, this value provides a reference scale only; the comparison remains qualitative because the present solver uses a different turbulence model and near-wall treatment.

The turbulent U-bend topology should also be read in the same way as the provisional pipe-bend topology: it is an extracted optimization endpoint from the available run, not a fully converged final optimum. The available U-bend convergence history is reported in Figure~\ref{fig:alexandersen_ubend_histories} of Appendix~\ref{app:C_topology_optimization_benchmarks}, using the same combined objective and volume-residual format as the pipe-bend diagnostic.

\subsection{CFD Validation}
\label{subsec:alexandersen_ubend_cfd}
The extracted U-bend topology is validated with the same sharp-geometry workflow as the pipe bend. The Ansys-SA solve checks the design under the same closure family used in the optimization loop, while the Ansys SST $k-\omega$ solve tests whether the extracted turbulent design remains robust under a two-equation model that is more sensitive to separation and adverse pressure gradients. Table~\ref{tab:alexandersen_ubend_validation_metrics} summarizes the FEniCS-SA optimization diagnostic and the two sharp-geometry Ansys validation solves for the turbulent U-bend design; the supporting validation mesh hierarchy is documented in Table~\ref{tab:app_B_alexandersen_ubend_turb_mesh} of Appendix~\ref{app:C_topology_optimization_benchmarks}.

\begin{table}[!ht]
	\centering
	\footnotesize
	\caption{CFD validation metrics for the extracted turbulent U-bend topology.}
	\label{tab:alexandersen_ubend_validation_metrics}
	\renewcommand{\arraystretch}{1.2}
	\begin{tabularx}{\textwidth}{l c c c X}
		\hline
		\textbf{Solver} & \textbf{Model} & \textbf{$\Phi_D$ [W/m]} & \textbf{$\Delta p$ [Pa]} & \textbf{Comment} \\
		\hline
		FEniCS diffuse design & SA & $0.0936$ & $0.8250$ & Optimization-state diagnostic \\
		Ansys sharp geometry & SA & $0.0785$ & $0.8071$ & Cross-code comparison with the optimization closure family \\
		Ansys sharp geometry & SST $k-\omega$ & $0.1016$ & $1.1973$ & Independent turbulence-closure check \\
		\hline
	\end{tabularx}
\end{table}

Using the Ansys-SA result as the reference, the SST $k-\omega$ validation predicts larger losses for the same extracted turbulent geometry: $\Phi_D$ increases by approximately $29.4\%$ and $\Delta p$ by approximately $48.3\%$. This trend is physically plausible because the fixed separator forces a compact return flow, where SST $k-\omega$ is more responsive to adverse pressure gradients and local recirculation downstream of the separator. The U-bend validation is therefore interpreted as a closure-sensitivity result for the forced $180^\circ$ return flow rather than as closure-independent agreement.

\subsection{Comparison with Laminar Design}
\label{subsec:alexandersen_ubend_laminar}
The laminar reference uses the same U-bend mesh, inlet and outlet ports, separator geometry, objective type, and volume fraction, but it is optimized at $Re=1$ without the SA turbulence model. The extracted laminar and turbulent layouts are then evaluated under the same target turbulent operating conditions. This comparison tests whether the SA-optimized turbulent return path remains advantageous when the sharp geometry is re-solved with an independent turbulence closure.

Figure~\ref{fig:alexandersen_ubend_laminar_topologies} compares the laminar and turbulent optimized U-bend topologies. The laminar-design mesh-independence study supporting the SST $k-\omega$ result is reported in Table~\ref{tab:app_C_alexandersen_ubend_lam_mesh} of Appendix~\ref{app:C_topology_optimization_benchmarks}.

\begin{figure}[!ht]
	\centering
	\begin{subfigure}[b]{0.47\textwidth}
		\centering
		\includegraphics[width=\textwidth]{UBend-Laminar-Topology-Jp}
		\caption{Projected density field $\bar{\gamma}$: laminar-optimized design ($Re=1$).}
		\label{fig:alexandersen_ubend_laminar_topo}
	\end{subfigure}
	\hfill
	\begin{subfigure}[b]{0.50\textwidth}
		\centering
		\includegraphics[width=\textwidth]{UBend-Turbulent-Topology-Jp}
		\caption{Projected density field $\bar{\gamma}$: turbulent-optimized design ($Re=5{,}000$).}
		\label{fig:alexandersen_ubend_turbulent_topo}
	\end{subfigure}
	\caption{Laminar and turbulent topology comparison for the U-bend benchmark.}
	\label{fig:alexandersen_ubend_laminar_topologies}
\end{figure}

Figure~\ref{fig:alexandersen_ubend_laminar_velocity_comparison} shows SST $k-\omega$ velocity contours through the extracted U-bend designs, plotted with the same color scale. The fields show broadly similar return-flow organization, with a high-speed outer path around the bend and a low-momentum band along the inner side downstream of the separator. The close-up vector diagnostic in Figure~\ref{fig:app_C_alexandersen_ubend_recirc_sst} of Appendix~\ref{app:C_topology_optimization_benchmarks} shows local recirculation in this low-momentum region, so the velocity fields are interpreted as qualitative evidence of curvature-induced separation and recovery rather than as the primary performance metric.

\begin{figure}[!ht]
	\centering
	\begin{subfigure}[b]{0.47\textwidth}
		\centering
		\includegraphics[width=\textwidth]{UBend-Laminar-VelocityMagnitude-SSTkomega}
		\caption{Velocity magnitude field $u$: laminar-optimized design.}
		\label{fig:alexandersen_ubend_laminar_velocity_sst}
	\end{subfigure}
	\hfill
	\begin{subfigure}[b]{0.47\textwidth}
		\centering
		\includegraphics[width=\textwidth]{UBend-Turbulent-VelocityMagnitude-SSTkomega}
		\caption{Velocity magnitude field $u$: turbulent-optimized design.}
		\label{fig:alexandersen_ubend_turbulent_velocity_sst}
	\end{subfigure}
	\caption{SST $k-\omega$ velocity-contour comparison for the extracted U-bend designs under identical turbulent operating conditions.}
	\label{fig:alexandersen_ubend_laminar_velocity_comparison}
\end{figure}

Table~\ref{tab:alexandersen_ubend_laminar_perf} gives the corresponding laminar-design comparison under identical turbulent operating conditions.

\begin{table}[!ht]
	\centering
	\footnotesize
	\caption{Performance comparison of the extracted U-bend topologies under identical turbulent operating conditions.}
	\label{tab:alexandersen_ubend_laminar_perf}
	\renewcommand{\arraystretch}{1.2}
	\begin{tabularx}{\textwidth}{X c c c X}
		\hline
		\textbf{Design origin} & \textbf{Validation model} & \textbf{$\Phi_D$ [W/m]} & \textbf{$\Delta p$ [Pa]} & \textbf{Relative performance} \\
		\hline
		Laminar TO, $Re=1$ & SA & $0.1137$ & $0.9646$ & Baseline under SA \\
		Turbulent TO, SA & SA & $0.0785$ & $0.8071$ & $31.0\%$ lower $\Phi_D$; $16.3\%$ lower $\Delta p$ \\
		Laminar TO, $Re=1$ & SST $k-\omega$ & $0.1069$ & $1.1496$ & Baseline under SST $k-\omega$ \\
		Turbulent TO, SA & SST $k-\omega$ & $0.1016$ & $1.1973$ & $5.0\%$ lower $\Phi_D$; $4.1\%$ higher $\Delta p$ \\
		\hline
	\end{tabularx}
\end{table}

Under Ansys-SA, the turbulent U-bend design gives the expected improvement over the laminar reference. Under SST $k-\omega$, however, the two designs are much closer: the turbulent design has a slightly lower viscous dissipation but a slightly higher pressure drop. The U-bend therefore supports the same cautious message as the velocity fields: the SA-optimized return path is beneficial under the design closure, but the $180^\circ$ turn remains sensitive to the turbulence model used for sharp-geometry validation.

\section{Diffuser Benchmark}
\label{sec:yoon_diffuser}
The diffuser benchmark follows the third numerical example of Yoon~\cite{Yoon2016}. It adds a downstream-boundary-condition sensitivity to the closure-sensitivity checks already used for the bend cases. A velocity-outlet baseline follows the published setup, while the primary result replaces the prescribed outlet velocity profile with a pressure outlet so that the outlet velocity distribution is solved rather than imposed.

\subsection{Velocity-Outlet Boundary-Condition Sensitivity}
\label{subsec:yoon_diffuser_velocity_outlet}
The velocity-outlet baseline prescribes parabolic velocity profiles at both the full-height inlet and the center-third outlet. It produces a more strongly diverging channel near the outlet, which is interpreted as partly boundary-condition-driven. Its geometry, FEniCS fields, and Ansys diagnostic contour are reported in Appendix~\ref{app:C_topology_optimization_benchmarks}; the pressure-outlet variant is used as the primary diffuser result below.

\subsection{Pressure-Outlet Setup and Optimization Results}
\label{subsec:yoon_diffuser_pressure_setup}
\label{subsec:yoon_diffuser_pressure_fenics}
The pressure-outlet variant changes only the downstream boundary treatment relative to the velocity-outlet baseline documented in Appendix~\ref{app:C_topology_optimization_benchmarks}. Its local geometry and boundary-condition layout are shown in Figure~\ref{fig:yoon_diffuser_pressure_setup}. The inlet remains the shifted parabolic velocity profile with $U_{\max,\mathrm{in}}=3~\mathrm{m/s}$ and the same $Re=3{,}000$ definition, while the center outlet is assigned a fixed static pressure $p=0~\mathrm{Pa}$. No outlet velocity profile is prescribed. This makes the outlet mass flux and velocity distribution part of the solution and removes the direct Dirichlet forcing that affected the velocity-outlet topology.

\begin{figure}[!ht]
	\centering
	\includegraphics[width=0.95\textwidth]{Diffuser-Geometry-BCs-PreOut}
	\caption{Pressure-outlet diffuser geometry and boundary conditions \protect\ChatGPTFigureRef.}
	\label{fig:yoon_diffuser_pressure_setup}
\end{figure}

Figure~\ref{fig:yoon_diffuser_pressure_fenics_results} reports the pressure-outlet FEniCS-SA topology and velocity field. Compared with the velocity-outlet result in Appendix~\ref{app:C_topology_optimization_benchmarks}, the outlet region is less strongly widened. The optimized pressure-outlet objective is also lower, decreasing from $J_D=3.1164\times10^5~\mathrm{W/m}$ to $J_D=2.7664\times10^5~\mathrm{W/m}$, or by about $11.2\%$. This is not a direct performance ranking between equivalent designs, because the boundary-value problem has changed. The conclusion is narrower: removing the outlet Dirichlet velocity profile reduces the amount of geometry change required near the outlet and gives a less forced diffuser design.

\begin{figure}[!ht]
	\centering
	\begin{subfigure}[b]{0.47\textwidth}
		\centering
		\includegraphics[width=\textwidth]{Diffuser-Turbulent-Topology-PreOut}
		\caption{Projected density field $\bar{\gamma}$.}
		\label{fig:yoon_diffuser_pressure_topology}
	\end{subfigure}
	\hfill
	\begin{subfigure}[b]{0.50\textwidth}
		\centering
		\includegraphics[width=\textwidth]{Diffuser-Turbulent-Velocity-FEniCS-PreOut}
		\caption{Velocity magnitude and vector field.}
		\label{fig:yoon_diffuser_pressure_velocity}
	\end{subfigure}
	\caption{FEniCS-SA optimization result for the pressure-outlet diffuser at $V_f=0.30$. Final objective value $J_D=2.7664\times10^5~\mathrm{W/m}$.}
	\label{fig:yoon_diffuser_pressure_fenics_results}
\end{figure}

The corresponding pressure-outlet convergence history through this common comparison point is reported in Figure~\ref{fig:yoon_diffuser_pressure_convergence} of Appendix~\ref{app:C_topology_optimization_benchmarks}.

\subsection{Pressure-Outlet CFD Validation}
\label{subsec:yoon_diffuser_pressure_cfd}
The extracted pressure-outlet topology is validated with the same sharp-geometry workflow as the pipe bend and U-bend. The Ansys-SA solve checks the transfer from the diffuse FEniCS design to a sharp-wall geometry under the same closure family, while the Ansys SST $k-\omega$ solve tests whether the ranking is robust under an independent two-equation model. Because the optimization objective is $J_D$, the validation emphasizes both physical viscous dissipation and static pressure drop.

\begin{table}[!ht]
	\centering
	\footnotesize
	\caption{CFD validation metrics for the extracted pressure-outlet diffuser topology.}
	\label{tab:yoon_diffuser_validation_metrics}
	\renewcommand{\arraystretch}{1.2}
		\begin{tabularx}{\textwidth}{l c c >{\raggedright\arraybackslash}X}
		\hline
		\textbf{Evaluation} & \textbf{$\Phi_D$ [W/m]} & \textbf{$\Delta p_s$ [Pa]} & \textbf{Interpretation} \\
		\hline
		FEniCS-SA diffuse topology & $2.6307\times10^5$ & $2.6022\times10^5$ & Diffuse-domain diagnostic; total objective $J_D=2.7664\times10^5~\mathrm{W/m}$ \\
		Ansys-SA sharp geometry & $1.1885\times10^5$ & $1.1541\times10^5$ & Sharp-wall validation with the optimization closure family \\
		Ansys SST $k-\omega$ sharp geometry & $1.2834\times10^4$ & $5.0006\times10^4$ & Independent closure check; much lower modeled dissipation \\
		\hline
	\end{tabularx}
\end{table}

The FEniCS-SA values in Table~\ref{tab:yoon_diffuser_validation_metrics} are approximately $2.21$ times larger in $\Phi_D$ and $2.25$ times larger in $\Delta p_s$ than the Ansys-SA sharp-geometry values, so this case cannot be claimed as quantitative cross-code agreement. The result is interpreted as a transfer limitation between the diffuse Brinkman-domain optimization diagnostic and the sharp-wall extracted geometry.

The Ansys validation rows are also strongly closure-sensitive: for the turbulent pressure-outlet design, the SA pressure drop is $2.31$ times the SST $k-\omega$ value, and the SA dissipation is $9.26$ times the SST $k-\omega$ value. This is the opposite trend from the U-bend: here SST $k-\omega$ predicts lower modeled losses, while SA appears more loss-producing in the adverse-pressure-gradient diffuser flow. Because the split is large, it is treated as model-form uncertainty rather than evidence that either closure gives a closure-independent physical truth. Additional velocity-magnitude, turbulent-viscosity-ratio, and local dissipation-density fields were inspected but are not retained as main figures because they do not change the interpretation. The appendix wall-shear diagnostic and SST $k-\omega$ vector close-up in Figures~\ref{fig:app_C_yoon_diffuser_wall_shear_sign_map} and~\ref{fig:app_C_yoon_diffuser_separation_sst} identify local reverse-flow regions in the pressure-outlet geometry; the integrated pressure-drop and dissipation metrics remain the primary evidence for this diffuser.

\subsection{Pressure-Outlet Laminar-Design Comparison}
\label{subsec:yoon_diffuser_pressure_laminar}
The laminar reference for the final diffuser comparison uses the same pressure-outlet boundary treatment, objective type, and volume fraction, but it is optimized at $Re=1$ without the SA turbulence model. The extracted laminar and turbulent pressure-outlet layouts are then evaluated under the same target turbulent operating conditions. This comparison isolates whether the SA-optimized pressure-outlet design remains beneficial when the final sharp geometry is re-solved with the same closure and with SST $k-\omega$.

Table~\ref{tab:yoon_diffuser_pressure_laminar_perf} compares the extracted laminar and turbulent pressure-outlet diffuser topologies under identical turbulent operating conditions. The topology and velocity fields were inspected, but they do not by themselves reveal the performance ranking clearly enough to justify a separate figure. Because the diffuser is optimized with $J_D$, lower $\Phi_D$ is the primary performance indicator, while the static pressure drop remains the most direct hydraulic interpretation of the extracted sharp geometry. Both SA and SST $k-\omega$ validation solves are included to quantify whether the final ranking changes with the validation model.

\begin{table}[!ht]
	\centering
	\footnotesize
	\caption{Performance comparison of the extracted pressure-outlet diffuser topologies under identical turbulent operating conditions.}
	\label{tab:yoon_diffuser_pressure_laminar_perf}
	\renewcommand{\arraystretch}{1.2}
		\begin{tabularx}{\textwidth}{l c c >{\raggedright\arraybackslash}X}
		\hline
		\textbf{Evaluation} & \textbf{$\Phi_D$ [W/m]} & \textbf{$\Delta p_s$ [Pa]} & \textbf{Relative performance} \\
		\hline
		Laminar TO, Ansys-SA & $1.3495\times10^5$ & $1.3979\times10^5$ & Baseline under SA \\
		Turbulent TO, Ansys-SA & $1.1885\times10^5$ & $1.1541\times10^5$ & $11.9\%$ lower $\Phi_D$; $17.4\%$ lower $\Delta p_s$ \\
		Laminar TO, Ansys SST $k-\omega$ & $1.1289\times10^4$ & $3.8332\times10^4$ & Baseline under SST $k-\omega$ \\
		Turbulent TO, Ansys SST $k-\omega$ & $1.2834\times10^4$ & $5.0006\times10^4$ & $13.7\%$ higher $\Phi_D$; $30.5\%$ higher $\Delta p_s$ \\
		\hline
	\end{tabularx}
\end{table}

The laminar-design comparison reinforces this closure-sensitive interpretation. Under Ansys-SA, the turbulent topology reduces the pressure drop by $17.4\%$ and the viscous dissipation by $11.9\%$ relative to the pressure-outlet laminar topology. Under SST $k-\omega$, the ranking reverses: the turbulent topology gives a $30.5\%$ higher pressure drop and a $13.7\%$ higher dissipation. The appendix wall-shear and vector diagnostics identify where local reverse flow appears in the pressure-outlet geometry, but the case-specific conclusion comes from the scalar ranking reversal. The diffuser is therefore retained as a closure-sensitivity case for the present SA-based framework.

\section{Conclusion}
\label{sec:ch6_conclusion}
This chapter answers the third research question by evaluating whether the turbulent topology optimization framework produces designs that remain credible after independent CFD validation. The pipe bend provides the strongest validation evidence among the reported cases: the Appendix~\ref{app:C_topology_optimization_benchmarks} frozen-sensitivity checks confirm the derivative supplied to MMA for $J_p$, and the extracted $70$-iteration topology remains robust when validated in Ansys Fluent with both SA and SST $k-\omega$. The FEniCS optimization endpoint itself is provisional rather than converged, but the extracted geometry is not narrowly tuned to the SA model.

The U-bend extends the assessment to a forced $180^\circ$ return path and gives a more cautious result. Like the pipe-bend topology, the turbulent U-bend topology is an optimization endpoint rather than a fully converged final optimum, so the result is interpreted as validation of the extracted geometry and not as a convergence claim. For the turbulent topology alone, replacing Ansys-SA by SST $k-\omega$ increases $\Phi_D$ by $29.4\%$ and $\Delta p$ by $48.3\%$. Compared with the laminar reference, the turbulent design improves both metrics under SA, but under SST $k-\omega$ it gives only a $5.0\%$ lower dissipation and a $4.1\%$ higher pressure drop. The appendix close-up vector diagnostic explains this case-specific sensitivity through local recirculation in the low-momentum inner return region.

The diffuser adds a separate caution: outlet treatment and validation closure both affect the final ranking. The velocity-outlet baseline shows that a prescribed outlet profile can influence the optimized outlet geometry, so the pressure-outlet variant is used as the primary diffuser result. For this pressure-outlet diffuser, the turbulent SA design improves on the laminar design under Ansys-SA, reducing $\Delta p_s$ by $17.4\%$ and $\Phi_D$ by $11.9\%$. The same two geometries evaluated with SST $k-\omega$ give the opposite ranking, with the turbulent design increasing $\Delta p_s$ by $30.5\%$ and $\Phi_D$ by $13.7\%$. Overall, the pipe bend provides the cleanest transfer case, while the U-bend and diffuser show why extracted turbulent designs must still be checked under an independent closure before making closure-independent claims.

\chapter{Conclusion and Future Work}
\label{ch:conclusion}

\section{Conclusion}
The objective of this thesis was to build and assess a flow-only turbulent topology optimization framework for internal configurations relevant to energy and thermal-management systems. The framework does not solve the energy equation and does not yet optimize conjugate heat transfer. Instead, it addresses the hydraulic prerequisite for those applications: generating low-loss internal flow paths under turbulent operating conditions. This narrower scope is deliberate, because turbulent momentum transport, wall-distance treatment, Brinkman penalization, geometry extraction, and adjoint sensitivity assembly already form a difficult optimization problem.

The work was organized as a sequence of checks. Chapter~\ref{ch:turbulence_modeling} assessed the useful range of the Spalart-Allmaras (SA) closure. Chapters~\ref{ch:fenics_implementation} and~\ref{ch:sa_verification} implemented and verified the FEniCS-SA solver. Chapters~\ref{ch:topology_optimization_methodology} and~\ref{ch:benchmark_results} embedded that solver in a density-based topology optimization loop with Brinkman penalization, filtering, projection, MMA, topology-dependent SA and wall-distance penalties, a frozen-turbulence adjoint, sensitivity checks, and independent sharp-geometry CFD validation.

\subsection*{Research Question 1}
\textbf{How accurately does the Spalart-Allmaras turbulence model predict curved and separated internal turbulent flows, and where do its physical limitations appear?}
\\ \\
The answer is conditional. SA is defensible as a practical design-model closure for attached wall-bounded and mildly curved internal flows, including the moderate-curvature VMFL048 U-bend evidence. It is not closure-independent truth for strongly curved or separated flows.

The Ansys comparisons in this thesis make that boundary concrete. The VMFL048 U-bend at $De\approx15{,}000$ remains credible, while the tighter VMFL030 pipe bend at $De\approx18{,}000$ already shows clear SA deterioration in the selected profile comparison. The backward-facing-step evidence gives the same warning for separated recovery: SA can predict plausible integral behavior, but SST $k-\omega$ follows the local velocity recovery more closely. The limiting issue is therefore curvature- and separation-driven Reynolds-stress anisotropy under a scalar eddy-viscosity closure, not a single Reynolds-number threshold.

\subsection*{Research Question 2}
 \textbf{Does the custom FEniCS implementation of the Spalart-Allmaras model reproduce Ansys Fluent predictions for the selected turbulent internal-flow benchmarks?}
 \\ \\
 Within the two-dimensional, wall-resolved scope of the thesis, the answer is yes. The solver reproduces the governing ingredients needed by the optimizer: incompressible flow coupling, SA transport, eddy-viscosity reconstruction, wall-distance evaluation, and near-wall resolution.

In the periodic channel, the fine-mesh bulk velocity differs from the Ansys reference by $2.67\%$, while the exported profile errors are $E_{\infty}=2.66\%$ and $E_{L_2}=2.64\%$. In the U-bend verification case, the pressure drop differs from Ansys by $2.79\%$, and the bend-exit profile gives $E_{\infty}=8.50\%$ and $E_{L_2}=2.75\%$. These results verify the FEniCS implementation as a sufficiently consistent forward solver for the topology optimization studies, while leaving the model-form limitations from Chapter~\ref{ch:turbulence_modeling} in place.

\subsection*{Research Question 3}
\textbf{To what extent do SA-driven topology optimizations produce final designs whose predicted performance remains consistent under independent Ansys Fluent validation with SA and SST $k-\omega$, and how do they compare with laminar-optimized designs under identical turbulent operating conditions?}
\\ \\
Chapter~\ref{ch:benchmark_results} demonstrates that the reliability of an SA-optimized design depends fundamentally on the flow physics dictated by the target geometry.

The pipe bend represents the cleanest transfer case, showing that the framework successfully produces robust, closure-independent designs in regimes dominated by attached flow and moderate curvature. Even as a provisional $70$-iteration endpoint, the extracted pipe-bend topology remains highly credible: replacing Ansys-SA by SST $k-\omega$ increases pressure drop by only about $4.5\%$ and viscous dissipation by $8.5\%$. Under identical SST $k-\omega$ operating conditions, the turbulent design decisively outperforms the laminar-optimized reference, reducing pressure drop by $72.0\%$ and viscous dissipation by $47.0\%$. Furthermore, the derivative supplied to MMA is internally verified, with coordinate finite differences giving a maximum relative discrepancy of $5.92\times10^{-5}$ and the directional Taylor test showing approximately second-order behavior.

Conversely, the U-bend and diffuser benchmarks expose the severe limitations of relying on a one-equation closure for topology optimization in separation-dominated flows. Because the quantitative results for these cases are highly sensitive to the validation model, the exact performance metrics lose their predictive value. For geometries forcing strong recirculation (such as the $180^\circ$ return of the U-bend) or adverse pressure gradients (the diffuser), the optimizer heavily tunes the layout to the SA model's specific, and often overly forgiving, prediction of flow separation and recovery. When these extracted shapes are cross-examined using SST $k-\omega$, which is notably more responsive to adverse pressure gradients, the predicted flow structures disagree and the performance deteriorates significantly.

Consequently, the advantage of a turbulent-optimized design over a laminar baseline is not absolute. In attached flows, the turbulent formulation is a clear necessity and delivers massive performance gains. But in heavily separated regimes, the optimized geometry becomes so deeply coupled to the SA model's simplifying assumptions that the turbulent designs offer only inconclusive, marginal gains over a simple laminar surrogate, or can even perform worse when evaluated under a different turbulence closure.

\subsection*{Overall Conclusion}
The central conclusion is balanced. A frozen-SA density-based topology optimization framework can produce credible low-loss turbulent internal-flow design candidates when the flow remains within the closure's practical range and when the extracted geometry is independently validated. It is not a substitute for final CFD validation, experimental testing, or a higher-fidelity turbulence model in separated, strongly three-dimensional flows. Independent sharp-geometry CFD validation is the mandatory engineering step between an optimized density field and anything that could be manufactured.

Within its validated envelope, the framework can still be valuable as an early-stage design and research tool for internal hydraulic layouts where pressure loss or dissipated mechanical energy is the main design constraint. It can generate turbulence-aware topology candidates, compare them with laminar surrogate designs, test objective and penalization choices, and expose whether a candidate remains robust when the turbulence closure is changed. The contribution of this thesis is therefore a reproducible bridge between turbulence-model knowledge, open-source finite-element optimization, and engineering judgment about when a turbulent topology-optimized design is credible.

\section{Limitations}
The limitations below define how far the conclusions can be taken. They are separated from the future-work actions so that the same point is not repeated twice.
\begin{enumerate}
	\item \textbf{Simplified turbulence physics and two-dimensional scope:} The framework uses a wall-resolved one-equation SA closure in 2D planar domains. This suppresses Dean vortices, secondary motion, Reynolds-stress anisotropy, and three-dimensional separation structures that can be important in physical internal-flow devices.
	\item \textbf{Frozen adjoint and sensitivities:} The finite-difference and Taylor checks verify the derivative used by MMA under the same frozen-turbulence assumption. They do not verify the derivative of a fully coupled turbulence, wall-distance, and geometry-extraction pipeline. The optimizer is therefore internally consistent with its chosen approximation, but the approximation remains a modeling decision.
	\item \textbf{Diffuse-to-sharp geometry transfer:} The optimized design is a continuous Brinkman-penalized density field, while validation is performed on an extracted sharp-wall geometry. The chosen threshold and smoothing procedure affect wall placement, local curvature, and separation behavior, so FEniCS-Ansys discrepancies cannot always be attributed to solver physics alone.
	\item \textbf{Validation level:} Independent Ansys Fluent validation with SA and SST $k-\omega$ is a necessary engineering step, but it is still numerical validation. Experimental pressure-drop and flow-field measurements would be required before using these layouts as physically validated device designs.
	\item \textbf{Optimization robustness:} The reported turbulent pipe-bend and U-bend topologies should be interpreted as optimization endpoints rather than converged final results. Although pressure drop and dissipation are physically related in steady internal flows, domain-integrated dissipation can provide a less localized optimization signal than an average inlet-pressure objective.
\end{enumerate}

\section{Future Work}
The verified turbulent fluid mechanics framework developed in this thesis establishes a foundation for several targeted extensions:

\subsection{Benchmark Completion and Optimization Robustness}
The short $70$-iteration, non-continuation pipe-bend run and the $275$-iteration turbulent U-bend endpoint reported in Chapter~\ref{ch:benchmark_results} are useful as provisional validation states, but they should be replaced by converged or deliberately polished runs before making final quantitative claims about the FEniCS optima. Future runs should use a longer MMA budget together with continuation of the RAMP, projection, move-limit, and topology-wall penalty parameters. In particular, the turbulent-design U-bend should be repeated as a more robust optimization campaign tuned for the forced $180^\circ$ return-flow geometry. The same future campaign should rerun the frozen-sensitivity diagnostic on the exact continuation settings used for the final design.

\subsection{Conjugate Heat Transfer (CHT) and Thermal Management}
The ultimate engineering application for high-Reynolds topology optimization is not merely fluid routing, but active thermal management. Historically, the immense computational cost of resolving Conjugate Heat Transfer (CHT) has forced optimization frameworks to rely on laminar approximations (e.g., natural convection as demonstrated by Alexandersen et al. \cite{Alexandersen2014}) or simplified flow proxies such as Darcy models \cite{Zhao2018}. Even when high-resolution, density-based solvers are successfully deployed for microelectronics cooling, they typically circumvent the complexities of turbulent momentum transport \cite{Barroca2023}. 

However, high-performance heat sinks still require models that capture turbulent convection. Recent work has introduced turbulence into thermal topology optimization using extended finite element (XFEM) level-set approaches with the SA model \cite{Maute2023}, computationally reduced multi-layer thermofluid approximations \cite{Zhao2021}, and standard density-based frameworks coupled to two-equation $k-\epsilon$ closures \cite{Sun2023}.

With the turbulent momentum and SA transport equations implemented and verified within the current FEniCS architecture, the framework provides a practical starting point for turbulent CHT optimization. Extending the current methodology to CHT will require implementing a turbulent Prandtl number ($Pr_t$) approach to model the turbulent heat flux vector, alongside deriving the corresponding thermal adjoint sensitivities to minimize peak temperatures or maximize overall heat transfer coefficients. Furthermore, the capacity for optimizing two-fluid heat exchangers, currently dominated by laminar assumptions \cite{Andreasen2020}, could be improved by adapting this turbulent solver.

\subsection{Three-Dimensional Topologies}
As noted in the limitations, turbulent secondary flows are inherently 3D phenomena. To fully evaluate the geometric limits of the turbulence models and design physically realizable internal-flow components, the current 2D FEniCS implementation must be extended to three dimensions. The mathematical structure of the governing equations and adjoint derivations remains similar, but the computational cost changes by orders of magnitude because wall-resolved turbulent optimization requires fine near-wall resolution over a three-dimensional design domain. Large-scale topology optimization studies have demonstrated that parallel finite-element optimization at extreme design-variable counts is possible when the formulation is paired with efficient domain decomposition and scalable solvers \cite{Aage2017}. A 3D turbulent extension of the present work should therefore prioritize parallel mesh generation, robust algebraic multigrid preconditioning, and automated sharp-geometry extraction.

\subsection{Adjoint Model Hierarchies}
This thesis deliberately uses the frozen-turbulence adjoint as the robust baseline. Future frameworks should investigate stable ways to add selected SA transport sensitivities, especially in separated flows where the constant-eddy-viscosity assumption is least accurate. Existing turbulent topology optimization studies already span several levels of adjoint complexity, from frozen or constant-eddy-viscosity approximations \cite{Dilgen2018,Zymaris2009} to more complete RANS sensitivity formulations and implementation studies \cite{Margetis2024,Wu2023}. A practical next step is therefore not necessarily to jump directly to a fully coupled turbulence adjoint, but to build an adjoint hierarchy: frozen SA sensitivities as the baseline, selected SA transport sensitivities as an intermediate model, and two-equation closures such as SST $k-\omega$ for cases where curvature, separation, and anisotropic shear dominate the design response.


% If you have appendices:
\appendixpage*          % if wanted
\appendix
\include{appendix-A-SA-Validity-Cases-Details}
% ... and so on until
\chapter{Spalart-Allmaras FEniCS Solver Support}
\label{app:sa_fenics_solver_support}

This appendix supports Chapter~\ref{ch:fenics_implementation}. It collects standard Spalart--Allmaras closure definitions used by the FEniCS implementation, while keeping the main chapter focused on the weak forms, Picard coupling, and solver structure.

\section{SA Closure Functions and Calibration Constants}
\label{app:sa_closure_terms}
The FEniCS implementation uses the vorticity-based Spalart--Allmaras closure definition. The production term relies on a modified vorticity magnitude $\tilde{S}$ that maintains correct near-wall scaling while avoiding loss of production when the local vorticity becomes small \cite{Spalart1994,NasaTurb}. First,
\begin{equation}
	S = \sqrt{2\Omega_{ij}\Omega_{ij}}
	\quad \text{and} \quad
	\bar{S} = \frac{\tilde{\nu}}{\kappa^2 y^2} f_{v2}.
\end{equation}
Here $S$ is the scalar magnitude of the vorticity tensor, $\kappa$ is the von Karman constant, and $y$ is the distance to the nearest wall. During nonlinear iterations, the code evaluates the protected form of $\tilde{S}$ recommended in the SA reference implementation \cite{NasaTurb} so that the production term remains well defined when $S+\bar{S}$ becomes small. Small numerical floors are added in the implementation to prevent division by zero. The function $f_{v2}$ is a near-wall modifier:
\begin{equation}
	f_{v2} = 1 - \frac{\chi}{1 + \chi f_{v1}} \quad \text{where} \quad \chi = \frac{\tilde{\nu}}{\nu} \quad \text{and} \quad f_{v1} = \frac{\chi^3}{\chi^3 + c_{v1}^3}.
\end{equation}

The wall-destruction term is modulated by the non-dimensional function $f_w$,
\begin{equation}
	f_w = g \left( \frac{1 + c_{w3}^6}{g^6 + c_{w3}^6} \right)^{1/6},
\end{equation}
where
\begin{equation}
	g = r + c_{w2}(r^6 - r) \quad \text{and} \quad r = \min \left( \frac{\tilde{\nu}}{\tilde{S} \kappa^2 y^2}, 10 \right).
\end{equation}
The upper bound of $10$ on $r$ improves numerical stability during solver initialization when $\tilde{S}$ may momentarily approach zero.

The original formulation by Spalart and Allmaras \cite{Spalart1994} includes additional geometric trip functions ($f_{t1}$ and $f_{t2}$) designed to trigger transition from laminar to turbulent flow at prescribed coordinates. Following the topology-optimization methodology of Yoon~\cite{Yoon2016}, these transition functions are omitted here because the topological boundaries are unknown a priori and evolve during optimization.

The empirical calibration constants are
\begin{align*}
	c_{b1} &= 0.1355 & \sigma &= 2/3 & c_{b2} &= 0.622 \\
	\kappa &= 0.41 & c_{w2} &= 0.3 & c_{w3} &= 2.0 \\
	c_{v1} &= 7.1 & c_{v2} &= 0.7 & c_{v3} &= 0.9 \\
	c_{w1} &= \frac{c_{b1}}{\kappa^2} + \frac{1 + c_{b2}}{\sigma}.
\end{align*}

\chapter{Spalart-Allmaras Solver Verification Details}
\label{app:B_sa_solver_verification}

This appendix supports Chapter~\ref{ch:sa_verification}. It contains concise FEniCS setup tables, the mesh-independence studies used for the channel benchmark, and the selected FEniCS and Ansys mesh information used for the U-bend verification case. For traceability, the active configuration file is named once near each FEniCS setup; the tables retain only the governing geometry, boundary-condition, and solver settings rather than duplicating the full Python configuration. The backward-facing-step material is placed in Appendix~\ref{app:A_backstep_benchmark}, because it supports the Chapter~\ref{ch:turbulence_modeling} turbulence-model discussion rather than the Chapter~\ref{ch:sa_verification} forward-solver verification.

Unless stated otherwise, all grid-independence relative differences in this appendix are reported as positive percentage magnitudes. For a mesh-dependent metric $\phi$, the coarse--medium difference is computed as $|\phi_{\mathrm{coarse}}-\phi_{\mathrm{medium}}|/|\phi_{\mathrm{coarse}}|\times100\%$, while the medium--fine difference is computed as $|\phi_{\mathrm{medium}}-\phi_{\mathrm{fine}}|/|\phi_{\mathrm{fine}}|\times100\%$.

\section{Verification Benchmark 1: Channel Flow}
\label{sec:app_A_channel}

\subsection{Physical and Geometric Parameters}
\label{subsec:app_A_channel_physical}
The first verification benchmark consists of a pressure-driven periodic 2D channel. To align with fundamental turbulence research conventions, the characteristic length is defined as the full channel height ($H$), and the Reynolds number ($Re_H$) is computed strictly based on this geometric parameter rather than the hydraulic diameter. The active FEniCS setup is summarized from \texttt{ConfigChannel\_SpalartAllmaras\_Steady.py}, and the key physical setup is detailed in Table~\ref{tab:app_A_channel_physical}.

\begin{table}[htbp]
	\centering
	\footnotesize
	\caption{Geometric, physical, and boundary-condition parameters for the periodic channel flow benchmark.}
	\label{tab:app_A_channel_physical}
	\renewcommand{\arraystretch}{1.2}
	\begin{tabularx}{\textwidth}{l X}
		\hline
		\textbf{Parameter} & \textbf{Channel flow setting} \\
		\hline
		Benchmark type & Pressure-driven, streamwise-periodic 2D channel \\
		Domain dimensions & Length $L=2.0$ m; full channel height $H=2.0$ m \\
		Characteristic length & Full channel height $H=2.0$ m \\
		Reference velocity & $U_{\mathrm{ref}}=18.5$ m/s, utilized to define the target Reynolds number \\
		Kinematic viscosity & $\nu = 1.81818 \times 10^{-3}$ m$^2$/s \\
		Reynolds number & $Re_H \approx 2.03 \times 10^4$ \\
			Body force & $\mathbf{f}=(1.0,0)$ m/s$^2$ in the kinematic-pressure convention \\
			Driving condition & Streamwise body-force mode, equivalent to a pressure drop of $2$ Pa over the periodic length $L=2$ m \\
			Velocity boundary conditions & Periodic velocity between inlet/outlet planes; no-slip on top and bottom walls \\
			Pressure boundary conditions & Periodic pressure field with no inlet/outlet pressure Dirichlet condition; pressure mean normalized during IPCS \\
			SA boundary conditions & Periodic $\tilde{\nu}$ between inlet/outlet planes; $\tilde{\nu}=0$ at walls \\
			Initial fields & $\mathbf{u}_0=(18.5,0)$ m/s, $p_0=0$ Pa, $\tilde{\nu}_0=5.0 \times 10^{-3}$ m$^2$/s \\
		Verification mesh & \texttt{Fine\_FEniCS}; first layer $\Delta y_1 = 1.709 \times 10^{-3}$ m \\
		\hline
	\end{tabularx}
\end{table}

\subsection{Numerical Simulation Parameters}
\label{subsec:app_A_channel_numerical}
The channel flow was resolved using an Incremental Pressure Correction Scheme (IPCS) embedded within a Picard fixed-point iteration loop. The key numerical solver settings, convergence tolerances, and relaxation factors used in FEniCS are summarized in Table~\ref{tab:app_A_channel_numerical}.

\begin{table}[htbp]
	\centering
	\scriptsize
	\caption{Numerical simulation parameters for the FEniCS steady Spalart-Allmaras channel benchmark.}
	\label{tab:app_A_channel_numerical}
	\renewcommand{\arraystretch}{1.2}
	\begin{tabularx}{\textwidth}{l c X}
		\hline
		\textbf{Parameter} & \textbf{Symbol} & \textbf{Setting} \\
		\hline
		Velocity space & $\mathbf{u}$ & Vector CG2, periodic in streamwise direction \\
		Pressure space & $p$ & Scalar CG1 \\
		SA working-variable space & $\tilde{\nu}$ & Scalar CG1, periodic in streamwise direction \\
		Quadrature degree & -- & $2$ \\
			Maximum outer Picard steps & $N_{\mathrm{Picard,max}}$ & $400$ \\
			SA solves per Picard step & -- & $1$ \\
			Outer velocity tolerance & $\epsilon_{\mathbf{u}}$ & $1.0 \times 10^{-5}$ \\
			Outer pressure tolerance & $\epsilon_p$ & $1.0 \times 10^{-6}$ \\
			Outer SA tolerance & $\epsilon_{\tilde{\nu}}$ & $1.0 \times 10^{-6}$ \\
			Outer diagnostic convergence & -- & Midline bulk velocity window of $100$ Picard steps with relative tolerance $1.0\times10^{-3}$ \\
			IPCS pseudo-time step & $\Delta t$ & $5.0 \times 10^{-3}$ s \\
			Maximum IPCS steps per Picard step & -- & $500$ \\
			IPCS velocity tolerance & -- & $1.0 \times 10^{-6}$ \\
			IPCS pressure tolerance & -- & $1.0 \times 10^{-6}$; pressure convergence is diagnostic in body-force mode \\
			IPCS pressure normalization & -- & Enabled; the pressure correction is mean-normalized \\
		Velocity relaxation factor & $\alpha_u$ & $0.3$ \\
		Pressure relaxation factor & $\alpha_p$ & $0.3$ \\
		SA relaxation factor & $\alpha_{\tilde{\nu}}$ & $0.3$ \\
		SA lower bound & $\tilde{\nu}_{\min}$ & $1.0 \times 10^{-12}$ \\
		Wall-distance equation & -- & Relaxed Eikonal, $\sigma_w=0.1$, $G_0=20$ \\
		Wall-distance Newton settings & -- & Relative tolerance $1.0 \times 10^{-8}$; absolute tolerance $1.0 \times 10^{-10}$; maximum $80$ iterations; relaxation $0.5$ \\
		Linear solver strategy & -- & IPCS blocks: direct MUMPS; SA transport: default solver \\
		\hline
	\end{tabularx}
\end{table}

\subsection{Spatial Verification and Grid Independence}
\label{subsec:app_A_channel_mesh}
To guarantee that the numerical results were independent of the spatial discretization, a formal grid independence study was conducted. Three distinct boundary-layer inflated meshes (Coarse, Medium, and Fine) were evaluated. The geometric properties, node counts, and resolution parameters of this mesh hierarchy are explicitly detailed in Table~\ref{tab:app_A_channel_meshes}. Crucially, the first-layer height was kept constant across all grids to strictly enforce the $y^+ \approx 1$ requirement, while the bulk and inflation-layer resolutions were systematically refined. 

\begin{table}[htbp]
	\centering
	\footnotesize
	\caption{Wall-resolved channel mesh hierarchy for the FEniCS grid-independence study.}
	\label{tab:app_A_channel_meshes}
	\renewcommand{\arraystretch}{1.2}
	\setlength{\tabcolsep}{5pt}
	\begin{tabular*}{\textwidth}{@{\extracolsep{\fill}}lccc@{}}
		\hline
		\textbf{Parameter} & \textbf{Coarse} & \textbf{Medium} & \textbf{Fine} \\
		\hline
		\multicolumn{4}{@{}l}{\textit{Mesh settings}} \\
		Streamwise cells & $10$ & $20$ & $40$ \\
		Wall-normal cells & $42$ & $51$ & $68$ \\
		$\Delta y_1$ [m] & $1.709 \times 10^{-3}$ & $1.709 \times 10^{-3}$ & $1.709 \times 10^{-3}$ \\
		$\Delta y_{\min}$ [m] & $1.709 \times 10^{-3}$ & $1.709 \times 10^{-3}$ & $1.709 \times 10^{-3}$ \\
		$\Delta y_{\max}$ [m] & $7.707 \times 10^{-2}$ & $8.454 \times 10^{-2}$ & $7.589 \times 10^{-2}$ \\
		Inflation layers per wall & $10$ & $14$ & $18$ \\
		Growth ratio & $1.45$ & $1.35$ & $1.25$ \\
		Vertices & $473$ & $1{,}092$ & $2{,}829$ \\
		Elements & $840$ & $2{,}040$ & $5{,}440$ \\
		\hline
		\multicolumn{4}{@{}l}{\textit{Grid metric}} \\
		$U_{\mathrm{bulk}}$ [m/s] & $16.923$ & $18.489$ & $18.472$ \\
		$\epsilon_{\mathrm{rel}}(U_{\mathrm{bulk}})$ [\%] & -- & $9.25$ & $0.09$ \\
		\hline
	\end{tabular*}
	\vspace{2pt}
	\begin{minipage}{0.98\textwidth}
		\scriptsize \textit{Note.} For the $\epsilon_{\mathrm{rel}}$ row, the Medium column gives the Coarse--Medium change and the Fine column gives the Medium--Fine change.
	\end{minipage}
\end{table}

Visual confirmation of this spatial convergence is provided by the wall-normal velocity profiles plotted in Figure~\ref{fig:app_A_channel_mesh_independence}. In each case, the profile is extracted at the streamwise mid-plane, $x=0.5L$, with the wall-normal coordinate $y$ varied across the full channel height. The \texttt{Fine\_FEniCS} grid was consequently selected for the final verification against the Ansys Fluent baseline, while the coarser grids demonstrate that the wall-normal velocity profile is already asymptotically converged.

\begin{figure}[htpb]
	\centering
	\includegraphics[width=\textwidth]{Channel-FEniCS-Mesh-Independence}
	\caption{FEniCS wall-normal velocity profiles at $x=0.5L$ for the Coarse, Medium, and Fine grids.}
	\label{fig:app_A_channel_mesh_independence}
\end{figure}

Furthermore, to ensure a rigorous cross-code validation, an equivalent grid independence study was performed for the commercial Ansys Fluent baseline. The Ansys domain was discretized using a comparable mapped mesh hierarchy, maintaining a strict $y^+ \approx 1$ resolution at the walls. In addition to the residual monitors, the bulk velocity was monitored and required to change by less than $1.0 \times 10^{-3}$ over $250$ iterations before accepting a converged solution. The continuity residual criterion was set to $1.0 \times 10^{-5}$, while the momentum and SA working-variable residual criteria were set to $1.0 \times 10^{-6}$. Mass conservation was also checked from the inlet and outlet mass-flow reports. The properties of these meshes are documented in Table~\ref{tab:app_A_channel_meshes_ansys}. The wall-normal velocity profiles extracted from the Ansys solver at the same mid-channel location, illustrated in Figure~\ref{fig:app_A_channel_mesh_indep_ansys}, exhibited identical asymptotic convergence, confirming that the baseline reference data is mathematically robust and completely independent of spatial discretization errors.

\begin{table}[htbp]
	\centering
	\footnotesize
	\caption{Wall-resolved channel mesh hierarchy used for the Ansys grid-independence study.}
	\label{tab:app_A_channel_meshes_ansys}
	\renewcommand{\arraystretch}{1.2}
	\setlength{\tabcolsep}{5pt}
	\begin{tabular*}{\textwidth}{@{\extracolsep{\fill}}lccc@{}}
		\hline
		\textbf{Parameter} & \textbf{Coarse} & \textbf{Medium} & \textbf{Fine} \\
		\hline
		\multicolumn{4}{@{}l}{\textit{Mesh settings}} \\
		Mesh type & MultiZone & MultiZone & MultiZone \\
		Face type & All quad & All quad & All quad \\
		$N_x$ & $64$ & $160$ & $220$ \\
		$N_y$ & $80$ & $144$ & $180$ \\
		Bias type & Two-sided & Two-sided & Two-sided \\
		Bias factor & $59$ & $28$ & $21$ \\
		Approx. $\Delta y_1$ [m] & $1.700 \times 10^{-3}$ & $1.700 \times 10^{-3}$ & $1.700 \times 10^{-3}$ \\
		Nodes & $7{,}740$ & $26{,}448$ & $43{,}010$ \\
		Elements & $7{,}565$ & $26{,}123$ & $42{,}594$ \\
		\hline
		\multicolumn{4}{@{}l}{\textit{Grid metric}} \\
		$U_{\mathrm{bulk}}$ [m/s] & $18.305$ & $18.114$ & $17.977$ \\
		$\epsilon_{\mathrm{rel}}(U_{\mathrm{bulk}})$ [\%] & -- & $1.04$ & $0.76$ \\
		\hline
	\end{tabular*}
	\vspace{2pt}
	\begin{minipage}{0.98\textwidth}
		\scriptsize \textit{Note.} $N_x$ and $N_y$ are the streamwise and wall-normal edge divisions. For the $\epsilon_{\mathrm{rel}}$ row, the Medium column gives the Coarse--Medium change and the Fine column gives the Medium--Fine change.
	\end{minipage}
\end{table}

\begin{figure}[htpb]
	\centering
	\includegraphics[width=\textwidth]{Channel-Ansys-Mesh-Independence}
	\caption{Ansys wall-normal velocity profiles at $x=0.5L$ for the Coarse, Medium, and Fine grids.}
	\label{fig:app_A_channel_mesh_indep_ansys}
\end{figure}

\FloatBarrier

\section{Verification Benchmark 2: U-Bend Geometry}
\label{sec:app_A_ubend}

\subsection{Physical and Geometric Parameters}
\label{subsec:app_A_ubend_physical}
The second verification benchmark evaluates a highly convective U-bend. Because this FEniCS domain is a 2D planar approximation of the symmetry mid-plane of the 3D Ansys VMFL048 case, the 3D physical pipe diameter ($D$) is explicitly retained as the characteristic length scale. This ensures the Reynolds ($Re_D$) and Dean ($De$) numbers remain non-dimensionally consistent with the commercial baseline. The straight inlet and outlet legs are shorter than in the original VMFL048 Ansys mesh to reduce meshing and computational effort; both the FEniCS and Ansys runs use this same short-leg U-bend geometry. The active FEniCS setup is summarized from \texttt{ConfigUBend\_SpalartAllmaras\_Steady.py}, and the configuration is detailed in Table~\ref{tab:app_A_ubend_physical}.

\begin{table}[htbp]
	\centering
	\footnotesize
	\caption{Geometric, physical, and boundary-condition parameters for the U-bend benchmark.}
	\label{tab:app_A_ubend_physical}
	\renewcommand{\arraystretch}{1.2}
	\begin{tabularx}{\textwidth}{l X}
		\hline
		\textbf{Parameter} & \textbf{U-bend setting} \\
		\hline
		Benchmark type & Short-leg 2D symmetry-plane U-bend based on Ansys VMFL048 \\
		Domain dimensions & Pipe diameter $D=0.028$ m; centreline bend radius $R_c=0.125$ m; straight-leg length $30D=0.840$ m \\
		Characteristic length & Physical pipe diameter $D=0.028$ m \\
		Reference / inlet velocity & Uniform prescribed inlet velocity $\mathbf{u}_{in}=(0,-1.42)$ m/s \\
		Kinematic viscosity & $\nu = 8.9 \times 10^{-7}$ m$^2$/s \\
		Reynolds number & $Re_D \approx 4.47 \times 10^4$ \\
		Dean number & $De \approx 1.50 \times 10^4$ \\
		Body force & $\mathbf{f}=(0,0)$ m/s$^2$ \\
		Driving condition & Prescribed inlet velocity and zero relative static pressure at the outlet \\
		Velocity boundary conditions & No-slip on all curved and straight pipe walls; natural velocity condition at pressure outlet \\
		Pressure boundary conditions & $p=0$ Pa at outlet; pressure left free at inlet and walls \\
		SA boundary conditions & $\tilde{\nu}_{in}=9.34 \times 10^{-5}$ m$^2$/s at inlet; $\tilde{\nu}=0$ at walls; natural condition at outlet \\
			Initial fields & $\mathbf{u}_0=(0,-1.42)$ m/s, $p_0=0$ Pa, $\tilde{\nu}_0=\tilde{\nu}_{in}$ \\
			Verification mesh & \texttt{Medium\_WallResolved}; first layer $\Delta y_1 = 1.20 \times 10^{-5}$ m \\
		\hline
	\end{tabularx}
\end{table}

\subsection{Numerical Simulation Parameters}
\label{subsec:app_A_ubend_numerical}
Due to the adverse pressure gradients and secondary flow separation induced by the streamline curvature, the U-bend represents a significantly more demanding numerical challenge than the channel. Consequently, the FEniCS solver required a smaller pseudo-time step, tighter under-relaxation, and advanced iterative preconditioners to maintain stability, as detailed in Table~\ref{tab:app_A_ubend_numerical}.

\begin{table}[htbp]
	\centering
	\footnotesize
	\caption{Numerical simulation parameters for the highly convective FEniCS U-bend benchmark.}
	\label{tab:app_A_ubend_numerical}
	\renewcommand{\arraystretch}{1.2}
	\begin{tabularx}{\textwidth}{l c X}
		\hline
		\textbf{Parameter} & \textbf{Symbol} & \textbf{Setting} \\
		\hline
		Velocity space & $\mathbf{u}$ & Vector CG2 \\
		Pressure space & $p$ & Scalar CG1 \\
		SA working-variable space & $\tilde{\nu}$ & Scalar CG1 \\
		Quadrature degree & -- & $4$ \\
			Maximum outer Picard steps & $N_{\mathrm{Picard,max}}$ & $1000$ \\
		SA solves per Picard step & -- & $1$ \\
		Outer velocity tolerance & $\epsilon_{\mathbf{u}}$ & $1.0 \times 10^{-4}$ \\
		Outer pressure tolerance & $\epsilon_p$ & $1.0 \times 10^{-4}$ \\
		Outer SA tolerance & $\epsilon_{\tilde{\nu}}$ & $1.0 \times 10^{-6}$ \\
			IPCS pseudo-time step & $\Delta t$ & $1.0 \times 10^{-4}$ s \\
				Maximum IPCS steps per Picard step & -- & $200$ \\
			Maximum IPCS steps for final flow & -- & $600$ \\
			IPCS velocity tolerance & -- & $1.0 \times 10^{-4}$ \\
			IPCS pressure tolerance & -- & $2.0 \times 10^{-3}$ \\
			Pressure-drop convergence monitor & -- & Final-window length $100$ with relative tolerance $5.0\times10^{-4}$; flow convergence required \\
			Velocity relaxation factor & $\alpha_u$ & $0.3$ \\
			Pressure relaxation factor & $\alpha_p$ & $0.1$ \\
			SA relaxation factor & $\alpha_{\tilde{\nu}}$ & $0.01$ \\
			SA pseudo-time stabilization & -- & SA transport pseudo-time step $5.0\times10^{-2}$ s \\
			SA lower bound & $\tilde{\nu}_{\min}$ & $1.0 \times 10^{-12}$ \\
			Wall-distance equation & -- & Relaxed Eikonal, $\sigma_w=0.1$, $G_0=20$ \\
			Wall-distance Newton settings & -- & Relative tolerance $1.0 \times 10^{-8}$; absolute tolerance $1.0 \times 10^{-10}$; maximum $80$ iterations; relaxation $0.5$ \\
			Linear solver strategy & -- & IPCS blocks: GMRES with ILU velocity/correction preconditioning and Hypre-AMG pressure preconditioning; SA transport: default solver \\
			Restart handling & -- & Restart from saved state enabled when checkpoint files are present; otherwise the configured initial fields are used \\
		\hline
	\end{tabularx}
\end{table}

\subsection{Selected FEniCS Mesh and Ansys Grid Independence}
\label{subsec:app_A_ubend_mesh}
	Managing the highly stretched, non-orthogonal inflation layers around the curved arcs of the U-bend required custom mesh generation via Gmsh. The current FEniCS comparison setup points to the wall-resolved \texttt{Medium\_WallResolved} grid documented in Table~\ref{tab:app_A_ubend_meshes}. Similar to the channel benchmark, the wall-adjacent element height is kept at $1.20 \times 10^{-5}$ m to satisfy the low-Reynolds SA wall constraints. A separate coarse--medium--fine FEniCS mesh-convergence study is therefore not claimed for this U-bend verification case.

The reported FEniCS static pressure-drop metric is computed from boundary-measure-weighted mean pressures over the inlet and outlet sections, rather than from local point samples. Let $\Gamma_{\mathrm{in}}$ and $\Gamma_{\mathrm{out}}$ denote the inlet and outlet boundaries, and let $p_k$ be the kinematic pressure solved in FEniCS. The physical static pressure drop reported in pascals is
\begin{equation}
	\Delta p_s
	= \rho\left(\overline{p}_{k,\mathrm{in}}-\overline{p}_{k,\mathrm{out}}\right)
	= \rho\left[
	\frac{1}{|\Gamma_{\mathrm{in}}|}\int_{\Gamma_{\mathrm{in}}} p_k \, \mathrm{d}\Gamma
	-
	\frac{1}{|\Gamma_{\mathrm{out}}|}\int_{\Gamma_{\mathrm{out}}} p_k \, \mathrm{d}\Gamma
	\right],
	\label{eq:app_A_ubend_fenics_delta_p}
\end{equation}
where $|\Gamma|=\int_{\Gamma}1\,\mathrm{d}\Gamma$ is the FEniCS boundary measure and $\rho=997$ kg/m$^3$. In the implementation this corresponds to integrating over the marked inlet and outlet facets with \texttt{ds(2)} and \texttt{ds(3)}, respectively. The density factor is required because the IPCS formulation stores pressure in kinematic form, while the benchmark table reports the static pressure drop in Pa. Although the outlet pressure is prescribed as zero relative pressure, the outlet term is still evaluated as a boundary average so that the diagnostic follows the same inlet-minus-outlet definition on every grid.

	\begin{table}[htbp]
		\centering
		\footnotesize
		\caption{Selected wall-resolved FEniCS U-bend mesh used for comparison with the Ansys Fluent verification baseline.}
		\label{tab:app_A_ubend_meshes}
		\renewcommand{\arraystretch}{1.2}
		\begin{tabularx}{0.82\textwidth}{l X}
			\hline
				\textbf{Parameter} & \textbf{\texttt{Medium\_WallResolved} setting} \\
				\hline
				$h_{\mathrm{bulk}}$ [m] & $1.40 \times 10^{-3}$ \\
				$\Delta y_1$ [m] & $1.20 \times 10^{-5}$ \\
				Inflation layers & $14$ \\
				Growth ratio & $1.30$ \\
				$t_{\mathrm{infl}}$ [m] & $1.53 \times 10^{-3}$ \\
				Vertices & $73{,}978$ \\
				Elements & $144{,}952$ \\
			\hline
		\end{tabularx}
	\end{table}

For the independent Ansys Fluent baseline, a separate mesh-independence study was executed in the same shortened U-bend geometry. Using a parallel unstructured meshing strategy with strict $y^+ \approx 1$ boundary layer inflation (documented in Table~\ref{tab:app_A_ubend_meshes_ansys}), the Ansys velocity profiles at the bend exit across the same pipe-diameter line, shown in Figure~\ref{fig:app_A_ubend_mesh_indep_ansys}, were checked for spatial convergence. For each Ansys solve, convergence was accepted only after the scaled continuity, momentum, and Spalart--Allmaras residuals reached $1.0\times10^{-6}$ and the monitored static pressure drop changed by less than a relative tolerance of $5.0\times10^{-4}$ over the final $100$ iterations. The inlet and outlet mass-flow rates were also compared for each grid to verify the global mass balance before accepting the pressure-drop values. The \texttt{Fine\_Ansys} grid is retained as the commercial verification baseline, while the deliberately coarser \texttt{Coarse\_Ansys} level is selected to bracket the mesh-convergence trend.

\begin{table}[htbp]
	\centering
		\footnotesize
		\caption{Wall-resolved U-bend mesh hierarchy used for the Ansys grid-independence study.}
		\label{tab:app_A_ubend_meshes_ansys}
		\renewcommand{\arraystretch}{1.2}
		\setlength{\tabcolsep}{5pt}
		\begin{tabular*}{\textwidth}{@{\extracolsep{\fill}}lccc@{}}
			\hline
			\textbf{Parameter} & \textbf{Coarse} & \textbf{Medium} & \textbf{Fine} \\
			\hline
			\multicolumn{4}{@{}l}{\textit{Mesh settings}} \\
			Mesh type & Quad-dom. & Quad-dom. & Quad-dom. \\
			$h_{\mathrm{face}}$ [m] & $4.00 \times 10^{-3}$ & $2.00 \times 10^{-3}$ & $1.25 \times 10^{-3}$ \\
			Curvature angle & $22^\circ$ & $12^\circ$ & $8^\circ$ \\
			Inflation option & First layer & First layer & First layer \\
			$\Delta y_1$ [m] & $1.20 \times 10^{-5}$ & $1.20 \times 10^{-5}$ & $1.20 \times 10^{-5}$ \\
			Inflation layers & $10$ & $16$ & $22$ \\
			Growth rate & $1.35$ & $1.22$ & $1.16$ \\
			$t_{\mathrm{infl}}$ [m] & $6.55 \times 10^{-4}$ & $1.59 \times 10^{-3}$ & $2.17 \times 10^{-3}$ \\
			Nodes & $14{,}525$ & $47{,}707$ & $106{,}460$ \\
			Elements & $13{,}984$ & $46{,}640$ & $104{,}763$ \\
			\hline
			\multicolumn{4}{@{}l}{\textit{Grid metric}} \\
			Static pressure drop $\Delta p_s$ [Pa] & $814.109$ & $796.804$ & $792.995$ \\
			$\epsilon_{\mathrm{rel}}(\Delta p_s)$ [\%] & -- & $2.13$ & $0.48$ \\
			\hline
		\end{tabular*}
		\vspace{2pt}
		\begin{minipage}{0.98\textwidth}
			\scriptsize \textit{Note.} Quad-dom. = quad-dominant. For the $\epsilon_{\mathrm{rel}}$ row, the Medium column gives the Coarse--Medium change and the Fine column gives the Medium--Fine change.
		\end{minipage}
	\end{table}

\begin{figure}[htpb]
	\centering
	\includegraphics[width=\textwidth]{UBend-Ansys-Mesh-Independence}
	\caption{Ansys U-bend exit velocity profiles across $D=0.028$ m for the Coarse, Medium, and Fine grids.}
	\label{fig:app_A_ubend_mesh_indep_ansys}
\end{figure}

\FloatBarrier

\chapter{Topology-Optimization Methodology Support}
\label{app:topology_methodology_support}

This appendix gathers supporting details for Chapter~\ref{ch:topology_optimization_methodology}: interpolation curves, passive-region handling, filtering conventions, inlet turbulence conversion, and solver controls. These details are kept separate from the main methodology chapter so the central formulation remains readable.

\section{RAMP Brinkman Interpolation}
\label{sec:app_topology_ramp_heaviside}
For visualization, the Brinkman resistance in Eq.~\ref{eq:brinkman_interpolation} can be normalized as
\begin{equation}
	\hat{\alpha}(\bar{\gamma}) =
	\frac{\alpha(\bar{\gamma})-\alpha_{fluid}}
	{\alpha_{solid}-\alpha_{fluid}} .
	\label{eq:normalized_brinkman_resistance}
\end{equation}
When $\alpha_{fluid}=0$, this becomes $\hat{\alpha}=\alpha(\bar{\gamma})/\alpha_{solid}$. With the $q$ convention used in this thesis, small $q$ values make most intermediate densities comparatively fluid-like except very close to $\bar{\gamma}=0$, while larger $q$ values increase the resistance over a wider part of the intermediate-density range.

\begin{figure}[!ht]
	\centering
	\includegraphics[width=0.72\textwidth]{RAMP-Interpolation-Curve}
	\caption{Normalized RAMP Brinkman resistance from Eq.~\ref{eq:brinkman_interpolation} for $q=\{0.01,0.05,0.1,0.5,1.0\}$  \ChatGPTFigureRef.}
	\label{fig:ramp_interpolation_curve}
\end{figure}

\section{Active and Passive Density Regions}
\label{sec:app_topology_active_passive_regions}
Figure~\ref{fig:filter_domain_schematic} shows how the optimization code separates cells that MMA is allowed to change from cells that are prescribed by the benchmark geometry. The active cells are design variables. Passive-fluid extensions, for example inlet and outlet ducts, are held at $\bar{\gamma}=1$ but are still included in the flow solve. Passive-solid regions, when used, are held at $\bar{\gamma}=0$.

\begin{figure}[!ht]
	\centering
	\includegraphics[width=0.80\textwidth]{Filter-Domain-Schematic}
	\caption{Active design region with passive inlet and outlet extensions \ChatGPTFigureRef.}
	\label{fig:filter_domain_schematic}
\end{figure}

\section{DG0 Density Filtering and Active-Mask Normalization}
\label{sec:app_topology_filtering}
The DG0 density filter of Eq.~\ref{eq:dg_helmholtz_filter} is applied before the projection step. Near prescribed inlet and outlet extensions, a direct filter would blend the fixed-fluid cells into the neighboring active cells. The active-mask normalization avoids that leakage.

The implementation therefore filters the active contribution and divides by the filtered active mask:
\begin{equation}
	\tilde{\gamma}
	=
	\frac{\mathcal{H}(M\gamma)}
	{\mathcal{H}(M)} .
	\label{eq:active_mask_filter}
\end{equation}
Here $\mathcal{H}$ denotes the filter operator and $M$ is the active-cell indicator. After the Heaviside projection, the prescribed passive-cell bounds are restored. The sensitivity chain uses the same mapping in reverse, so MMA receives derivatives only for the active raw density variables.

\section{Heaviside Projection Curve}
\label{sec:app_topology_heaviside}
Figure~\ref{fig:heaviside_projection_curve} plots the Heaviside projection used in Eq.~\ref{eq:tanh_projection}. Larger values of $\beta$ make the projected field switch more abruptly between solid-like and fluid-like values.

\begin{figure}[!ht]
	\centering
	\includegraphics[width=0.68\textwidth]{Heaviside-Projection-Curve}
	\caption{Heaviside projection from Eq.~\ref{eq:tanh_projection} at $\eta=0.5$  \ChatGPTFigureRef.}
	\label{fig:heaviside_projection_curve}
\end{figure}

\section{Inlet Turbulence Conversion}
\label{sec:app_topology_inlet_turbulence}
When an inlet turbulence intensity $I$ and turbulence length scale $\ell$ are provided, the code first estimates an eddy-viscosity ratio using the standard $k$--$\varepsilon$ inlet relations,
\begin{equation}
	k=\frac{3}{2}(UI)^2,
	\qquad
	\epsilon=C_\mu^{3/4}\frac{k^{3/2}}{\ell},
	\qquad
	\frac{\nu_t}{\nu}
	=
	C_\mu^{1/4}\sqrt{\frac{3}{2}}\,
	I\,\frac{U\ell}{\nu}.
	\label{eq:inlet_turbulence_to_viscosity_ratio}
\end{equation}
If the Reynolds number and length-scale ratio are supplied instead, $U\ell/\nu$ is replaced by $Re\,(\ell/L)$.

The resulting ratio is converted to the SA inlet working variable by solving the SA viscosity relation
\begin{equation}
	\nu_t = \tilde{\nu} f_{v1},
	\qquad
	f_{v1}=\frac{\chi^3}{\chi^3+c_{v1}^3},
	\qquad
	\chi=\frac{\tilde{\nu}}{\nu}.
	\label{eq:sa_viscosity_ratio_conversion}
\end{equation}
This conversion is only an inlet initialization and boundary-condition device. The SA transport equation still determines the interior $\tilde{\nu}$ field during the forward solve.

\section{Forward-Solver Controls}
\label{sec:app_topology_forward_controls}
The optimization runs use IPCS and Picard updates to obtain the forward state before assembling the frozen-viscosity residual for the adjoint. Some benchmarks also activate SNES recovery schedules, convection continuation, or accepted-residual criteria when intermediate designs are difficult to solve. These controls are forward-solver devices. The frozen adjoint in Chapter~\ref{ch:topology_optimization_methodology} is formed from the final steady residual in Eq.~\ref{eq:penalized_momentum_weak}, not by differentiating every IPCS, Picard, or recovery substep.

\chapter{Topology-Optimization Benchmark Details}
\label{app:C_topology_optimization_benchmarks}
\label{app:C_laminar_turbulent_topology_comparison}

This appendix supports Chapter~\ref{ch:benchmark_results}. It combines the turbulent and laminar topology-optimization support material in one place and groups it by benchmark case. Within each benchmark, the ordering is: turbulent-TO FEniCS settings and diagnostics, frozen-sensitivity checks where applicable, turbulent-design Ansys validation-mesh information, laminar-TO FEniCS settings and diagnostics, and laminar-design Ansys mesh independence where a completed hierarchy is available. For traceability, configuration file names are mentioned once in the relevant setup text; the tables summarize only the most important current settings.

Unless stated otherwise, all mesh-independence relative differences in this appendix are reported as positive percentage magnitudes. For a mesh-dependent metric $\phi$, the coarse--medium difference is computed as $|\phi_{\mathrm{coarse}}-\phi_{\mathrm{medium}}|/|\phi_{\mathrm{coarse}}|\times100\%$, while the medium--fine difference is computed as $|\phi_{\mathrm{medium}}-\phi_{\mathrm{fine}}|/|\phi_{\mathrm{fine}}|\times100\%$.

\section{Validation and Meshing Strategy}
\label{sec:app_B_mesh_strategy}
The FEniCS meshes used during topology optimization are treated as fixed optimization discretizations rather than as a separate mesh-convergence hierarchy. This distinction is important because refining the FEniCS mesh would also change the $DG_0$ design space, the filter stencil, the number of MMA design variables, and the diffuse Brinkman interface location. The pipe-bend optimization uses a rectangular-coordinate triangular mesh with near-wall clustering based on a smooth-wall $y^+\approx1$ estimate at the prescribed inlet Reynolds number. The U-bend optimization uses the corresponding wall-resolved triangular Gmsh mesh, with a near-wall size field based on the same $y^+\approx1$ estimate and a far-field target size of $h_{\max}=0.007L$ used to set the density-filter scale. The diffuser benchmark from Yoon~\cite{Yoon2016} uses a wall-resolved rectangular-coordinate triangular mesh over the square diffuser domain, refined around the full-height inlet, the center-third outlet, and the no-slip walls. In all cases, the mesh is chosen to provide sufficient wall-resolved SA and design resolution for the optimization problem; quantitative grid-convergence evidence is instead reported for the extracted sharp geometries in Ansys Fluent.

For each extracted turbulent topology optimization design with a completed Ansys mesh hierarchy, the validation mesh is refined through coarse, medium, and fine resolutions. Both SA and SST $k-\omega$ validation runs use the same extracted geometry, inlet condition, outlet condition, and fluid properties. The mesh-independence comparison is reported for both turbulence models where available because near-wall behavior and separated-flow predictions may respond differently to grid refinement. For the U-bend benchmark, the turbulent- and laminar-design validation tables use the same coarse--medium--fine mesh-sizing settings, while the final node and element counts differ because the extracted sharp-wall geometries are different.

Completed Ansys Fluent mesh-independence runs in this appendix use the same convergence and acceptance criteria. The static pressure drop must change by less than $5\times10^{-4}$ over the final $100$ iterations, while the viscous dissipation must change by less than $10^{-3}$ over the final $100$ iterations. The fine mesh is accepted only when the medium--fine relative difference is below $1\%$ for static pressure drop and below $2\%$ for viscous dissipation for both SA and SST $k-\omega$. The reported refinement differences are computed as $|\phi_{\mathrm{coarse}}-\phi_{\mathrm{medium}}|/|\phi_{\mathrm{coarse}}|$ and $|\phi_{\mathrm{medium}}-\phi_{\mathrm{fine}}|/|\phi_{\mathrm{fine}}|$ for each validation metric.

\section{Shared Turbulence-Topology Coupling Conventions}
\label{sec:app_B_shared_coupling_conventions}
All three turbulent topology optimization benchmarks use the same basic coupling idea between the density field, the SA working variable, and the reciprocal wall-distance equation. Topology-created solids act as walls for the SA working-variable penalty and the reciprocal wall-distance solve. The wall-density source is therefore \texttt{design}, the wall-distance mode is \texttt{reciprocal\_penalized}, and both the SA and wall-distance penalties use the power-law solid indicator defined in Chapter~\ref{ch:topology_optimization_methodology}. The same frozen-turbulence adjoint convention is also used: the forward loop updates $\tilde{\nu}$, $\nu_t$, and $G$, while the derivative freezes $\tilde{\nu}$ and $G$ during gradient assembly.

The pipe-bend and U-bend benchmarks then share a Bayat-style pressure-loss setup, while the diffuser follows the Yoon dissipation benchmark. The common mechanisms and the benchmark-specific parameter values are summarized in Table~\ref{tab:app_B_shared_coupling_conventions}.

The pipe-bend finite-difference sensitivity verification in Section~\ref{subsec:alexandersen_pipebend_sensitivity} uses the first pressure-objective stage of the current smooth-continuation setup and stores its own finite-difference output tables. It should therefore be read as a frozen-adjoint implementation check, not as evidence of convergence for the reported optimization state.

\begin{table}[htbp]
	\centering
	\footnotesize
	\caption{Shared turbulence-topology coupling conventions and benchmark-specific parameter values.}
	\label{tab:app_B_shared_coupling_conventions}
	\renewcommand{\arraystretch}{1.2}
	\begin{tabularx}{\textwidth}{l X X}
		\hline
		\textbf{Parameter} & \textbf{Pipe bend and U-bend} & \textbf{Diffuser} \\
		\hline
		SA wall-distance mode & \texttt{reciprocal\_penalized} & \texttt{reciprocal\_penalized} \\
		SA wall-density source & \texttt{design} & \texttt{design} \\
		SA penalty interpolation & Power-law solid indicator & Power-law solid indicator \\
		Wall-distance penalty interpolation & Power-law solid indicator & Power-law solid indicator \\
		Wall-distance recovery & Reciprocal recovery from the relaxed Eikonal variable $G$ & Reciprocal recovery from the relaxed Eikonal variable $G$ \\
		Adjoint treatment & Frozen-turbulence adjoint; $\tilde{\nu}$ and $G$ frozen during gradient assembly & Frozen-turbulence adjoint; $\tilde{\nu}$ and $G$ frozen during gradient assembly \\
		Solid threshold for wall source & $\bar{\gamma}=0.50$ & No separate threshold parameter in the diffuser configuration \\
			SA working-variable penalty & Pipe bend: fixed amplitude $10^3$, exponent $n=4$; U-bend: amplitude continuation ending at $10^5$, exponent $n=3$ & Power-law solid indicator with amplitude $10^5$ and exponent $n=3$ \\
			Wall-distance penalty & Pipe bend: fixed amplitude $10^3$, exponent $n=4$; U-bend: amplitude continuation ending at $10^5$, exponent $n=3$ & Power-law solid indicator with amplitude $10^5$ and exponent $n=3$ \\
			Wall-distance parameters & Pipe bend: $\sigma_w=0.01$, $G_0=20$; U-bend: $\sigma_w=0.1$, $G_0=20$ & $\sigma_w=0.1$, $G_0=20$ \\
		SA inlet treatment & Turbulence-intensity estimate: $I=7.5\%$, $\ell/H=0.05$, eddy-viscosity ratio capped at $50$ & Prescribed $\tilde{\nu}=1$ at inlet and outlet \\
		Optimization objective & Average inlet pressure with fixed outlet pressure $p=0$ & Dissipation objective $J_D$; the velocity-outlet baseline prescribes inlet and outlet profiles, while the primary pressure-outlet case prescribes the inlet profile and fixes outlet static pressure \\
		\hline
	\end{tabularx}
\end{table}

\FloatBarrier

\section{Pipe Bend Benchmark}
\label{sec:app_B_config_alexandersen_pipebend}

The pipe-bend case follows the setup of Bayat et al.~\cite{Bayat2026}. The reference paper uses a standard $k$-$\varepsilon$ model with wall functions; the present thesis uses the wall-resolved SA implementation and evaluates the $V_f=0.25$ design. The active turbulent FEniCS configuration is the pressure-objective smooth-continuation path in \texttt{Config\_PipeBendAlexandersen\_Frozen.py}, resumed from the reconstructed iteration-$70$ checkpoint that completed the first $q=1/150$, $\beta=4$ stage. This reported state is therefore still treated as an in-progress continuation state rather than as a converged final topology-optimization result.

\begin{table}[htbp]
	\centering
	\footnotesize
	\caption{Geometric, physical, and boundary-condition parameters for the pipe-bend benchmark.}
	\label{tab:app_B_alexandersen_pipebend_physical}
	\renewcommand{\arraystretch}{1.2}
	\begin{tabularx}{\textwidth}{l X}
		\hline
		\textbf{Parameter} & \textbf{Setting} \\
		\hline
		Reference case & Pipe-bend benchmark from Bayat et al.~\cite{Bayat2026} \\
		Reference figure for setup & Bayat et al. Fig.~10 \\
		Reference velocity figure & Bayat et al. Fig.~15 \\
		Reference turbulence model & Standard $k$-$\varepsilon$ with wall functions \\
		Turbulence model in this thesis & Spalart-Allmaras \\
		Design-domain dimensions & Square design region $L\times L$ with $L=1$ \\
		Full computational domain & $x/L\in[-0.2,1]$, $y/L\in[-0.2,1]$ \\
		Non-design inlet duct & Length $0.2L$, height $0.2L$, inlet over $y/L\in[0.7,0.9]$ \\
		Non-design outlet duct & Width $0.2L$, length $0.2L$, outlet over $x/L\in[0.7,0.9]$ on the lower boundary \\
		Volume fraction & $V_f=0.25$ \\
		Reynolds number & $Re=5{,}000$, based on $U_{\mathrm{in}}H_{\mathrm{in}}\rho/\mu$ with $H_{\mathrm{in}}=0.2L$ \\
		Density and viscosity & $\rho=1.0$, $\mu=4.0\times10^{-5}$ \\
		Inlet condition & Uniform inlet velocity $\mathbf{u}=(1,0)$; SA inlet estimated from the shared turbulence-intensity settings in Table~\ref{tab:app_B_shared_coupling_conventions} \\
		Outlet condition & Pressure outlet $p=0$ on the lower outlet, with the normal/tangential component constraint used by the FEniCS configuration \\
		Wall conditions & No-slip velocity and $\tilde{\nu}=0$ on external walls; topology-created solids enter through the design-density penalties \\
		Validation metrics & Average inlet pressure, static pressure drop, $\Phi_D$, and relative changes between FEniCS-SA, Ansys-SA, and Ansys SST $k-\omega$ \\
		\hline
	\end{tabularx}
\end{table}

\begin{table}[htbp]
	\centering
	\scriptsize
	\caption{Numerical FEniCS solver, topology-optimization, and sensitivity-verification parameters for the reported pipe-bend state.}
	\label{tab:app_B_alexandersen_pipebend_numerical}
	\renewcommand{\arraystretch}{1.2}
	\begin{tabularx}{\textwidth}{l c X}
		\hline
		\textbf{Parameter} & \textbf{Symbol / option} & \textbf{Setting} \\
		\hline
		Flow finite elements & $(\mathbf{u},p)$ & Taylor--Hood $P_2/P_1$ continuous Galerkin elements \\
		SA working-variable space & $\tilde{\nu}$ & $P_1$ continuous Galerkin \\
		Design-density space & $\gamma$ & $DG_0$ cellwise constant density \\
		FEniCS optimization mesh & -- & \texttt{mesh\_yplus1.xdmf}; $39{,}128$ vertices, $77{,}334$ triangles, $h_{\max}=0.007L$ \\
		Objective & $J_p$ & Average inlet pressure \\
		Initial design density & $\gamma_0$ & $V_f=0.25$, matched to the filtered volume \\
			Brinkman penalty & $\alpha(\bar{\gamma})$ & RAMP form with $\alpha_{\mathrm{fluid}}=0$, $\alpha_{\mathrm{solid}}=10^3$ \\
				RAMP continuation & $q$ & Pressure-objective schedule $\{1/150,1/120,1/100,1/75,1/50,1/30,1/15\}$ \\
				Projection continuation & $\beta$ & Schedule $\{4,4.5,5,6,7,8,9\}$ with threshold $\eta=0.50$ \\
				MMA move limits & -- & Schedule $\{0.010,0.004,0.0035,0.003,0.0025,0.002,0.0015\}$ \\
				Optimization iterations & -- & Resume from reconstructed iteration-$70$ checkpoint; seven-stage budgets $\{70,80,90,110,130,160,200\}$ with stopping after $14$ iterations below objective-change tolerance $10^{-5}$ \\
			Filter radius & $r_f$ & $4h_{\max}=0.028L$ \\
				SA and wall-distance penalties & -- & Fixed amplitudes $10^3$ with power-law exponent $n=4$; reciprocal wall-distance parameter $\sigma_w=0.01$; no penalty continuation in the active pressure-objective path \\
			Forward flow solver & -- & IPCS flow solve for the frozen-SA Picard update, followed by a warm-started SNES Newton line-search flow solve \\
			Frozen NS-SA coupling & -- & One Picard update per design iteration with SA relaxation factor $0.15$ \\
			Linear solver & -- & MUMPS for SNES/state and adjoint solves; BiCGSTAB with ILU preconditioning for IPCS warm-start solves \\
		Adjoint treatment & -- & Frozen-turbulence adjoint in $(\mathbf{u}^*,p^*)$; $\tilde{\nu}$ and $G$ are frozen during gradient assembly \\
		Sensitivity verification samples & -- & Five active DG0 cells, seed $13$ \\
		Sensitivity verification step & $\Delta h$ & $10^{-6}$, central finite difference \\
		Sensitivity verification fields held fixed & -- & SA working variable $\tilde{\nu}^{0}$ and reciprocal wall-distance field $G^{0}$ \\
		\hline
	\end{tabularx}
\end{table}

\subsection{Reference Topology Overlay}
\label{subsec:app_C_alexandersen_pipebend_overlay}
Figure~\ref{fig:app_C_alexandersen_pipebend_paper_overlay} overlays the present FEniCS-SA pipe-bend topology with the corresponding Bayat et al.~\cite{Bayat2026} wall-function result. The comparison is qualitative, because the reference and present designs use different turbulence closures and near-wall treatments, but the main routed bend is nearly coincident.

\begin{figure}[htbp]
	\centering
	\includegraphics[width=0.5\textwidth]{PipeBend-Turbulent-Paper-vs-FEniCS}
	\caption{Qualitative topology overlay for the turbulent pipe-bend benchmark; red: present FEniCS-SA fluid path, yellow: Bayat et al.~\cite{Bayat2026} wall-function topology.}
	\label{fig:app_C_alexandersen_pipebend_paper_overlay}
\end{figure}

\FloatBarrier

\subsection{Frozen Sensitivity Verification}
\label{subsec:alexandersen_pipebend_sensitivity}
		The frozen adjoint is checked directly on the pressure-objective pipe-bend setup inherited by \texttt{Config\_PipeBendAlexandersen\_Frozen\_Sensitivity.py}. At the checked initial design state this corresponds to the first smooth-continuation stage in Table~\ref{tab:app_B_alexandersen_pipebend_numerical}: $\alpha_{\mathrm{solid}}=10^3$, $q=1/150$, $\beta=4$, move limit $0.010$, one frozen-SA Picard update with SA relaxation $0.15$, wall-distance parameter $\sigma_w=0.01$, and power-law SA/wall-distance penalty exponent $n=4$. No full turbulence-adjoint or wall-distance-adjoint comparison is made; the purpose is to compare the frozen derivative supplied to MMA with objective changes obtained by re-solving the perturbed velocity-pressure state while using the same derivative path as the current continuation configuration.

	The checked-in coordinate finite-difference result is performed at the initial pipe-bend design state, with base objective
	$J_0=1.8343209502$. In every perturbation, the SA working variable $\tilde{\nu}^{0}$ and reciprocal wall-distance field $G^{0}$ are held fixed. This matches the frozen-turbulence assumption used by the adjoint: the flow state is re-solved, but the turbulence and wall-distance updates are not differentiated.

	The sensitivity configuration evaluates five active $DG_0$ samples with seed $13$. For sampled cell $i$, the raw density vector $\boldsymbol{\gamma}$ is perturbed by $\pm\Delta h\,\mathbf{e}_i$, where $\mathbf{e}_i$ is the unit vector for that active design variable and $\Delta h=10^{-6}$. The central finite-difference derivative is
\begin{equation}
	S_i^{\mathrm{FD}}
	=
	\frac{
			J_p\!\left(\boldsymbol{\gamma}+\Delta h\,\mathbf{e}_i;
			\mathbf{u}^{+},p^{+},\tilde{\nu}^{0},G^{0}\right)
		-
			J_p\!\left(\boldsymbol{\gamma}-\Delta h\,\mathbf{e}_i;
			\mathbf{u}^{-},p^{-},\tilde{\nu}^{0},G^{0}\right)
	}{
		2\Delta h
	}.
	\label{eq:alexandersen_pipebend_fd_metric}
\end{equation}
The frozen-adjoint value for the same cell is denoted by $S_i^{\mathrm{fr}}$; the superscripts $\mathrm{FD}$ and $\mathrm{fr}$ denote finite-difference and frozen-adjoint quantities, respectively. The absolute and relative discrepancies are computed as
\begin{equation}
	\Delta_i^{\mathrm{abs}} =
	\left|S_i^{\mathrm{fr}}-S_i^{\mathrm{FD}}\right|,
	\qquad
	\varepsilon_i^{\mathrm{fr}}
	=
	\frac{\Delta_i^{\mathrm{abs}}}
	{\max\!\left(\left|S_i^{\mathrm{fr}}\right|,\left|S_i^{\mathrm{FD}}\right|,10^{-30}\right)} .
	\label{eq:alexandersen_pipebend_fd_relative_error}
\end{equation}

\begin{table}[!ht]
	\centering
	\scriptsize
	\caption{Coordinate sensitivity check for the frozen pipe-bend configuration. Values are taken from \texttt{SensitivityVerificationTable.tsv}.}
	\label{tab:alexandersen_pipebend_fd_adjoint_comparison}
	\renewcommand{\arraystretch}{1.2}
	\begin{tabular}{l c c c c}
		\hline
		\textbf{Cell} &
		\shortstack{\textbf{Finite}\\\textbf{difference}\\$S_i^{\mathrm{FD}}$} &
		\shortstack{\textbf{Frozen}\\\textbf{adjoint}\\$S_i^{\mathrm{fr}}$} &
		\shortstack{\textbf{Absolute}\\\textbf{error}\\$\Delta_i^{\mathrm{abs}}$} &
		\shortstack{\textbf{Relative}\\\textbf{error}\\$\varepsilon_i^{\mathrm{fr}}$} \\
		\hline
			$CV_1$ & $-2.7011\times10^{-5}$ & $-2.7010\times10^{-5}$ & $1.5875\times10^{-10}$ & $5.8774\times10^{-6}$ \\
			$CV_2$ & $-2.1146\times10^{-6}$ & $-2.1148\times10^{-6}$ & $1.2513\times10^{-10}$ & $5.9171\times10^{-5}$ \\
			$CV_3$ & $-1.6664\times10^{-4}$ & $-1.6664\times10^{-4}$ & $4.2394\times10^{-10}$ & $2.5440\times10^{-6}$ \\
			$CV_4$ & $-1.7209\times10^{-5}$ & $-1.7209\times10^{-5}$ & $4.3775\times10^{-10}$ & $2.5437\times10^{-5}$ \\
			$CV_5$ & $ 2.3641\times10^{-5}$ & $ 2.3641\times10^{-5}$ & $2.0497\times10^{-10}$ & $8.6700\times10^{-6}$ \\
		\hline
	\end{tabular}
\end{table}

	The coordinate checks show close agreement between the frozen-adjoint derivative and the re-solved central finite differences across all five samples. The maximum absolute discrepancy is $4.38\times10^{-10}$, and the maximum relative discrepancy is $5.92\times10^{-5}$. The sign of the derivative also agrees in all five samples.

	The same sampled active cells are also used for a directional Taylor check with steps $(10^{-3},3\times10^{-4},10^{-4},3\times10^{-5})$ and seed $29$. The directional derivative is
	$9.6668564775\times10^{-5}$, and the central finite differences remain within $8.66\times10^{-8}$ relative error over the tested step range. The symmetric Taylor remainder decays with approximately second-order behavior, as reported in Table~\ref{tab:alexandersen_pipebend_taylor_check}.

\begin{table}[!ht]
	\centering
	\scriptsize
	\caption{Directional Taylor check for the frozen pipe-bend sensitivity configuration. Values are taken from \texttt{TaylorCheckLog.tsv}.}
	\label{tab:alexandersen_pipebend_taylor_check}
	\renewcommand{\arraystretch}{1.2}
	\begin{tabular}{c c c c c}
		\hline
		\textbf{Step $\epsilon$} &
		\shortstack{\textbf{Frozen}\\\textbf{directional}\\\textbf{derivative}} &
		\shortstack{\textbf{Central}\\\textbf{finite}\\\textbf{difference}} &
		\shortstack{\textbf{Directional}\\\textbf{relative}\\\textbf{error}} &
		\shortstack{\textbf{Symmetric}\\\textbf{Taylor}\\\textbf{order}} \\
		\hline
		$1.0\times10^{-3}$ & $9.6668564775\times10^{-5}$ & $9.6668565153\times10^{-5}$ & $3.9055\times10^{-9}$ & -- \\
		$3.0\times10^{-4}$ & $9.6668564775\times10^{-5}$ & $9.6668563524\times10^{-5}$ & $1.2939\times10^{-8}$ & $2.0006$ \\
		$1.0\times10^{-4}$ & $9.6668564775\times10^{-5}$ & $9.6668562044\times10^{-5}$ & $2.8252\times10^{-8}$ & $2.0077$ \\
		$3.0\times10^{-5}$ & $9.6668564775\times10^{-5}$ & $9.6668573146\times10^{-5}$ & $8.6596\times10^{-8}$ & $2.1339$ \\
		\hline
	\end{tabular}
\end{table}

		This check verifies the implemented frozen-turbulence derivative for the first-stage pipe-bend sensitivity setup that generated the stored TSV file. The conclusion is deliberately limited to that derivative: the SA working variable and reciprocal wall-distance field are held fixed during the checks, matching the optimization adjoint; consequently, the result should not be interpreted as validation of a fully coupled turbulence sensitivity or as convergence evidence for the resumed pipe-bend continuation state.

\FloatBarrier

\subsection{Turbulent FEniCS Optimization Diagnostics}
\label{subsec:app_B_alexandersen_pipebend_turb_diagnostics}
Figure~\ref{fig:alexandersen_pipebend_histories} records the convergence history for the first turbulent pipe-bend pressure-objective stage. The reported endpoint is the reconstructed iteration-$70$ checkpoint that seeds the current smooth-continuation configuration. The pressure objective is normalized by its initial value $J_{p,0}$, while the volume constraint is defined as the relative residual
\begin{equation}
	r_V^{(k)} = 100 \frac{v_f^{(k)} - V_f}{V_f},
	\label{eq:alexandersen_pipebend_volume_residual}
\end{equation}
where $v_f^{(k)}$ is the realized fluid volume fraction at MMA iteration $k$. With this sign convention, the dashed line at $r_V=0$ represents the active volume constraint. Values below this line are feasible and correspond to a design using slightly less fluid volume than the prescribed limit.

\begin{figure}[htbp]
	\centering
	\includegraphics[width=0.92\textwidth]{PipeBend-Turbulent-Jp-Convergence-History}
	\caption[Convergence history for the first turbulent pipe-bend stage]{Convergence history for the first $70$-iteration turbulent pipe-bend $J_p$ stage, showing $J_p/J_{p,0}$ and the relative volume-fraction residual $r_V$ from Eq.~\eqref{eq:alexandersen_pipebend_volume_residual}.}
	\label{fig:alexandersen_pipebend_histories}
\end{figure}

The monotonic objective reduction supports using the checkpoint scalar values for the provisional turbulent pipe-bend discussion in Chapter~\ref{ch:benchmark_results}. Since the current checked-in setup resumes from this stage into additional continuation, no final convergence claim is made from the first-stage history alone.

\FloatBarrier

The extracted turbulent pipe-bend geometry is meshed in Ansys Fluent with a patch-conforming triangular bulk mesh and wall-resolved inflation on all no-slip walls. The first-layer height is kept fixed from the smooth-wall $y^+\approx1$ estimate, while the mesh hierarchy refines the bulk face size, curvature angle, and inflation distribution.

\begin{table}[htbp]
	\centering
	\footnotesize
	\caption{Ansys Fluent mesh-independence table for the extracted turbulent pipe-bend benchmark geometry.}
	\label{tab:app_B_alexandersen_pipebend_mesh}
	\renewcommand{\arraystretch}{1.2}
	\setlength{\tabcolsep}{5pt}
	\begin{tabular*}{\textwidth}{@{\extracolsep{\fill}}lccc@{}}
		\hline
		\textbf{Parameter} & \textbf{Coarse} & \textbf{Medium} & \textbf{Fine} \\
		\hline
		\multicolumn{4}{@{}l}{\textit{Mesh settings}} \\
		Mesh type & Tri. & Tri. & Tri. \\
		$h_{\mathrm{bulk}}$ [m] & $1.10 \times 10^{-2}$ & $5.00 \times 10^{-3}$ & $3.25 \times 10^{-3}$ \\
		Curvature angle & $15^\circ$ & $7^\circ$ & $4.5^\circ$ \\
		Near-wall & WR infl. & WR infl. & WR infl. \\
		$\Delta y_1$ [m] & $6.00 \times 10^{-4}$ & $6.00 \times 10^{-4}$ & $6.00 \times 10^{-4}$ \\
		Inflation layers & $8$ & $14$ & $18$ \\
		Growth rate & $1.15$ & $1.10$ & $1.08$ \\
		$t_{\mathrm{infl}}$ [m] & $8.24 \times 10^{-3}$ & $1.68 \times 10^{-2}$ & $2.25 \times 10^{-2}$ \\
		Nodes & $5{,}147$ & $21{,}343$ & $44{,}042$ \\
		Elements & $7{,}551$ & $32{,}780$ & $68{,}743$ \\
		\hline
		\multicolumn{4}{@{}l}{\textit{Validation metrics}} \\
		$\Phi_D^{\mathrm{SA}}$ [W/m] & $0.0143$ & $0.0151$ & $0.0153$ \\
		$\epsilon_{\mathrm{rel}}(\Phi_D^{\mathrm{SA}})$ [\%] & -- & $5.38$ & $1.32$ \\
		$\Phi_D^{\mathrm{SST}\,k\text{-}\omega}$ [W/m] & $0.0150$ & $0.0164$ & $0.0166$ \\
		$\epsilon_{\mathrm{rel}}(\Phi_D^{\mathrm{SST}\,k\text{-}\omega})$ [\%] & -- & $9.04$ & $1.22$ \\
		Static pressure drop $\Delta p_s^{\mathrm{SA}}$ [Pa] & $0.1423$ & $0.1449$ & $0.1457$ \\
		$\epsilon_{\mathrm{rel}}(\Delta p_s^{\mathrm{SA}})$ [\%] & -- & $1.84$ & $0.55$ \\
		Static pressure drop $\Delta p_s^{\mathrm{SST}\,k\text{-}\omega}$ [Pa] & $0.1481$ & $0.1516$ & $0.1522$ \\
		$\epsilon_{\mathrm{rel}}(\Delta p_s^{\mathrm{SST}\,k\text{-}\omega})$ [\%] & -- & $2.36$ & $0.40$ \\
		Selected validation mesh & \multicolumn{3}{c}{Fine\_Ansys} \\
		\hline
	\end{tabular*}
	\vspace{2pt}
	\begin{minipage}{0.98\textwidth}
		\scriptsize \textit{Note.} WR = wall-resolved; $\Phi_D$ = viscous dissipation; $\Delta p_s$ = static pressure drop. For $\epsilon_{\mathrm{rel}}$ rows, the Medium column gives the Coarse--Medium change and the Fine column gives the Medium--Fine change.
	\end{minipage}
\end{table}

\FloatBarrier

\subsection{Laminar FEniCS Configuration Differences}
\label{subsec:app_C_alexandersen_pipebend_lam_config}
	The laminar pipe-bend reference is a completed FEniCS topology-optimization run from \texttt{Config\_PipeBendAlexandersen\_Laminar.py}. It reuses the turbulent pipe-bend setup in Tables~\ref{tab:app_B_alexandersen_pipebend_physical} and~\ref{tab:app_B_alexandersen_pipebend_numerical} wherever not stated otherwise: the geometry, wall-resolved mesh, non-design inlet and outlet ducts, uniform inlet velocity, pressure-outlet treatment, $V_f=0.25$ volume constraint, average inlet-pressure objective $J_p$, filter radius, and projection threshold are unchanged. The Brinkman resistance and continuation controls are laminar-specific and are listed in Table~\ref{tab:app_C_alexandersen_pipebend_lam_differences}.

\begin{table}[htbp]
	\centering
	\footnotesize
	\caption{Laminar-specific differences from the turbulent pipe-bend FEniCS setup above.}
	\label{tab:app_C_alexandersen_pipebend_lam_differences}
	\renewcommand{\arraystretch}{1.2}
	\begin{tabularx}{\textwidth}{l X}
		\hline
		\textbf{Parameter} & \textbf{Laminar setting} \\
		\hline
			Flow model & Laminar incompressible Navier--Stokes; the SA working variable, reciprocal wall-distance equation, topology-dependent SA penalties, and frozen-turbulence coupling are omitted \\
			Reynolds number & $Re=1$, based on $U_{\mathrm{in}}H_{\mathrm{in}}\rho/\mu$ instead of the turbulent $Re=5{,}000$ setup \\
			Density and viscosity & $\rho=1.0$ and $\mu=0.2$, with the same $U_{\mathrm{in}}=1$ and $H_{\mathrm{in}}=0.2L$ as the turbulent case \\
			Brinkman penalty & RAMP form with $\alpha_{\mathrm{fluid}}=0$ and $\alpha_{\mathrm{solid}}=10^5$ \\
			Flow finite elements & Taylor--Hood $P_2/P_1$ for $(\mathbf{u},p)$; no separate $P_1$ SA working-variable space is used \\
			Forward and adjoint solves & SNES/MUMPS laminar Newton forward solves with relative tolerance $10^{-6}$, absolute tolerance $10^{-8}$, and maximum forward iteration count $220$; the adjoint is solved as a linear system with MUMPS \\
			Continuation and optimizer controls & Seven-stage schedule $q=\{0.01,0.02,0.04,0.08,0.08,0.12,0.20\}$, $\beta=\{1,2,4,8,8,12,16\}$, move limits $\{0.010,0.008,0.006,0.004,0.003,0.002,0.0015\}$, and inner-iteration budgets $\{12,25,45,80,120,160,220\}$; stop if objective change is below $10^{-6}$ for ten iterations \\
		\hline
	\end{tabularx}
\end{table}

\subsection{Laminar FEniCS Optimization Diagnostics}
\label{subsec:app_C_alexandersen_pipebend_lam_diagnostics}
Figure~\ref{fig:app_C_alexandersen_pipebend_lam_convergence} records the completed laminar pipe-bend optimization history. The normalized inlet-pressure objective drops rapidly in the first continuation stage and then remains nearly flat, while the volume-fraction residual repeatedly returns to the active constraint after continuation updates.

\begin{figure}[htbp]
	\centering
	\includegraphics[width=\textwidth]{PipeBend-Laminar-Jp-Convergence-History}
	\caption{FEniCS convergence history for the laminar pipe-bend topology-optimization run, showing $J_p/J_{p,0}$ and the relative volume-fraction residual.}
	\label{fig:app_C_alexandersen_pipebend_lam_convergence}
\end{figure}

\FloatBarrier

\subsection{Laminar-Design Mesh Independence}
\label{subsec:app_C_alexandersen_pipebend_lam_mesh}
Table~\ref{tab:app_C_alexandersen_pipebend_lam_mesh} reports the Ansys mesh-independence study for the extracted laminar pipe-bend design. This table is separate from the turbulent-design mesh study in Table~\ref{tab:app_B_alexandersen_pipebend_mesh} of Section~\ref{sec:app_B_config_alexandersen_pipebend} because the laminar and turbulent topologies generally produce different wall shapes and near-wall mesh requirements.

\begin{table}[htbp]
	\centering
	\footnotesize
	\caption{Ansys mesh-independence table for the extracted laminar pipe-bend design.}
	\label{tab:app_C_alexandersen_pipebend_lam_mesh}
	\renewcommand{\arraystretch}{1.2}
	\setlength{\tabcolsep}{5pt}
	\begin{tabular*}{\textwidth}{@{\extracolsep{\fill}}lccc@{}}
		\hline
		\textbf{Parameter} & \textbf{Coarse} & \textbf{Medium} & \textbf{Fine} \\
		\hline
		\multicolumn{4}{@{}l}{\textit{Mesh settings}} \\
		Mesh type & Tri. & Tri. & Tri. \\
		$h_{\mathrm{bulk}}$ [m] & $4.50 \times 10^{-3}$ & $2.60 \times 10^{-3}$ & $2.00 \times 10^{-3}$ \\
		Curvature angle & $8^\circ$ & $4^\circ$ & $3^\circ$ \\
		$\Delta y_1$ [m] & $1.20 \times 10^{-3}$ & $1.20 \times 10^{-3}$ & $1.20 \times 10^{-3}$ \\
		Inflation layers & $12$ & $15$ & $16$ \\
		Growth rate & $1.06$ & $1.045$ & $1.04$ \\
		$t_{\mathrm{infl}}$ [m] & $2.02 \times 10^{-2}$ & $2.49 \times 10^{-2}$ & $2.62 \times 10^{-2}$ \\
		Nodes & $22{,}494$ & $56{,}468$ & $89{,}804$ \\
		Elements & $35{,}948$ & $93{,}828$ & $153{,}163$ \\
		\hline
		\multicolumn{4}{@{}l}{\textit{Validation metrics}} \\
		$\Phi_D^{\mathrm{SST}\,k\text{-}\omega}$ [W/m] & $0.0293$ & $0.0310$ & $0.0313$ \\
		$\epsilon_{\mathrm{rel}}(\Phi_D^{\mathrm{SST}\,k\text{-}\omega})$ [\%] & -- & $5.82$ & $0.95$ \\
		Static pressure drop $\Delta p_s^{\mathrm{SST}\,k\text{-}\omega}$ [Pa] & $0.5229$ & $0.5420$ & $0.5444$ \\
		$\epsilon_{\mathrm{rel}}(\Delta p_s^{\mathrm{SST}\,k\text{-}\omega})$ [\%] & -- & $3.65$ & $0.45$ \\
		Selected comparison mesh & \multicolumn{3}{c}{Fine\_Ansys} \\
		\hline
	\end{tabular*}
	\vspace{2pt}
	\begin{minipage}{0.98\textwidth}
			\scriptsize \textit{Note.} $\Phi_D$ = viscous dissipation; $\Delta p_s$ = static pressure drop.
	\end{minipage}
\end{table}

\FloatBarrier

\section{U-Bend Benchmark}
\label{sec:app_B_config_alexandersen_ubend}

The U-bend case follows Fig.~11 and Table~4 of Bayat et al.~\cite{Bayat2026}. It uses the same turbulent topology optimization framework as the pipe bend, but forces a $180^\circ$ return path around a fixed internal separator. The active turbulent FEniCS setup is summarized from \texttt{Config\_UBendAlexandersen\_Frozen.py}.

\begin{table}[htbp]
	\centering
	\footnotesize
	\caption{Geometric, physical, and boundary-condition parameters for the U-bend benchmark.}
	\label{tab:app_B_alexandersen_ubend_physical}
	\renewcommand{\arraystretch}{1.2}
	\begin{tabularx}{\textwidth}{l X}
		\hline
		\textbf{Parameter} & \textbf{Setting} \\
		\hline
		Reference case & U-bend benchmark from Bayat et al.~\cite{Bayat2026} \\
		Reference figure/table for setup & Bayat et al. Fig.~11 and Table~4 \\
		Reference turbulence model & Standard $k$-$\varepsilon$ with wall functions \\
		Turbulence model in this thesis & Spalart-Allmaras \\
		Design-domain dimensions & Square design region $L\times L$ with $L=1$ \\
		Full computational domain & $x/L\in[-0.2,1]$, $y/L\in[0,1]$ \\
		Port geometry & Two left-side ports of height $0.2L$; inlet over $y/L\in[0.55,0.75]$, outlet over $y/L\in[0.25,0.45]$ \\
		Fixed separator & Bar thickness $0.10L$, rounded-tip radius $0.05L$, extending from $x/L=-0.2$ to tip at $x/L=0.5$ \\
		Volume fraction & $V_f=0.27$ \\
		Reynolds number & $Re=5{,}000$, based on $U_{\mathrm{in}}H_{\mathrm{port}}\rho/\mu$ with $H_{\mathrm{port}}=0.2L$ \\
		Density and viscosity & $\rho=1.0$, $\mu=4.0\times10^{-5}$ \\
		Inlet condition & Uniform inlet velocity $\mathbf{u}=(1,0)$; SA inlet estimated from the shared turbulence-intensity settings in Table~\ref{tab:app_B_shared_coupling_conventions} \\
		Outlet condition & Pressure outlet $p=0$ on the lower left port, with the component constraint used by the FEniCS configuration \\
		Wall conditions & No-slip velocity and $\tilde{\nu}=0$ on external walls and the fixed separator; topology-created solids enter through the design-density penalties \\
		Validation metrics & Average inlet pressure, static pressure drop, $\Phi_D$, and relative changes between FEniCS-SA, Ansys-SA, and Ansys SST $k-\omega$ \\
		\hline
	\end{tabularx}
\end{table}

\begin{table}[htbp]
	\centering
	\footnotesize
	\caption{Numerical FEniCS solver and topology-optimization parameters for the U-bend benchmark.}
	\label{tab:app_B_alexandersen_ubend_numerical}
	\renewcommand{\arraystretch}{1.2}
	\begin{tabularx}{\textwidth}{l c X}
		\hline
		\textbf{Parameter} & \textbf{Symbol / option} & \textbf{Setting} \\
		\hline
		Flow finite elements & $(\mathbf{u},p)$ & Taylor--Hood $P_2/P_1$ continuous Galerkin elements \\
		SA working-variable space & $\tilde{\nu}$ & $P_1$ continuous Galerkin \\
		Design-density space & $\gamma$ & $DG_0$ cellwise constant density \\
		Reference FEniCS mesh & -- & \texttt{mesh\_yplus1.xdmf}; wall-resolved triangular Gmsh mesh with $45{,}280$ vertices, $85{,}780$ triangles, near-wall refinement based on $y^+\approx1$, and far-field $h_{\max}=0.007L$ \\
		Objective & $J_p$ & Average inlet pressure \\
		Initial design density & $\gamma_0$ & $V_f=0.27$, matched to the filtered volume \\
			Brinkman penalty & $\alpha(\bar{\gamma})$ & RAMP form with $\alpha_{\mathrm{fluid}}=0$, $\alpha_{\mathrm{solid}}=10^3$ for the active $J_p$ path \\
			RAMP continuation & $q$ & $[0.02,\ 0.04,\ 0.08,\ 0.12,\ 0.18]$ \\
			Projection continuation & $\beta$ & $[1.5,\ 3,\ 5,\ 8,\ 12]$ with threshold $\eta=0.50$ \\
			MMA move limits & -- & $[0.012,\ 0.010,\ 0.007,\ 0.005,\ 0.003]$ \\
			Maximum optimization iterations & -- & Five continuation stages with $[60,\ 70,\ 90,\ 110,\ 140]$ MMA iterations; stop if objective change is below $10^{-5}$ for ten iterations \\
			Filter radius & $r_f$ & $4h_{\max}=0.028L$ \\
			SA and wall-distance penalty continuation & -- & Shared amplitude schedule $[10^3,\ 3\times10^3,\ 10^4,\ 3\times10^4,\ 10^5,\ 10^5,\ 10^5]$; the first five entries are used by the active five-stage continuation \\
			Forward flow solver & -- & IPCS flow solves for frozen-SA Picard updates, followed by a warm-started SNES Newton line-search final flow solve \\
				Frozen NS-SA coupling & -- & Three Picard updates per design iteration with SA relaxation factor $0.04$ for the active $J_p$ path \\
			Linear solver & -- & MUMPS for SNES/state and adjoint solves; BiCGSTAB with ILU preconditioning for IPCS warm-start solves \\
		Adjoint treatment & -- & Frozen-turbulence adjoint in $(\mathbf{u}^*,p^*)$; $\tilde{\nu}$ and $G$ are frozen during gradient assembly \\
		\hline
	\end{tabularx}
\end{table}

No separate U-bend sensitivity-verification table is reported. The U-bend uses the same average inlet-pressure objective $J_p$ and the same frozen-turbulence $(\mathbf{u},p)$ adjoint implementation as the pipe-bend pressure-loss benchmark, for which the direct finite-difference and Taylor checks are reported above. The omitted U-bend check would therefore mainly repeat the same derivative path, while adding a final-flow-solver robustness issue specific to the tighter U-bend continuation rather than providing independent adjoint evidence.

\subsection{Turbulent FEniCS Optimization Diagnostics}
\label{subsec:app_B_alexandersen_ubend_turb_diagnostics}
Figure~\ref{fig:alexandersen_ubend_histories} records the available convergence history for the turbulent U-bend topology through $275$ global MMA iterations. The pressure objective is normalized by its initial value $J_{p,0}$, and the volume trace uses the same relative residual definition as Eq.~\eqref{eq:alexandersen_pipebend_volume_residual}. The run should be interpreted as an optimization diagnostic rather than a converged final result: the normalized objective decreases smoothly over the available iterations, while the volume residual shows continuation-related excursions before returning toward the active constraint.

\begin{figure}[htbp]
	\centering
	\includegraphics[width=0.92\textwidth]{UBend-Turbulent-Jp-Convergence-History}
	\caption[Convergence history for the provisional turbulent U-bend run]{Convergence history for the $275$-iteration turbulent U-bend $J_p$ run, showing $J_p/J_{p,0}$ and the relative volume-fraction residual $r_V$.}
	\label{fig:alexandersen_ubend_histories}
\end{figure}

The next step for this benchmark is therefore a more robust turbulent-design U-bend optimization campaign, with solver and continuation settings tuned for the forced $180^\circ$ return flow and run to a defined convergence criterion.

\FloatBarrier

\subsection{Turbulent-Design Mesh Independence}
\label{subsec:app_B_alexandersen_ubend_turb_mesh}
Table~\ref{tab:app_B_alexandersen_ubend_turb_mesh} documents the Ansys validation mesh hierarchy for the extracted turbulent U-bend design. The coarse and fine mesh-sizing settings match those used for the laminar-design U-bend mesh-independence table in Table~\ref{tab:app_C_alexandersen_ubend_lam_mesh}, while the medium grid is refined for the turbulent-design geometry. The extracted sharp-wall geometry, final node and element counts, and resulting validation metrics therefore differ between the two studies.

\begin{table}[htbp]
	\centering
	\footnotesize
	\caption{Ansys mesh-independence table for the extracted turbulent U-bend design.}
	\label{tab:app_B_alexandersen_ubend_turb_mesh}
	\renewcommand{\arraystretch}{1.2}
	\setlength{\tabcolsep}{5pt}
	\begin{tabular*}{\textwidth}{@{\extracolsep{\fill}}lccc@{}}
		\hline
		\textbf{Parameter} & \textbf{Coarse} & \textbf{Medium} & \textbf{Fine} \\
		\hline
		\multicolumn{4}{@{}l}{\textit{Mesh settings}} \\
		Mesh type & Tri. & Tri. & Tri. \\
		$h_{\mathrm{bulk}}$ [m] & $6.00 \times 10^{-3}$ & $3.25 \times 10^{-3}$ & $2.75 \times 10^{-3}$ \\
		Curvature angle & $8^\circ$ & $6^\circ$ & $4^\circ$ \\
		$\Delta y_1$ [m] & $6.00 \times 10^{-4}$ & $6.00 \times 10^{-4}$ & $6.00 \times 10^{-4}$ \\
		Inflation layers & $12$ & $18$ & $20$ \\
		Growth rate & $1.10$ & $1.07$ & $1.06$ \\
		$t_{\mathrm{infl}}$ [m] & $1.28 \times 10^{-2}$ & $2.04 \times 10^{-2}$ & $2.21 \times 10^{-2}$ \\
		Nodes & $15{,}250$ & $44{,}627$ & $59{,}252$ \\
		Elements & $22{,}878$ & $68{,}769$ & $91{,}806$ \\
		\hline
		\multicolumn{4}{@{}l}{\textit{Validation metrics}} \\
		$\Phi_D^{\mathrm{SA}}$ [W/m] & $0.0791$ & $0.0785$ & $0.0785$ \\
		$\epsilon_{\mathrm{rel}}(\Phi_D^{\mathrm{SA}})$ [\%] & -- & $0.76$ & $0.00$ \\
		$\Phi_D^{\mathrm{SST}\,k\text{-}\omega}$ [W/m] & $0.0961$ & $0.1009$ & $0.1016$ \\
		$\epsilon_{\mathrm{rel}}(\Phi_D^{\mathrm{SST}\,k\text{-}\omega})$ [\%] & -- & $4.99$ & $0.69$ \\
		Static pressure drop $\Delta p_s^{\mathrm{SA}}$ [Pa] & $0.8725$ & $0.8123$ & $0.8071$ \\
		$\epsilon_{\mathrm{rel}}(\Delta p_s^{\mathrm{SA}})$ [\%] & -- & $6.90$ & $0.64$ \\
		Static pressure drop $\Delta p_s^{\mathrm{SST}\,k\text{-}\omega}$ [Pa] & $1.2139$ & $1.1936$ & $1.1973$ \\
		$\epsilon_{\mathrm{rel}}(\Delta p_s^{\mathrm{SST}\,k\text{-}\omega})$ [\%] & -- & $1.67$ & $0.31$ \\
		Selected comparison mesh & \multicolumn{3}{c}{Fine\_Ansys} \\
		\hline
	\end{tabular*}
	\vspace{2pt}
	\begin{minipage}{0.98\textwidth}
		\scriptsize \textit{Note.} $\Phi_D$ = viscous dissipation; $\Delta p_s$ = static pressure drop. For $\epsilon_{\mathrm{rel}}$ rows, the Medium column gives the Coarse--Medium change and the Fine column gives the Medium--Fine change.
	\end{minipage}
\end{table}

\FloatBarrier

\subsection{Laminar FEniCS Configuration Differences}
\label{subsec:app_C_alexandersen_ubend_lam_config}
	The laminar U-bend reference is a completed FEniCS topology-optimization run from \texttt{Config\_UBendAlexandersen\_Laminar.py}. It reuses the turbulent U-bend setup in Tables~\ref{tab:app_B_alexandersen_ubend_physical} and~\ref{tab:app_B_alexandersen_ubend_numerical} wherever not stated otherwise: the geometry, fixed separator, wall-resolved mesh, non-design ports, uniform inlet velocity, pressure-outlet treatment, $V_f=0.27$ volume constraint, average inlet-pressure objective $J_p$, filter radius, and projection threshold are unchanged. The Brinkman resistance and continuation controls are laminar-specific and are listed in Table~\ref{tab:app_C_alexandersen_ubend_lam_differences}.

\begin{table}[htbp]
	\centering
	\footnotesize
	\caption{Laminar-specific differences from the turbulent U-bend FEniCS setup above.}
	\label{tab:app_C_alexandersen_ubend_lam_differences}
	\renewcommand{\arraystretch}{1.2}
	\begin{tabularx}{\textwidth}{l X}
		\hline
		\textbf{Parameter} & \textbf{Laminar setting} \\
		\hline
			Flow model & Laminar incompressible Navier--Stokes; the SA working variable, reciprocal wall-distance equation, topology-dependent SA penalties, and frozen-turbulence coupling are omitted \\
			Reynolds number & $Re=1$, based on $U_{\mathrm{in}}H_{\mathrm{port}}\rho/\mu$ instead of the turbulent $Re=5{,}000$ setup \\
			Density and viscosity & $\rho=1.0$ and $\mu=0.2$, with the same $U_{\mathrm{in}}=1$ and $H_{\mathrm{port}}=0.2L$ as the turbulent case \\
			Brinkman penalty & RAMP form with $\alpha_{\mathrm{fluid}}=0$ and $\alpha_{\mathrm{solid}}=10^5$ \\
			Flow finite elements & Taylor--Hood $P_2/P_1$ for $(\mathbf{u},p)$; no separate $P_1$ SA working-variable space is used \\
			Forward and adjoint solves & SNES/MUMPS laminar Newton forward solves with relative tolerance $10^{-6}$, absolute tolerance $10^{-8}$, and maximum forward iteration count $260$; the adjoint is solved as a linear system with MUMPS \\
			Continuation and optimizer controls & Seven-stage schedule $q=\{0.01,0.02,0.04,0.08,0.08,0.12,0.20\}$, $\beta=\{1,2,4,8,8,12,16\}$, move limits $\{0.010,0.008,0.006,0.004,0.003,0.002,0.0015\}$, and inner-iteration budgets $\{12,25,45,80,120,160,220\}$; stop if objective change is below $10^{-6}$ for ten iterations \\
		\hline
	\end{tabularx}
\end{table}

\subsection{Laminar FEniCS Optimization Diagnostics}
\label{subsec:app_C_alexandersen_ubend_lam_diagnostics}
Figure~\ref{fig:app_C_alexandersen_ubend_lam_convergence} records the completed laminar U-bend optimization history. As in the pipe-bend case, the objective reduction occurs mainly at the start of the run, after which the optimizer preserves a nearly active volume constraint while only small continuation-related corrections remain.

\begin{figure}[htbp]
	\centering
	\includegraphics[width=\textwidth]{UBend-Laminar-Jp-Convergence-History}
	\caption{FEniCS convergence history for the laminar U-bend topology-optimization run, showing $J_p/J_{p,0}$ and the relative volume-fraction residual.}
	\label{fig:app_C_alexandersen_ubend_lam_convergence}
\end{figure}

\FloatBarrier

\subsection{Laminar-Design Mesh Independence}
\label{subsec:app_C_alexandersen_ubend_lam_mesh}
Table~\ref{tab:app_C_alexandersen_ubend_lam_mesh} reports the Ansys mesh-independence study for the extracted laminar U-bend design. The turbulent-design validation table in Table~\ref{tab:app_B_alexandersen_ubend_turb_mesh} uses the same coarse and fine sizing settings, but a refined medium grid.
For the $Re=5{,}000$ U-bend operating point, the inlet-port height is $0.2L$ and $U_{\mathrm{in}}=1~\mathrm{m\,s^{-1}}$, giving $\nu=4.0\times10^{-5}~\mathrm{m^2\,s^{-1}}$. Using the same smooth-wall estimate as the FEniCS mesh generator gives $u_\tau\approx6.59\times10^{-2}~\mathrm{m\,s^{-1}}$ and a wall distance of $y_{+}=1$ at $6.07\times10^{-4}~\mathrm{m}$. The Ansys mesh hierarchy therefore keeps the first layer fixed at $6.00\times10^{-4}~\mathrm{m}$ and refines the bulk triangular mesh, curvature resolution, and inflation-layer distribution around the separator tip and $180^\circ$ return wall.

\begin{table}[htbp]
	\centering
	\footnotesize
	\caption{Ansys mesh-independence table for the extracted laminar U-bend design.}
	\label{tab:app_C_alexandersen_ubend_lam_mesh}
	\renewcommand{\arraystretch}{1.2}
	\setlength{\tabcolsep}{5pt}
	\begin{tabular*}{\textwidth}{@{\extracolsep{\fill}}lccc@{}}
		\hline
		\textbf{Parameter} & \textbf{Coarse} & \textbf{Medium} & \textbf{Fine} \\
		\hline
		\multicolumn{4}{@{}l}{\textit{Mesh settings}} \\
		Mesh type & Tri. & Tri. & Tri. \\
		$h_{\mathrm{bulk}}$ [m] & $6.00 \times 10^{-3}$ & $4.00 \times 10^{-3}$ & $2.75 \times 10^{-3}$ \\
		Curvature angle & $8^\circ$ & $6^\circ$ & $4^\circ$ \\
		$\Delta y_1$ [m] & $6.00 \times 10^{-4}$ & $6.00 \times 10^{-4}$ & $6.00 \times 10^{-4}$ \\
		Inflation layers & $12$ & $16$ & $20$ \\
		Growth rate & $1.10$ & $1.08$ & $1.06$ \\
		$t_{\mathrm{infl}}$ [m] & $1.28 \times 10^{-2}$ & $1.82 \times 10^{-2}$ & $2.21 \times 10^{-2}$ \\
		Nodes & $16{,}106$ & $32{,}580$ & $62{,}622$ \\
		Elements & $24{,}460$ & $50{,}067$ & $98{,}084$ \\
		\hline
		\multicolumn{4}{@{}l}{\textit{Validation metrics}} \\
		$\Phi_D^{\mathrm{SST}\,k\text{-}\omega}$ [W/m] & $0.101$ & $0.105$ & $0.107$ \\
		$\epsilon_{\mathrm{rel}}(\Phi_D^{\mathrm{SST}\,k\text{-}\omega})$ [\%] & -- & $3.96$ & $1.90$ \\
		Static pressure drop $\Delta p_s^{\mathrm{SST}\,k\text{-}\omega}$ [Pa] & $1.125$ & $1.141$ & $1.150$ \\
		$\epsilon_{\mathrm{rel}}(\Delta p_s^{\mathrm{SST}\,k\text{-}\omega})$ [\%] & -- & $1.42$ & $0.79$ \\
		Selected comparison mesh & \multicolumn{3}{c}{Fine\_Ansys} \\
		\hline
	\end{tabular*}
	\vspace{2pt}
	\begin{minipage}{0.98\textwidth}
		\scriptsize \textit{Note.} $\Phi_D$ = viscous dissipation; $\Delta p_s$ = static pressure drop. For $\epsilon_{\mathrm{rel}}$ rows, the Medium column gives the Coarse--Medium change and the Fine column gives the Medium--Fine change.
	\end{minipage}
\end{table}

\FloatBarrier

\subsection{Recirculation Diagnostic}
\label{subsec:app_C_alexandersen_ubend_recirc}
Figure~\ref{fig:app_C_alexandersen_ubend_recirc_sst} gives a close-up of the SST $k-\omega$ velocity field near the separator and inner return path for the extracted laminar and turbulent U-bend designs. The vectors show local recirculation in the low-momentum region downstream of the separator, supporting the qualitative separation-and-recovery interpretation used in Chapter~\ref{ch:benchmark_results}. The performance ranking is still taken from the integrated pressure-drop and dissipation metrics.

\begin{figure}[htbp]
	\centering
	\begin{subfigure}[b]{0.98\textwidth}
		\centering
		\includegraphics[width=\textwidth]{UBend-Laminar-Recirculation}
		\caption{Laminar-optimized design.}
		\label{fig:app_C_alexandersen_ubend_laminar_recirc_sst}
	\end{subfigure}
	\vspace{0.6em}
	\begin{subfigure}[b]{0.98\textwidth}
		\centering
		\includegraphics[width=\textwidth]{UBend-Turbulent-Recirculation}
		\caption{Turbulent-optimized design.}
		\label{fig:app_C_alexandersen_ubend_turbulent_recirc_sst}
	\end{subfigure}
	\caption{SST $k-\omega$ velocity-vector close-up near the U-bend separator and inner return path.}
	\label{fig:app_C_alexandersen_ubend_recirc_sst}
\end{figure}

\FloatBarrier

\section{Diffuser Benchmark}
\label{sec:app_B_config_yoon_diffuser}

The diffuser case follows Fig.~20--22 of Yoon~\cite{Yoon2016}. The original velocity-outlet formulation prescribes parabolic velocity profiles at both the inlet and the outlet and uses the dissipation objective $J_D$. Chapter~\ref{ch:benchmark_results} uses that case as a boundary-condition sensitivity baseline, then treats the pressure-outlet variant as the primary diffuser result. The active turbulent FEniCS setup is summarized from \texttt{Config\_DiffuserYoon\_Frozen.py}. Both variants use the shared \texttt{design}, \texttt{reciprocal\_penalized}, power-law, and frozen-adjoint conventions summarized in Table~\ref{tab:app_B_shared_coupling_conventions}.

\begin{table}[htbp]
	\centering
	\footnotesize
	\caption{Geometric, physical, and boundary-condition parameters for the diffuser benchmark.}
	\label{tab:app_B_yoon_diffuser_physical}
	\renewcommand{\arraystretch}{1.2}
	\begin{tabularx}{\textwidth}{l X}
		\hline
		\textbf{Parameter} & \textbf{Setting} \\
		\hline
		Reference case & Diffuser benchmark from Yoon~\cite{Yoon2016} \\
		Reference figures for setup and topology & Yoon Fig.~20--22 \\
		Reference turbulence model & Spalart-Allmaras \\
		Turbulence model in this thesis & Spalart-Allmaras with frozen-turbulence adjoint \\
		Design-domain dimensions & Square design region $L\times L$ with $L=1~\mathrm{m}$ \\
		Full computational domain & $x/L\in[0,1]$, $y/L\in[0,1]$ \\
			Inlet geometry & FEniCS inlet is the full left boundary, $x/L=0$ and $y/L\in[0,1]$ \\
			Outlet geometry & FEniCS pressure outlet is the center third of the right boundary, $x/L=1$ and $y/L\in[1/3,2/3]$; right boundary outside the outlet is wall \\
		Volume fraction & $V_f=0.30$ fluid, corresponding to Yoon's $70\%$ solid mass usage \\
		Reynolds number & $Re=3{,}000$, based on $U_{\max,\mathrm{in}}L\rho/\mu$ \\
		Density and viscosity & $\rho=1000~\mathrm{kg/m^3}$, $\mu=1.0~\mathrm{Pa\,s}$ \\
			Inlet condition & Parabolic velocity profile with $U_{\max,\mathrm{in}}=3~\mathrm{m/s}$; $\tilde{\nu}=1$ for the FEniCS-SA solve \\
			Outlet condition & Velocity-outlet baseline: parabolic outflow profile with maximum velocity set by incompressible flow-rate conservation and pressure pinned at $(0,0)$ for gauge fixing. Pressure-outlet primary case: fixed static pressure $p=0~\mathrm{Pa}$ on the center-third outlet with no prescribed outlet velocity. \\
		Wall conditions & No-slip velocity and $\tilde{\nu}=0$ on the top and bottom walls and on the right boundary outside the outlet \\
		Brinkman solid resistance & $\alpha_{\mathrm{solid}}=10^9$, $\alpha_{\mathrm{fluid}}=0$ \\
		Optimization objective & Dissipation objective $J_D$, including physical viscous dissipation and Brinkman drag over the design domain \\
		Validation metrics & $J_D$, static pressure drop, $\Phi_D$, and relative changes between FEniCS-SA, Ansys-SA, and Ansys SST $k-\omega$ \\
		\hline
	\end{tabularx}
\end{table}

For the body-fitted Ansys Fluent evaluations, the inlet parabolic profile is shifted to the extracted inlet segment. In the velocity-outlet baseline, the outlet profile is also shifted to the extracted outlet segment and its maximum velocity is adjusted by incompressible flow-rate conservation. In the pressure-outlet variant, the outlet velocity distribution is solved rather than prescribed.

\begin{table}[htbp]
	\centering
	\footnotesize
	\caption{Numerical FEniCS solver and topology-optimization parameters for the diffuser benchmark.}
	\label{tab:app_B_yoon_diffuser_numerical}
	\renewcommand{\arraystretch}{1.2}
	\begin{tabularx}{\textwidth}{l c X}
		\hline
		\textbf{Parameter} & \textbf{Symbol / option} & \textbf{Setting} \\
		\hline
		Flow finite elements & $(\mathbf{u},p)$ & Taylor--Hood $P_2/P_1$ continuous Galerkin elements \\
		SA working-variable space & $\tilde{\nu}$ & $P_1$ continuous Galerkin \\
		Design-density space & $\gamma$ & $DG_0$ cellwise constant density \\
		Reference FEniCS mesh & -- & \texttt{mesh\_yoon\_yplus1.xdmf}; wall-resolved triangular mesh with $38{,}745$ vertices, $76{,}704$ triangles, far-field $h_{\max}=L/180=5.56\times10^{-3}L$, and first-cell width $1.82\times10^{-3}L$ \\
		Objective & $J_D$ & Dissipation objective over the design domain \\
		Initial design density & $\gamma_0$ & $V_f=0.30$, matched to the filtered volume \\
		Brinkman penalty & $\alpha(\bar{\gamma})$ & RAMP form with $\alpha_{\mathrm{fluid}}=0$, $\alpha_{\mathrm{solid}}=10^9$ \\
			RAMP continuation & $q$ & 8-stage schedule: 0.01, 0.02, 0.04, 0.08, 0.08, 0.12, 0.20, 0.35 \\
			Projection continuation & $\beta$ & 8-stage schedule: 1, 2, 4, 8, 8, 12, 16, 32, with threshold $\eta=0.50$ \\
			MMA move limits & -- & 8-stage schedule: 0.05, 0.04, 0.03, 0.02, 0.012, 0.010, 0.0075, 0.005 \\
			Maximum optimization iterations & -- & 8-stage schedule: 50, 60, 80, 100, 200, 180, 220, 260 MMA iterations; stop if objective change is below $10^{-6}$ for ten iterations \\
		Filter radius & $r_f$ & $4h_{\max}=2.22\times10^{-2}L$ \\
		SA working-variable penalty & -- & Power-law solid indicator with amplitude $10^5$ and exponent $n=3$ \\
		Wall-distance penalty & -- & Power-law solid indicator with amplitude $10^5$ and exponent $n=3$ \\
		Wall-distance equation & -- & Reciprocal penalized relaxed Eikonal equation with $\sigma_w=0.1$ and $G_0=20$ \\
			Forward flow solver & -- & IPCS flow solves for frozen-SA Picard updates, followed by an SNES Newton line-search final flow solve \\
		Frozen NS-SA coupling & -- & Three Picard updates per design iteration with SA relaxation factor $0.15$ \\
			Linear solver & -- & MUMPS for SNES/state and adjoint solves; BiCGSTAB with ILU preconditioning for IPCS Picard updates \\
		Adjoint treatment & -- & Frozen-turbulence adjoint in $(\mathbf{u}^*,p^*)$; $\tilde{\nu}$ and $G$ are frozen during gradient assembly \\
		Sensitivity verification samples & -- & Five active DG0 cells, seed $13$ \\
		Sensitivity verification step & $\Delta h$ & $10^{-6}$, central finite difference \\
		Sensitivity verification fields held fixed & -- & SA working variable $\tilde{\nu}^{0}$ and reciprocal wall-distance field $G^{0}$ \\
		\hline
	\end{tabularx}
\end{table}

\subsection{Velocity-Outlet Boundary-Condition Baseline}
\label{subsec:app_C_yoon_diffuser_velocity_outlet}
The geometry and boundary conditions follow Fig.~20 of Yoon~\cite{Yoon2016}. The design domain is a square of side length $L=1~\mathrm{m}$. A parabolic inflow enters through the full left boundary with maximum speed $U_{\max,\mathrm{in}}=3~\mathrm{m/s}$, while a parabolic outflow is prescribed on the center third of the right boundary with maximum speed $U_{\max,\mathrm{out}}=9~\mathrm{m/s}$. The remaining external boundaries are treated as no-slip walls. The Reynolds number is $Re=3{,}000$ based on $U_{\max,\mathrm{in}}L\rho/\mu$, and the present density convention converts Yoon's $70\%$ solid mass usage into a maximum fluid volume fraction of $V_f=0.30$.
Figure~\ref{fig:yoon_diffuser_velout_setup} shows the corresponding local geometry and boundary-condition layout used for the velocity-outlet baseline.

\begin{figure}[!ht]
	\centering
	\includegraphics[width=0.8\textwidth]{Diffuser-Geometry-BCs}
	\caption{Velocity-outlet diffuser baseline geometry and boundary conditions. Adapted from Yoon~\cite{Yoon2016}.}
	\label{fig:yoon_diffuser_velout_setup}
\end{figure}

Figure~\ref{fig:yoon_diffuser_velout_fenics_results} reports the FEniCS-SA topology and velocity field for the velocity-outlet baseline. The outlet region widens before the prescribed outlet segment. This feature should not be interpreted as a purely geometry-driven diffuser optimum: the optimizer must deliver the imposed parabolic outflow over the fixed outlet boundary, and opening the channel near that boundary reduces the shear and redistribution needed to satisfy the Dirichlet profile.

\begin{figure}[!ht]
	\centering
	\begin{subfigure}[b]{0.47\textwidth}
		\centering
		\includegraphics[width=\textwidth]{Diffuser-Turbulent-Topology-VelOut}
		\caption{Projected density field $\bar{\gamma}$.}
		\label{fig:yoon_diffuser_velout_topology}
	\end{subfigure}
	\hfill
	\begin{subfigure}[b]{0.50\textwidth}
		\centering
		\includegraphics[width=\textwidth]{Diffuser-Turbulent-Velocity-FEniCS-VelOut}
		\caption{Velocity magnitude and vector field.}
		\label{fig:yoon_diffuser_velout_velocity}
	\end{subfigure}
	\caption{FEniCS-SA result for the velocity-outlet diffuser baseline at $V_f=0.30$. Final objective value $J_D=3.1164\times10^5~\mathrm{W/m}$.}
	\label{fig:yoon_diffuser_velout_fenics_results}
\end{figure}

Figure~\ref{fig:yoon_diffuser_velout_ansys_separation} shows the corresponding Ansys-SA velocity contour for the extracted velocity-outlet design. The contour is used here only as a diagnostic of the baseline geometry: the widened outlet region contains low-speed separated or recovering flow, which makes the original diffuser result strongly sensitive to the outflow treatment and turbulence closure.

\begin{figure}[!ht]
	\centering
	\includegraphics[width=0.75\textwidth]{Diffuser-Turbulent-VelMagn-VelOut-SA}
	\caption{Ansys-SA velocity-magnitude contour for the extracted velocity-outlet diffuser baseline.}
	\label{fig:yoon_diffuser_velout_ansys_separation}
\end{figure}

\FloatBarrier

\subsection{Pressure-Outlet Turbulent FEniCS Optimization Diagnostics}
\label{subsec:app_C_yoon_diffuser_pressure_turb_diagnostics}
Figure~\ref{fig:yoon_diffuser_pressure_convergence} gives the pressure-outlet diffuser optimization history through the common $1150$-iteration comparison point used in Chapter~\ref{ch:benchmark_results}. The dissipation objective is normalized by its initial value, while the volume trace is plotted as the relative residual $100(v_f^{(k)}-V_f)/V_f$ so that constraint activity can be read on the same iteration axis.

\begin{figure}[htbp]
	\centering
	\includegraphics[width=\textwidth]{Diffuser-Turbulent-PreOut-Convergence-History}
	\caption{Optimization convergence history for the turbulent pressure-outlet diffuser benchmark, showing $J_D/J_{D,0}$ and the relative volume-fraction residual.}
	\label{fig:yoon_diffuser_pressure_convergence}
\end{figure}

\FloatBarrier

\subsection{Frozen Sensitivity Verification}
\label{subsec:yoon_diffuser_sensitivity}
	The frozen adjoint is checked directly on the pressure-outlet diffuser sensitivity configuration in \texttt{Config\_DiffuserYoon\_Frozen\_Sensitivity.py}. That file inherits the current pressure-outlet diffuser setup, including the dissipation objective $J_D$, the \texttt{design} wall-density source, the \texttt{reciprocal\_penalized} wall-distance convention, and the same power-law topology penalties used by the production diffuser run. The purpose of the check is therefore narrow and explicit: compare the frozen derivative supplied to MMA with objective changes obtained by re-solving the perturbed velocity-pressure state while holding the SA working variable and reciprocal wall-distance field fixed.

	The checked-in coordinate finite-difference result is performed at the initial diffuser design state, with base objective
	$J_0=1.1365575605\times10^{8}$. The sensitivity configuration evaluates five active $DG_0$ samples with seed $13$ and perturbation size $\Delta h=10^{-6}$. For sampled cell $i$, the raw density vector $\boldsymbol{\gamma}$ is perturbed by $\pm\Delta h\,\mathbf{e}_i$, where $\mathbf{e}_i$ is the unit vector for that active design variable, and the central finite-difference derivative is compared against the frozen-adjoint value.

\begin{table}[!ht]
	\centering
	\scriptsize
	\caption{Coordinate sensitivity check for the frozen pressure-outlet diffuser configuration. Values are taken from \texttt{SensitivityVerificationTable.tsv}.}
	\label{tab:yoon_diffuser_fd_adjoint_comparison}
	\renewcommand{\arraystretch}{1.2}
	\begin{tabular}{l c c c c}
		\hline
		\textbf{Cell} &
		\shortstack{\textbf{Finite}\\\textbf{difference}\\$S_i^{\mathrm{FD}}$} &
		\shortstack{\textbf{Frozen}\\\textbf{adjoint}\\$S_i^{\mathrm{fr}}$} &
		\shortstack{\textbf{Absolute}\\\textbf{error}\\$\Delta_i^{\mathrm{abs}}$} &
		\shortstack{\textbf{Relative}\\\textbf{error}\\$\varepsilon_i^{\mathrm{fr}}$} \\
		\hline
		$CV_1$ & $-2.1542\times10^{4}$ & $-2.1541\times10^{4}$ & $6.7475\times10^{-1}$ & $3.1322\times10^{-5}$ \\
		$CV_2$ & $-3.9061\times10^{3}$ & $-3.9059\times10^{3}$ & $2.4168\times10^{-1}$ & $6.1872\times10^{-5}$ \\
		$CV_3$ & $-1.4988\times10^{4}$ & $-1.4990\times10^{4}$ & $1.3499\times10^{0}$ & $9.0058\times10^{-5}$ \\
		$CV_4$ & $-7.6433\times10^{3}$ & $-7.6430\times10^{3}$ & $2.3322\times10^{-1}$ & $3.0513\times10^{-5}$ \\
		$CV_5$ & $-4.3901\times10^{3}$ & $-4.3896\times10^{3}$ & $5.7420\times10^{-1}$ & $1.3079\times10^{-4}$ \\
		\hline
	\end{tabular}
\end{table}

	The coordinate checks show consistent sign agreement and small relative errors across all five samples. The maximum relative discrepancy is $1.31\times10^{-4}$, while the largest absolute discrepancy is $1.35$. That absolute value is not used as the primary acceptance metric here: the sampled diffuser sensitivities themselves are of order $10^{4}$, and the base dissipation objective is of order $10^{8}$, so order-unity absolute derivative mismatches still correspond to small relative discrepancies in the actual derivative scale being checked.

	The same sampled active cells are also used for a directional Taylor check with steps $(10^{-3},3\times10^{-4},10^{-4},3\times10^{-5})$ and seed $29$. The directional derivative is
	$-1.3711285431\times10^{4}$, and the central finite differences remain within $1.09\times10^{-6}$ relative error over the tested step range. The symmetric Taylor remainder decays with approximately second-order behavior over the first three step reductions, as reported in Table~\ref{tab:yoon_diffuser_taylor_check}.

\begin{table}[!ht]
	\centering
	\scriptsize
	\caption{Directional Taylor check for the frozen pressure-outlet diffuser sensitivity configuration. Values are taken from \texttt{TaylorCheckLog.tsv}.}
	\label{tab:yoon_diffuser_taylor_check}
	\renewcommand{\arraystretch}{1.2}
	\begin{tabular}{c c c c c}
		\hline
		\textbf{Step $\epsilon$} &
		\shortstack{\textbf{Frozen}\\\textbf{directional}\\\textbf{derivative}} &
		\shortstack{\textbf{Central}\\\textbf{finite}\\\textbf{difference}} &
		\shortstack{\textbf{Directional}\\\textbf{relative}\\\textbf{error}} &
		\shortstack{\textbf{Symmetric}\\\textbf{Taylor}\\\textbf{order}} \\
		\hline
		$1.0\times10^{-3}$ & $-1.3711285431\times10^{4}$ & $-1.3711284466\times10^{4}$ & $7.0351\times10^{-8}$ & -- \\
		$3.0\times10^{-4}$ & $-1.3711285431\times10^{4}$ & $-1.3711284200\times10^{4}$ & $8.9732\times10^{-8}$ & $2.0079$ \\
		$1.0\times10^{-4}$ & $-1.3711285431\times10^{4}$ & $-1.3711295649\times10^{4}$ & $7.4528\times10^{-7}$ & $2.0128$ \\
		$3.0\times10^{-5}$ & $-1.3711285431\times10^{4}$ & $-1.3711300244\times10^{4}$ & $1.0804\times10^{-6}$ & $2.7943$ \\
		\hline
	\end{tabular}
\end{table}

	This check therefore verifies the implemented frozen-turbulence derivative for the pressure-outlet diffuser benchmark under the dissipation objective $J_D$. As with the bend cases, the conclusion is deliberately limited to that derivative: the SA working variable and reciprocal wall-distance field are held fixed during the checks, matching the optimization adjoint. It should not be interpreted as validation of a fully coupled turbulence sensitivity.

\FloatBarrier

The extracted turbulent diffuser geometry is meshed with the same wall-resolved triangular strategy as the laminar diffuser comparison in Table~\ref{tab:app_C_yoon_diffuser_lam_mesh}, but with a distinct coarse--medium--fine hierarchy. The first-layer height is kept fixed at $\Delta y_1=1.50\times10^{-3}~\mathrm{m}$ to match the laminar-design diffuser meshes, while the bulk size, curvature angle, and inflation-layer distribution are refined independently for the turbulent topology.

\begin{table}[htbp]
	\centering
	\footnotesize
	\caption{Ansys Fluent mesh-independence table for the extracted turbulent diffuser benchmark geometry.}
	\label{tab:app_B_yoon_diffuser_mesh}
	\renewcommand{\arraystretch}{1.2}
	\setlength{\tabcolsep}{5pt}
	\begin{tabular*}{\textwidth}{@{\extracolsep{\fill}}lccc@{}}
		\hline
		\textbf{Parameter} & \textbf{Coarse} & \textbf{Medium} & \textbf{Fine} \\
		\hline
		\multicolumn{4}{@{}l}{\textit{Mesh settings}} \\
		Mesh type & Tri. & Tri. & Tri. \\
		Near-wall & WR infl. & WR infl. & WR ST infl. \\
		$h_{\mathrm{bulk}}$ [m] & $7.50 \times 10^{-3}$ & $5.00 \times 10^{-3}$ & $4.00 \times 10^{-3}$ \\
		Curvature angle & $9^\circ$ & $6^\circ$ & $5^\circ$ \\
		$\Delta y_1$ [m] & $1.50 \times 10^{-3}$ & $1.50 \times 10^{-3}$ & $1.50 \times 10^{-3}$ \\
		Inflation layers & $10$ & $14$ & $16$ \\
		Growth rate & $1.11$ & $1.09$ & $1.08$ \\
		$t_{\mathrm{infl}}$ [m] & $2.51 \times 10^{-2}$ & $3.90 \times 10^{-2}$ & $4.55 \times 10^{-2}$ \\
		Nodes & $8{,}104$ & $16{,}166$ & $21{,}238$ \\
		Elements & $12{,}330$ & $24{,}570$ & $40{,}268$ \\
		\hline
		\multicolumn{4}{@{}l}{\textit{Validation metrics}} \\
		$\Phi_D^{\mathrm{SA}}$ [W/m] & $1.2921\times10^{5}$ & $1.2966\times10^{5}$ & $1.2995\times10^{5}$ \\
		$\epsilon_{\mathrm{rel}}(\Phi_D^{\mathrm{SA}})$ [\%] & -- & $0.35$ & $0.22$ \\
		$\Phi_D^{\mathrm{SST}\,k\text{-}\omega}$ [W/m] & $1.5296\times10^{4}$ & $1.5570\times10^{4}$ & $1.5725\times10^{4}$ \\
		$\epsilon_{\mathrm{rel}}(\Phi_D^{\mathrm{SST}\,k\text{-}\omega})$ [\%] & -- & $1.79$ & $0.99$ \\
		Static pressure drop $\Delta p_s^{\mathrm{SA}}$ [Pa] & $1.2460\times10^{5}$ & $1.2449\times10^{5}$ & $1.2422\times10^{5}$ \\
		$\epsilon_{\mathrm{rel}}(\Delta p_s^{\mathrm{SA}})$ [\%] & -- & $0.09$ & $0.22$ \\
		Static pressure drop $\Delta p_s^{\mathrm{SST}\,k\text{-}\omega}$ [Pa] & $4.2412\times10^{4}$ & $4.2034\times10^{4}$ & $4.1833\times10^{4}$ \\
		$\epsilon_{\mathrm{rel}}(\Delta p_s^{\mathrm{SST}\,k\text{-}\omega})$ [\%] & -- & $0.89$ & $0.48$ \\
		Selected validation mesh & \multicolumn{3}{c}{Fine\_Ansys} \\
		\hline
	\end{tabular*}
	\vspace{2pt}
	\begin{minipage}{0.98\textwidth}
		\scriptsize \textit{Note.} WR = wall-resolved; ST = smooth-transition; $\Phi_D$ = viscous dissipation; $\Delta p_s$ = static pressure drop. For $\epsilon_{\mathrm{rel}}$ rows, the Medium column gives the Coarse--Medium change and the Fine column gives the Medium--Fine change.
	\end{minipage}
\end{table}

\FloatBarrier

\subsection{Turbulent-Design Wall-Shear Diagnostic}
\label{subsec:app_C_yoon_diffuser_wall_shear}
Figure~\ref{fig:app_C_yoon_diffuser_wall_shear_sign_map} summarizes the downstream wall-shear sign changes for the extracted pressure-outlet turbulent diffuser on the fine Ansys mesh. The diagnostic is included only to support the separated-flow interpretation in Chapter~\ref{ch:benchmark_results}; the performance ranking itself is taken from the integrated pressure-drop and dissipation metrics.

\begin{figure}[htbp]
	\centering
	\makebox[\textwidth][c]{\includegraphics[width=1.08\textwidth]{Diffuser-Turbulent-WallShear-SignMap}}
	\caption{Downstream wall-shear sign map for the extracted pressure-outlet turbulent diffuser. Reversed wall shear denotes $\tau_{w,x}<0$. Interpolated locations: upper wall $x/L=0.8566$ (SST $k-\omega$), $0.8723$ (SA); lower wall $x/L=0.8406$, $0.9005$, $0.9795$ (SST $k-\omega$) and $0.8386$, $0.9109$, $0.9480$ (SA).  \ChatGPTFigureRef.}
	\label{fig:app_C_yoon_diffuser_wall_shear_sign_map}
\end{figure}

\FloatBarrier

\subsection{Additional Diffuser Separation Diagnostic}
\label{subsec:app_C_yoon_diffuser_separation_sst}
Figure~\ref{fig:app_C_yoon_diffuser_separation_sst} gives a close-up of the SST $k-\omega$ velocity vectors near the diffuser outlet for the extracted laminar and turbulent pressure-outlet designs. This figure is used as supporting visual evidence for the separated or recovering near-wall region discussed in Chapter~\ref{ch:benchmark_results}; the wall-shear sign map in Figure~\ref{fig:app_C_yoon_diffuser_wall_shear_sign_map} remains the primary separation diagnostic because it directly indicates local flow reversal at the wall.

\begin{figure}[htbp]
	\centering
	\begin{subfigure}[b]{\textwidth}
		\centering
		\includegraphics[width=\textwidth]{Diffuser-Laminar-Separation-SSTkomega}
		\caption{Laminar pressure-outlet design.}
		\label{fig:app_C_yoon_diffuser_laminar_separation_sst}
	\end{subfigure}

	\vspace{0.5em}

	\begin{subfigure}[b]{\textwidth}
		\centering
		\includegraphics[width=\textwidth]{Diffuser-Turbulent-Separation-SSTkomega}
		\caption{Turbulent pressure-outlet design.}
		\label{fig:app_C_yoon_diffuser_turbulent_separation_sst}
	\end{subfigure}
	\caption{SST $k-\omega$ velocity-vector close-up near the outlet of the extracted pressure-outlet diffuser designs.}
	\label{fig:app_C_yoon_diffuser_separation_sst}
\end{figure}

\FloatBarrier

\subsection{Laminar FEniCS Configuration Differences}
\label{subsec:app_C_yoon_diffuser_lam_config}
	The existing laminar diffuser reference is a completed FEniCS topology-optimization run from \texttt{Config\_DiffuserYoon\_Laminar.py}. It reuses the pressure-outlet turbulent diffuser setup in Tables~\ref{tab:app_B_yoon_diffuser_physical} and~\ref{tab:app_B_yoon_diffuser_numerical} wherever not stated otherwise: the square $L=1~\mathrm{m}$ domain, full-height parabolic inlet, center-third pressure outlet with $p=0~\mathrm{Pa}$, no prescribed outlet velocity, $V_f=0.30$ volume constraint, dissipation objective $J_D$, Brinkman limits, mesh, filter radius, and projection threshold are unchanged. The laminar Reynolds number, solver settings, and continuation controls are listed in Table~\ref{tab:app_C_yoon_diffuser_lam_differences}.

\begin{table}[htbp]
	\centering
	\footnotesize
	\caption{Laminar-specific differences from the pressure-outlet turbulent diffuser FEniCS setup above.}
	\label{tab:app_C_yoon_diffuser_lam_differences}
	\renewcommand{\arraystretch}{1.2}
	\begin{tabularx}{\textwidth}{l X}
		\hline
		\textbf{Parameter} & \textbf{Laminar setting} \\
		\hline
			Flow model & Laminar incompressible Navier--Stokes; the SA working variable, SA wall-distance equation, and frozen-turbulence coupling are omitted \\
			Reynolds number & $Re=1$, based on $U_{\max,\mathrm{in}}L\rho/\mu$ instead of the turbulent $Re=3{,}000$ setup \\
			Density and viscosity & $\rho=1000~\mathrm{kg/m^3}$ and $\mu=3000~\mathrm{Pa\,s}$, with the same $U_{\max,\mathrm{in}}=3~\mathrm{m/s}$ as the turbulent pressure-outlet case; no outlet velocity is prescribed \\
			Flow finite elements & Taylor--Hood $P_2/P_1$ for $(\mathbf{u},p)$; no separate $P_1$ SA working-variable space is used \\
			Forward and adjoint solves & SNES/MUMPS laminar Newton forward solves with method \texttt{newtontr}, relative tolerance $10^{-6}$, absolute tolerance $10^{-9}$, and maximum forward iteration count $80$; the adjoint is solved as a linear system with MUMPS \\
			Continuation and optimizer controls & Seven-stage schedule $q=\{0.01,0.02,0.04,0.08,0.08,0.12,0.20\}$, $\beta=\{1,2,4,8,8,12,16\}$, move limits $\{0.010,0.008,0.006,0.004,0.003,0.002,0.0015\}$, and inner-iteration budgets $\{12,25,45,80,120,160,220\}$; stop if objective change is below $10^{-6}$ for ten iterations \\
		\hline
	\end{tabularx}
\end{table}

\subsection{Laminar FEniCS Optimization Diagnostics}
\label{subsec:app_C_yoon_diffuser_lam_diagnostics}
Figure~\ref{fig:app_C_yoon_diffuser_lam_convergence} records the completed pressure-outlet laminar diffuser optimization history. The normalized objective drops rapidly during the early continuation stages and then remains nearly flat, while the volume-fraction residual stays close to the active constraint. This figure is kept in the appendix because it documents the laminar reference run rather than the turbulent performance comparison itself.

\begin{figure}[htbp]
	\centering
	\includegraphics[width=\textwidth]{Diffuser-Laminar-PreOut-Convergence-History}
	\caption{FEniCS convergence history for the pressure-outlet laminar diffuser topology-optimization run, showing $J_D/J_{D,0}$ and the relative volume-fraction residual.}
	\label{fig:app_C_yoon_diffuser_lam_convergence}
\end{figure}

\FloatBarrier

\subsection{Laminar-Design Mesh Independence}
\label{subsec:app_C_yoon_diffuser_lam_mesh}
Table~\ref{tab:app_C_yoon_diffuser_lam_mesh} reports the Ansys mesh-independence study for the extracted pressure-outlet laminar diffuser design. The turbulent diffuser mesh study is reported separately in Table~\ref{tab:app_B_yoon_diffuser_mesh} of Section~\ref{sec:app_B_config_yoon_diffuser}. The Ansys Fluent solves use the same shifted inlet-profile and pressure-outlet convention as the turbulent diffuser validation setup.

\begin{table}[htbp]
	\centering
	\footnotesize
	\caption{Ansys mesh-independence table for the extracted pressure-outlet laminar diffuser design.}
	\label{tab:app_C_yoon_diffuser_lam_mesh}
	\renewcommand{\arraystretch}{1.2}
	\setlength{\tabcolsep}{5pt}
	\begin{tabular*}{\textwidth}{@{\extracolsep{\fill}}lccc@{}}
		\hline
		\textbf{Parameter} & \textbf{Coarse} & \textbf{Medium} & \textbf{Fine} \\
		\hline
		\multicolumn{4}{@{}l}{\textit{Mesh settings}} \\
		Mesh type & Tri. & Tri. & Tri. \\
		Near-wall & WR infl. & WR infl. & WR infl. \\
		$h_{\mathrm{bulk}}$ [m] & $1.20 \times 10^{-2}$ & $5.50 \times 10^{-3}$ & $3.75 \times 10^{-3}$ \\
		Curvature angle & $15^\circ$ & $7^\circ$ & $5^\circ$ \\
		$\Delta y_1$ [m] & $1.50 \times 10^{-3}$ & $1.50 \times 10^{-3}$ & $1.50 \times 10^{-3}$ \\
		Inflation layers & $6$ & $12$ & $16$ \\
		Growth rate & $1.15$ & $1.10$ & $1.08$ \\
		$t_{\mathrm{infl}}$ [m] & $1.31 \times 10^{-2}$ & $3.19 \times 10^{-2}$ & $4.55 \times 10^{-2}$ \\
		Nodes & $3{,}317$ & $13{,}444$ & $25{,}485$ \\
		Elements & $5{,}085$ & $20{,}888$ & $39{,}582$ \\
		\hline
		\multicolumn{4}{@{}l}{\textit{Validation metrics}} \\
		$\Phi_D^{\mathrm{SST}\,k\text{-}\omega}$ [W/m] & $1.2643 \times 10^{4}$ & $1.2983 \times 10^{4}$ & $1.3077 \times 10^{4}$ \\
		$\epsilon_{\mathrm{rel}}(\Phi_D^{\mathrm{SST}\,k\text{-}\omega})$ [\%] & -- & $2.69$ & $0.72$ \\
		Static pressure drop $\Delta p_s^{\mathrm{SST}\,k\text{-}\omega}$ [Pa] & $3.6628 \times 10^{4}$ & $3.6693 \times 10^{4}$ & $3.6464 \times 10^{4}$ \\
		$\epsilon_{\mathrm{rel}}(\Delta p_s^{\mathrm{SST}\,k\text{-}\omega})$ [\%] & -- & $0.18$ & $0.62$ \\
		$\Phi_D^{\mathrm{SA}}$ [W/m] & $1.4296 \times 10^{5}$ & $1.4315 \times 10^{5}$ & $1.4333 \times 10^{5}$ \\
		$\epsilon_{\mathrm{rel}}(\Phi_D^{\mathrm{SA}})$ [\%] & -- & $0.13$ & $0.13$ \\
		Static pressure drop $\Delta p_s^{\mathrm{SA}}$ [Pa] & $1.4191 \times 10^{5}$ & $1.4276 \times 10^{5}$ & $1.4200 \times 10^{5}$ \\
		$\epsilon_{\mathrm{rel}}(\Delta p_s^{\mathrm{SA}})$ [\%] & -- & $0.60$ & $0.53$ \\
		Selected comparison mesh & \multicolumn{3}{c}{Fine\_Ansys} \\
		\hline
	\end{tabular*}
	\vspace{2pt}
	\begin{minipage}{0.98\textwidth}
		\scriptsize \textit{Note.} WR = wall-resolved; $\Phi_D$ = viscous dissipation; $\Delta p_s$ = static pressure drop. For $\epsilon_{\mathrm{rel}}$ rows, the Medium column gives the Coarse--Medium change and the Fine column gives the Medium--Fine change. The SA rows are re-evaluations for turbulence-model sensitivity.
	\end{minipage}
\end{table}

\FloatBarrier

\backmatter
% The bibliography comes after the appendices.
% You can replace the standard "abbrv" bibliography style by another one.
\bibliographystyle{abbrv}
\bibliography{references}
\include{CoC}
\end{document}
